% Lutte contre les problèmes liés à la régulation de la glycémie et à la pression artérielle — SVT Tle D — Cours signé OnBuch+
% Programme officiel MINESEC. Compiler avec : tectonic N13-regulation-glycemie-pression-arterielle.tex
\newcommand{\DOCMATIERE}{SVT}
\newcommand{\DOCNIVEAU}{Tle D}
\newcommand{\DOCMODULE}{Module II : L'Éducation à la Santé — La santé nutritionnelle}
\newcommand{\DOCLECON}{N13}
\newcommand{\DOCTITRE}{Lutte contre les problèmes liés à la régulation de la glycémie et à la pression artérielle}
% OnBuch+ — préambule « cours signés » ENRICHI (compilé avec tectonic / XeLaTeX)
% Système visuel « Pop » d'OnBuch+ : fond crème (jamais blanc), bordures encre
% épaisses, ombres « dures » décalées, titres Archivo Black en capitales.
% Version MATHÉMATIQUES : graphes (pgfplots/tikz), boîtes pédagogiques
% (définition, propriété, méthode, attention, le savais-tu, exercice résolu,
% exercice, TP, bilan), page de garde et signature OnBuch+ ancrée en bas.
%
% Macros attendues dans main.tex AVANT \input{preamble} :
%   \DOCMATIERE  \DOCNIVEAU  \DOCMODULE  \DOCLECON (numéro)  \DOCTITRE
\documentclass[11pt]{article}

\usepackage{fontspec}
\usepackage[french]{babel}
\usepackage[a4paper,margin=2cm,top=3.4cm,bottom=2.8cm,headheight=1.4cm,footskip=1.3cm]{geometry}
\usepackage{amsmath,amssymb,amsfonts}
\usepackage{mathtools} % \xmapsto, \coloneqq... souvent utilisés par les modèles
\usepackage{cancel} % simplifications barrées (bilans de forces)
% \abs{z} (module, valeur absolue) et \norm{u} (norme), utilisés par les modèles
\providecommand{\abs}{}\let\abs\relax\DeclarePairedDelimiter{\abs}{\lvert}{\rvert}
\providecommand{\norm}{}\let\norm\relax\DeclarePairedDelimiter{\norm}{\lVert}{\rVert}
\usepackage[table]{xcolor} % [table] : fournit aussi \rowcolors (alternance de lignes)
\usepackage{pagecolor}
\usepackage{tikz}
\usetikzlibrary{arrows.meta,positioning,calc,shapes.geometric,shapes.misc,shapes.arrows,decorations.pathmorphing,decorations.pathreplacing,decorations.markings,patterns,shadows,mindmap,trees,fit,backgrounds,matrix,intersections,angles,quotes}
\usepackage{pgfplots}
\usepackage[version=4]{mhchem} % équations nucléaires \ce{^{235}_{92}U}
\usepackage[european,siunitx]{circuitikz} % schémas électriques
\pgfplotsset{compat=1.18}
\usepgfplotslibrary{fillbetween,groupplots} % aires entre courbes (calcul intégral)
\usepackage{enumitem}
\usepackage{fancyhdr}
% [most] charge toutes les bibliothèques tcolorbox (listings, minted, raster,
% documentation...) et gonfle inutilement la mémoire TeX sur un document long
% (~300 boîtes) : on ne charge que ce qui est réellement utilisé.
\usepackage[skins,breakable]{tcolorbox}
\usepackage{graphicx}
\usepackage{colortbl}
\usepackage{booktabs}
\usepackage{tabularx}
\usepackage{diagbox} % case d'en-tête diagonale des tableaux à double entrée
% Colonnes extensibles centrée / gauche / droite (C, L, R), souvent utilisées par les modèles
\newcolumntype{C}{>{\centering\arraybackslash}X}
\newcolumntype{L}{>{\raggedright\arraybackslash}X}
\newcolumntype{R}{>{\raggedleft\arraybackslash}X}
\usepackage{array}
\usepackage{multicol}
\usepackage{siunitx}
% parse-numbers=false : accepte aussi \SI{2,5 \times 10^{-3}}{...}
\sisetup{locale=FR,output-decimal-marker={,},group-minimum-digits=4,parse-numbers=false}
\DeclareSIUnit\ohm{\ensuremath{\symup{\Omega}}} % U+2126 absent de la police
\DeclareSIUnit\division{div} % réglages d'oscilloscope (ms/div, V/div)
\DeclareSIUnit\tr{tr} % tours (vitesse de rotation, ex. \SI{1500}{\tr\per\minute})
\DeclareSIUnit\molar{M} % concentration molaire (ex. \SI{3}{\molar}, \SI{1}{\micro\molar})
\DeclareSIUnit\an{an} % année (le modèle confond parfois avec \year, un compteur TeX) — ex. \SI{5}{\kilo\meter\per\an}
\DeclareSIUnit\FCFA{FCFA} % franc CFA (ex. \SI{500}{\FCFA\per\kilo\gram})
\DeclareSIUnit\degreeBrix{\textdegree Brix} % teneur en sucre d'un sirop/jus (réfractométrie)
\DeclareSIUnit\month{mois} % le modèle utilise parfois \month, un compteur TeX, comme unité de durée
\usepackage{qrcode}
% \degree : commande fréquemment utilisée par les modèles pour un angle ou une
% température en dehors de \SI{}{} (non fournie par les paquets chargés).
\providecommand{\degree}{^{\circ}}
% \qed (fin de démonstration), utilisé par les modèles sans amsthm
\providecommand{\qed}{\hfill$\square$}
% \barre : double barre des valeurs interdites dans un tableau de variations (tkz-tab)
\providecommand{\barre}{\ensuremath{\|}}
% environnement proof (amsthm non chargé) utilisé par les modèles
\ifdefined\proof\else\newenvironment{proof}[1][Démonstration]{\par\noindent\textit{#1.}\ }{\qed\par}\fi
\usepackage{lastpage}
\usepackage{hyperref}
\hypersetup{hidelinks}

% ============================================================
% Palette « Pop » OnBuch+ — mêmes hex que la classe P (plus_ui.dart)
% ============================================================
\definecolor{poporange}{HTML}{F59321}
\definecolor{popdark}{HTML}{E07A0C}
\definecolor{popcream}{HTML}{FFF6E8}
\definecolor{popink}{HTML}{1C1714}
\definecolor{poppurple}{HTML}{7A5AE0}
\definecolor{popgreen}{HTML}{2E9E6A}
\definecolor{poppink}{HTML}{E0457B}
\definecolor{popblue}{HTML}{2F7DE1}
\definecolor{popgold}{HTML}{C98A12}
\definecolor{popmuted}{HTML}{8A7F76}
\definecolor{popline}{HTML}{EADFD2}
\definecolor{popgray}{HTML}{8A7F76}
\colorlet{popgrey}{popgray}
% Alias tolérants (le modèle nomme parfois les couleurs différemment)
\colorlet{popwhite}{white}
\colorlet{popblack}{popink}
\colorlet{popred}{poppink}
\colorlet{popyellow}{popgold}
% Teintes claires (fonds de tableaux/encadrés) — le modèle nomme parfois
% les couleurs avec un suffixe L (ex. popblueL) sans les avoir définies.
\colorlet{poporangeL}{poporange!20}
\colorlet{popdarkL}{popdark!20}
\colorlet{poppurpleL}{poppurple!20}
\colorlet{popgreenL}{popgreen!20}
\colorlet{poppinkL}{poppink!20}
\colorlet{popblueL}{popblue!20}
\colorlet{popgoldL}{popgold!20}
\colorlet{popredL}{popred!20}
\colorlet{popgrayL}{popgray!20}
% Teintes claires pour fonds de tableaux / zones
\colorlet{popblueL}{popblue!10!white}
\colorlet{poporangeL}{poporange!14!white}
\colorlet{popgreenL}{popgreen!12!white}
\colorlet{poppurpleL}{poppurple!12!white}
\colorlet{poppinkL}{poppink!12!white}
\colorlet{popgoldL}{popgold!14!white}

\pagecolor{popcream}

% ============================================================
% Typographie
% ============================================================
\defaultfontfeatures{Ligatures=TeX}
\setmainfont{PlusJakartaSans-Medium.ttf}[Path=../../../fonts/,BoldFont=PlusJakartaSans-Bold.ttf]
% Maths en sans-serif (Fira Math) avec chiffres et lettres droites de Plus
% Jakarta, pour que les formules se fondent dans le texte.
\usepackage[math-style=ISO,bold-style=ISO]{unicode-math}
\setmathfont{FiraMath-Regular.otf}[Path=../../../fonts/]
\setmathfont{PlusJakartaSans-Medium.ttf}[Path=../../../fonts/,range={up/num,up/{Latin,latin},\mathup}]
\usepackage{icomma} % virgule décimale sans espace en mode math
% Caractères mathématiques Unicode absents des polices de texte (Plus Jakarta,
% Archivo Black) : rendus par leur équivalent mathématique, en texte comme en
% formule — sinon ils s'affichent en carré vide (ex. « Arithmétique dans ℤ »).
\usepackage{newunicodechar}
\newunicodechar{ℝ}{\ensuremath{\mathbb{R}}}
\newunicodechar{ℤ}{\ensuremath{\mathbb{Z}}}
\newunicodechar{ℂ}{\ensuremath{\mathbb{C}}}
\newunicodechar{ℕ}{\ensuremath{\mathbb{N}}}
\newunicodechar{ℚ}{\ensuremath{\mathbb{Q}}}
\newunicodechar{𝔻}{\ensuremath{\mathbb{D}}}
\newunicodechar{α}{\ensuremath{\alpha}}
\newunicodechar{β}{\ensuremath{\beta}}
\newunicodechar{γ}{\ensuremath{\gamma}}
\newunicodechar{δ}{\ensuremath{\delta}}
\newunicodechar{ε}{\ensuremath{\varepsilon}}
\newunicodechar{θ}{\ensuremath{\theta}}
\newunicodechar{λ}{\ensuremath{\lambda}}
\newunicodechar{μ}{\ensuremath{\mu}}
\newunicodechar{σ}{\ensuremath{\sigma}}
\newunicodechar{τ}{\ensuremath{\tau}}
\newunicodechar{φ}{\ensuremath{\varphi}}
\newunicodechar{ψ}{\ensuremath{\psi}}
\newunicodechar{ω}{\ensuremath{\omega}}
\newunicodechar{Σ}{\ensuremath{\Sigma}}
% majuscules produites par \MakeUppercase dans les titres (α-aminés -> Α)
\newunicodechar{Α}{\ensuremath{\alpha}}
\newunicodechar{Β}{\ensuremath{\beta}}
\newunicodechar{Γ}{\ensuremath{\Gamma}}
\newunicodechar{Δ}{\ensuremath{\Delta}}
\newunicodechar{Ε}{\ensuremath{\varepsilon}}
\newunicodechar{Θ}{\ensuremath{\Theta}}
\newunicodechar{Λ}{\ensuremath{\Lambda}}
\newunicodechar{Μ}{\ensuremath{\mu}}
\newunicodechar{Τ}{\ensuremath{\tau}}
\newunicodechar{Φ}{\ensuremath{\Phi}}
\newunicodechar{Ψ}{\ensuremath{\Psi}}
\newunicodechar{Ω}{\ensuremath{\Omega}}
\newunicodechar{↦}{\ensuremath{\mapsto}}
\newunicodechar{∈}{\ensuremath{\in}}
\newunicodechar{∉}{\ensuremath{\notin}}
\newunicodechar{∧}{\ensuremath{\wedge}}
\newunicodechar{⇔}{\ensuremath{\Leftrightarrow}}
\newunicodechar{⟺}{\ensuremath{\Longleftrightarrow}}
\newunicodechar{⇒}{\ensuremath{\Rightarrow}}
\newunicodechar{⟹}{\ensuremath{\Longrightarrow}}
\newunicodechar{∘}{\ensuremath{\circ}}
\newunicodechar{∪}{\ensuremath{\cup}}
\newunicodechar{∩}{\ensuremath{\cap}}
\newunicodechar{≡}{\ensuremath{\equiv}}
\newunicodechar{⊂}{\ensuremath{\subset}}
\newunicodechar{⊥}{\ensuremath{\perp}}
\newunicodechar{∃}{\ensuremath{\exists}}
\newunicodechar{∀}{\ensuremath{\forall}}
\newunicodechar{∛}{\ensuremath{\sqrt[3]{\,}}}
\newunicodechar{∜}{\ensuremath{\sqrt[4]{\,}}}
\newunicodechar{⃗}{\ensuremath{{}^{\to}}}
\newunicodechar{ⁿ}{\ifmmode^{n}\else\textsuperscript{\ensuremath{n}}\fi}
\newunicodechar{ˣ}{\ifmmode^{x}\else\textsuperscript{\ensuremath{x}}\fi}
\newunicodechar{ᵇ}{\ifmmode^{b}\else\textsuperscript{\ensuremath{b}}\fi}
\newunicodechar{ʳ}{\ifmmode^{r}\else\textsuperscript{\ensuremath{r}}\fi}
\newunicodechar{ᵘ}{\ifmmode^{u}\else\textsuperscript{\ensuremath{u}}\fi}
\newunicodechar{ʸ}{\ifmmode^{y}\else\textsuperscript{\ensuremath{y}}\fi}
\newunicodechar{ᵃ}{\ifmmode^{a}\else\textsuperscript{\ensuremath{a}}\fi}
\newunicodechar{ᵉ}{\ifmmode^{e}\else\textsuperscript{\ensuremath{e}}\fi}
\newunicodechar{ᵏ}{\ifmmode^{k}\else\textsuperscript{\ensuremath{k}}\fi}
\newunicodechar{ᵖ}{\ifmmode^{p}\else\textsuperscript{\ensuremath{p}}\fi}
\newunicodechar{ᴺ}{\ifmmode^{N}\else\textsuperscript{\ensuremath{N}}\fi}
\newunicodechar{ᶜ}{\ifmmode^{c}\else\textsuperscript{\ensuremath{c}}\fi}
\newunicodechar{ʰ}{\ifmmode^{h}\else\textsuperscript{\ensuremath{h}}\fi}
\newunicodechar{ᵠ}{\ifmmode^{q}\else\textsuperscript{\ensuremath{q}}\fi}
\newunicodechar{ₙ}{\ifmmode_{n}\else\textsubscript{\ensuremath{n}}\fi}
\newunicodechar{ₐ}{\ifmmode_{a}\else\textsubscript{\ensuremath{a}}\fi}
\newunicodechar{ᵢ}{\ifmmode_{i}\else\textsubscript{\ensuremath{i}}\fi}
\newunicodechar{ᵣ}{\ifmmode_{r}\else\textsubscript{\ensuremath{r}}\fi}
\newunicodechar{ₖ}{\ifmmode_{k}\else\textsubscript{\ensuremath{k}}\fi}
\newunicodechar{ᵦ}{\ifmmode_{\beta}\else\textsubscript{\ensuremath{\beta}}\fi}
\newunicodechar{ₓ}{\ifmmode_{x}\else\textsubscript{\ensuremath{x}}\fi}
\newfontfamily\popdisplay{ArchivoBlack-Regular.ttf}[Path=../../../fonts/]
\newfontfamily\poplogo{SpaceGrotesk-Bold.ttf}[Path=../../../fonts/]
\newfontfamily\popsemi{PlusJakartaSans-SemiBold.ttf}[Path=../../../fonts/]
\newfontfamily\popxbold{PlusJakartaSans-ExtraBold.ttf}[Path=../../../fonts/]
\newfontfamily\popxboldls{PlusJakartaSans-ExtraBold.ttf}[Path=../../../fonts/,LetterSpace=3.0]

\setlist[enumerate]{leftmargin=1.7em,itemsep=5pt,topsep=4pt}
\setlist[itemize]{leftmargin=1.7em,itemsep=4pt,topsep=4pt,label={\color{poporange}\textbullet}}
\setlength{\parindent}{0pt}
\setlength{\parskip}{5pt}
\renewcommand{\arraystretch}{1.25}

% pgfplots : style Pop par défaut
\pgfplotsset{
  every axis/.append style={axis line style={popink,line width=1pt},tick style={popink},
    grid style={popline,line width=0.5pt},label style={font=\small},tick label style={font=\footnotesize}},
  popcurve/.style={poporange,line width=1.8pt,no markers,smooth},
}
% Même style disponible dans un \draw TikZ simple (hors pgfplots)
\tikzset{popcurve/.style={poporange,line width=1.8pt}}
\tikzset{popgrid/.style={popline,very thin}}
\colorlet{popcurve}{poporange} % le nom du style est parfois utilisé comme couleur

% ============================================================
% Pill / squiggle
% ============================================================
\newcommand{\poppill}[3]{%
  \tikz[baseline=-0.55ex]\node[draw=popink,line width=1.2pt,rounded corners=2.6pt,%
    fill=#2,inner xsep=6.5pt,inner ysep=3pt]{%
    \popxboldls\fontsize{7}{7}\selectfont\color{#3}\MakeUppercase{#1}};%
}
\newcommand{\popsquiggle}[1]{%
  \tikz[baseline=-0.4ex]{
    \draw[#1,line width=2.6pt,line cap=round]
      plot [smooth] coordinates {(0,0.09) (0.34,0.01) (0.68,0.21) (1.02,0.01) (1.36,0.10)};
  }%
}
% Mot-clé surligné (marqueur orange sous le texte)
\newcommand{\cle}[1]{{\popxbold\color{popdark}#1}}

% ============================================================
% En-tête / pied
% ============================================================
\pagestyle{fancy}
\fancyhf{}
\renewcommand{\headrulewidth}{0pt}
\renewcommand{\footrulewidth}{0pt}
\lhead{\raisebox{-0.15ex}{\poplogo\fontsize{13}{13}\selectfont\color{popink}OnBuch\color{poporange}+}}
\rhead{\poppill{\DOCMATIERE\ \textbullet\ \DOCNIVEAU\ \textbullet\ Leçon \DOCLECON}{white}{popink}}
\lfoot{\popxbold\fontsize{7.6}{7.6}\selectfont\color{popmuted}\MakeUppercase{Cours signé OnBuch+}\ \textbullet\ {\popsemi\DOCTITRE}}
\rfoot{\tikz[baseline=-0.6ex]\node[rounded corners=8.5pt,fill=popink,minimum height=17pt,minimum width=17pt,inner xsep=5pt,inner sep=0]{\popxbold\fontsize{8}{8}\selectfont\color{popcream}\thepage\,/\,\pageref*{LastPage}};}

% ============================================================
% Titres
% ============================================================
\newcommand{\chaptitre}[2]{%
  \begin{tcolorbox}[enhanced,colback=poporange,colframe=popink,boxrule=2.6pt,arc=11pt,
      drop shadow=popink,left=15pt,right=15pt,top=12pt,bottom=11pt,before skip=3pt,after skip=18pt]
    \poppill{#2}{popink}{popcream}\par\vspace{9pt}
    {\raggedright\hyphenpenalty=10000\exhyphenpenalty=10000\popdisplay\fontsize{22}{24}\selectfont\color{popcream}\MakeUppercase{#1}\par}
  \end{tcolorbox}
}
% Section numérotée : pastille orange avec le numéro + titre Archivo
\newcounter{popsec}
\newcommand{\coursec}[1]{%
  \stepcounter{popsec}\setcounter{popsub}{0}%
  \par\vspace{16pt}\noindent
  \begin{minipage}{\linewidth}
  \tikz[baseline=-0.7ex]\node[draw=popink,line width=1.6pt,fill=poporange,rounded corners=4pt,minimum size=22pt,inner sep=0]{\popdisplay\fontsize{12}{12}\selectfont\color{popcream}\thepopsec};%
  \hspace{8pt}{\popdisplay\fontsize{15}{17}\selectfont\color{popink}\MakeUppercase{#1}}\par
  \vspace{4pt}\noindent{\color{poporange}\rule{36pt}{3.6pt}}
  \end{minipage}\par\nopagebreak\vspace{8pt}
}
\newcounter{popsub}
\newcommand{\courssub}[1]{%
  \stepcounter{popsub}%
  \par\vspace{9pt}\noindent{\popxbold\fontsize{11.5}{13}\selectfont\color{popink}{\color{poporange}\thepopsec.\thepopsub}\hspace{6pt}#1}\par\nopagebreak\vspace{4pt}
}

% ============================================================
% Boîtes pédagogiques — même recette : fond blanc, bordure épaisse colorée,
% ombre dure encre, badge plein. Titre optionnel [#1].
% ============================================================
\tcbset{popbase/.style={breakable,enhanced,colback=white,boxrule=2.3pt,arc=9pt,
  shadow={3pt}{-3pt}{0pt}{popink},left=14pt,right=14pt,top=11pt,bottom=11pt,before skip=11pt,after skip=15pt,
  fontupper=\popsemi\fontsize{10.6}{14.5}\selectfont\color{popink},
  fontlower=\popsemi\fontsize{10.6}{14.5}\selectfont\color{popink}}}
% Coche dessinée (le glyphe ✓ n'existe pas dans les polices du document)
\newcommand{\popcheck}[1]{\tikz[baseline=-0.5ex]\draw[#1,line width=1.6pt,line cap=round,line join=round](-0.13,0.0)--(-0.04,-0.09)--(0.13,0.1);}
\providecommand{\coche}{\popcheck{popgreen}}
\providecommand{\resultat}[1]{\par\smallskip\noindent\fcolorbox{popgreen}{popgreenL}{\parbox{\dimexpr\linewidth-2\fboxsep-2\fboxrule\relax}{\textbf{Résultat.} #1}}\par\smallskip}
\newcommand{\popboxhead}[3]{\poppill{#1}{#2}{white}\ifx\relax#3\relax\else\hspace{6pt}{\popxbold\fontsize{10.6}{12}\selectfont\color{popink}#3}\fi\par\vspace{7pt}}

\newtcolorbox{aretenir}[1][]{popbase,colframe=popgreen,before upper={\popboxhead{À retenir}{popgreen}{#1}}}
\newtcolorbox{exemplebox}[1][]{popbase,colframe=poporange,before upper={\popboxhead{Exemple}{poporange}{#1}}}
\newtcolorbox{definition}[1][]{popbase,colframe=popblue,before upper={\popboxhead{Définition}{popblue}{#1}}}
\newtcolorbox{propriete}[1][]{popbase,colframe=poppurple,before upper={\popboxhead{Propriété}{poppurple}{#1}}}
\newtcolorbox{methode}[1][]{popbase,colframe=popgold,before upper={\popboxhead{Méthode}{popgold}{#1}}}
\newtcolorbox{attention}[1][]{popbase,colframe=poppink,before upper={\popboxhead{Attention}{poppink}{#1}}}
\newtcolorbox{experience}[1][]{popbase,colframe=popblue,colback=popblueL,before upper={\popboxhead{Expérience}{popblue}{#1}}}
% « Le savais-tu ? » : carte violette pleine (contexte, culture, Cameroun)
\newtcolorbox{savaistu}[1][]{popbase,colback=poppurple,colframe=popink,
  fontupper=\popsemi\fontsize{10.4}{14.2}\selectfont\color{white},
  before upper={\poppill{Le savais-tu\,?}{white}{poppurple}\ifx\relax#1\relax\else\hspace{6pt}{\popxbold\color{white}#1}\fi\par\vspace{7pt}}}
% Exercice résolu : énoncé en haut, \tcblower puis la solution sur fond crème
\newcounter{popexo}\newcounter{popexr}
\newtcolorbox{exoresolu}[1][]{popbase,colframe=poporange,colbacklower=popcream,
  segmentation style={popink,line width=1.2pt,dashed},
  before upper={\stepcounter{popexr}\popboxhead{Exercice résolu \thepopexr}{poporange}{#1}},
  before lower={\poppill{Solution}{popink}{popcream}\par\vspace{6pt}}}
% Exercice (non résolu en place) — la correction est dans la partie Corrigés
\newtcolorbox{exercice}[1][]{popbase,colframe=popink,
  before upper={\stepcounter{popexo}\popboxhead{Exercice \thepopexo}{popink}{#1}}}
% Corrigé
\newtcolorbox{corrige}[1][]{popbase,colframe=popgreen,colback=popgreenL,
  before upper={\popboxhead{Corrigé}{popgreen}{#1}}}
% Objectifs / prérequis / bilan
\newtcolorbox{objectifs}{popbase,colframe=popink,colback=poporangeL,
  before upper={\popboxhead{Objectifs de la leçon}{popink}{}}}
\newtcolorbox{prerequis}{popbase,colframe=popink,colback=popblueL,
  before upper={\popboxhead{Prérequis}{popblue}{}}}
\newtcolorbox{bilan}{popbase,colframe=popink,colback=white,boxrule=2.8pt,
  before upper={\popboxhead{Fiche bilan}{poporange}{L'essentiel à connaître pour l'examen}}}
% Figure encadrée avec légende
\newtcolorbox{popfigure}[1][]{enhanced,breakable=false,colback=white,colframe=popline,boxrule=1.4pt,arc=8pt,
  left=10pt,right=10pt,top=10pt,bottom=8pt,before skip=10pt,after skip=12pt,halign=center,
  lower separated=false,halign lower=center,
  fontlower=\popsemi\fontsize{9}{11.5}\selectfont\color{popmuted},
  #1}
\newcommand{\legende}[1]{\par\vspace{6pt}{\popsemi\fontsize{9}{11.5}\selectfont\color{popmuted}#1}}

% ============================================================
% Signature OnBuch+
% ============================================================
\newcommand{\poplogoinline}[1]{{\poplogo\fontsize{#1}{#1}\selectfont\color{popink}OnBuch\color{poporange}+}}
\newcommand{\signatureonbuch}{%
  \begin{tcolorbox}[enhanced,colback=popink,colframe=popink,boxrule=2.2pt,arc=10pt,
      drop shadow=poporange,left=14pt,right=14pt,top=10pt,bottom=10pt,before skip=0pt,after skip=0pt]
    \tikz[baseline=-0.7ex]\node[circle,fill=popgreen,draw=popcream,line width=1.3pt,minimum size=22pt,inner sep=0]{\popcheck{white}};%
    \hspace{9pt}\begin{minipage}[c]{0.62\linewidth}
      {\popxbold\fontsize{12}{13}\selectfont\color{popcream}Signé\ \textbullet\ {\poplogo OnBuch\color{poporange}+}}\par\vspace{2pt}
      {\raggedright\popsemi\fontsize{8}{10.5}\selectfont\color{popline}\MakeUppercase{Équipe pédagogique OnBuch+}\\\MakeUppercase{Conforme au programme officiel MINESEC}\par}
    \end{minipage}\hfill
    {\poplogo\fontsize{22}{22}\selectfont\color{popcream}OnBuch\color{poporange}+}
  \end{tcolorbox}%
}

% ============================================================
% Page de garde
% #1 titre · #2 matière · #3 niveau · #4 module · #5 n° leçon · #6 durée ·
% #7 description / objectifs (texte ou itemize)
% ============================================================
\newcommand{\pagegarde}[7]{%
  \thispagestyle{empty}
  \newgeometry{a4paper,margin=1.7cm,top=1.6cm,bottom=1.4cm}%
  \begin{tikzpicture}[remember picture,overlay]
    \fill[poporange!18] ([xshift=-3.2cm,yshift=-3.6cm]current page.north east) circle (4.6cm);
    \fill[poppurple!14] ([xshift=2.4cm,yshift=7.6cm]current page.south west) circle (3.8cm);
  \end{tikzpicture}%
  {\poplogo\fontsize{30}{30}\selectfont\color{popink}OnBuch\color{poporange}+}\hfill
  \raisebox{0.6ex}{\poppill{Cours signé \textbullet\ Premium}{popink}{popcream}}\par
  \vspace{1pt}\popsquiggle{poppurple}\par
  \vspace{8pt}
  \begin{tcolorbox}[enhanced,colback=poporange,colframe=popink,boxrule=2.8pt,arc=13pt,
      drop shadow=popink,left=18pt,right=18pt,top=12pt,bottom=13pt,after skip=10pt]
    \poppill{#2 \textbullet\ #3}{popink}{popcream}\hspace{5pt}\poppill{#4}{white}{popink}\par\vspace{8pt}
    {\popdisplay\fontsize{14}{14}\selectfont\color{popink}LEÇON #5}\par\vspace{4pt}
    {\raggedright\hyphenpenalty=10000\exhyphenpenalty=10000\popdisplay\fontsize{25}{27}\selectfont\color{popcream}\MakeUppercase{#1}\par}
    \vspace{8pt}
    {\popxbold\fontsize{9.5}{9.5}\selectfont\color{popink}\MakeUppercase{Durée conseillée : #6}}
  \end{tcolorbox}
  \begin{tcolorbox}[enhanced,colback=white,colframe=popink,boxrule=2.2pt,arc=10pt,
      drop shadow=popink,left=16pt,right=16pt,top=10pt,bottom=10pt,after skip=8pt]
    \poppill{Dans cette leçon}{poppurple}{white}\par\vspace{5pt}
    \popsemi\fontsize{9.3}{12.6}\selectfont\color{popink}#7
  \end{tcolorbox}
  \noindent\begin{minipage}[t]{0.487\linewidth}
    \begin{tcolorbox}[enhanced,colback=popgreen,colframe=popink,boxrule=2.2pt,arc=10pt,
        drop shadow=popink,left=11pt,right=11pt,top=10pt,bottom=10pt,height=3.1cm,valign=center]
      \tcbox[colback=white,colframe=popink,boxrule=1.3pt,arc=5pt,boxsep=0pt,nobeforeafter,
        left=3pt,right=3pt,top=3pt,bottom=3pt,valign=center]{\qrcode[height=1.9cm]{https://play.google.com/store/apps/details?id=com.luvvix.onbuch&hl=fr}}%
      \hspace{8pt}\begin{minipage}[c]{0.5\linewidth}
        {\popdisplay\fontsize{11}{12.5}\selectfont\color{white}\MakeUppercase{Télécharge OnBuch}}\par\vspace{4pt}
        {\raggedright\popsemi\fontsize{8.4}{11}\selectfont\color{white}Gratuit \textbullet\ Google Play\\Tuteur IA, annales, concours, résultats\par}
      \end{minipage}
    \end{tcolorbox}
  \end{minipage}\hfill
  \begin{minipage}[t]{0.487\linewidth}
    \begin{tcolorbox}[enhanced,colback=poppurple,colframe=popink,boxrule=2.2pt,arc=10pt,
        drop shadow=popink,left=11pt,right=11pt,top=10pt,bottom=10pt,height=3.1cm,valign=center]
      \tikz[baseline=-0.7ex]\node[circle,fill=white,draw=popink,line width=1.3pt,minimum size=1.6cm,inner sep=0]{\popdisplay\fontsize{24}{24}\selectfont\color{poppurple}+};%
      \hspace{8pt}\begin{minipage}[c]{0.6\linewidth}
        {\popdisplay\fontsize{11}{12.5}\selectfont\color{white}\MakeUppercase{Passe à OnBuch+}}\par\vspace{4pt}
        {\raggedright\popsemi\fontsize{8.4}{11}\selectfont\color{white}Tous les cours signés, Tuteur IA illimité, fiches PDF — onglet OnBuch+ de l'app\par}
      \end{minipage}
    \end{tcolorbox}
  \end{minipage}
  \vspace{8pt}
  \popcontenu
  \vfill
  \signatureonbuch
  \restoregeometry
  \newpage
}

% Rangée « Ce cours contient » (page de garde)
\newcommand{\poptile}[3]{% #1 couleur #2 gros texte #3 légende
  \begin{tcolorbox}[enhanced,colback=#1,colframe=popink,boxrule=2pt,arc=9pt,shadow={2.5pt}{-2.5pt}{0pt}{popink},
      left=6pt,right=6pt,top=9pt,bottom=9pt,halign=center,valign=center,height=1.75cm,width=\linewidth,nobeforeafter]
    {\popdisplay\fontsize{12.5}{13}\selectfont\color{white}\MakeUppercase{#2}}\par\vspace{3pt}
    {\centering\popsemi\fontsize{7.8}{9.5}\selectfont\color{white}#3\par}
  \end{tcolorbox}}
\newcommand{\popcontenu}{%
  \par\noindent{\popxboldls\fontsize{7.5}{7.5}\selectfont\color{popmuted}CE COURS SIGNÉ CONTIENT}\par\vspace{7pt}
  \noindent\begin{minipage}[t]{0.235\linewidth}\poptile{popblue}{Cours}{complet, illustré, pas à pas}\end{minipage}\hfill
  \begin{minipage}[t]{0.235\linewidth}\poptile{poporange}{Méthodes}{et exercices résolus}\end{minipage}\hfill
  \begin{minipage}[t]{0.235\linewidth}\poptile{poppink}{Exercices}{type examen + corrigés}\end{minipage}\hfill
  \begin{minipage}[t]{0.235\linewidth}\poptile{popgreen}{Fiche bilan}{l'essentiel à retenir}\end{minipage}\par}

% Sommaire
\newtcolorbox{sommaire}{enhanced,colback=white,colframe=popink,boxrule=2.2pt,arc=10pt,
  drop shadow=popink,left=18pt,right=18pt,top=13pt,bottom=13pt,before skip=6pt,after skip=20pt,
  before upper={\poppill{Sommaire}{poppurple}{white}\par\vspace{9pt}\popsemi\fontsize{10.6}{14.5}\selectfont\color{popink}}}

% Signature de fin de cours — poussée tout en bas de la dernière page
\newcommand{\findecours}{%
  \par\vfill\nopagebreak
  \begin{center}\popsquiggle{poporange}\end{center}\vspace{4pt}
  \signatureonbuch
}
\AtBeginDocument{\def\llbracket{\mathopen{[\mkern-3.5mu[}}\def\rrbracket{\mathclose{]\mkern-3.5mu]}}} % intervalles d'entiers (glyphe ⟦ absent de la police)
% Glyphes absents de Fira Math : redéfinis à partir de glyphes disponibles
\AtBeginDocument{%
  \def\setminus{\mathbin{\backslash}}\let\smallsetminus\setminus
  \def\vdots{\mathord{\vbox{\baselineskip3.2pt\lineskiplimit0pt\kern4pt\hbox{.}\hbox{.}\hbox{.}}}}%
  \def\ddots{\mathinner{\mkern1mu\raise6.5pt\hbox{.}\mkern2mu\raise3.6pt\hbox{.}\mkern2mu\raise0.7pt\hbox{.}\mkern1mu}}%
  \def\triangle{\mathord{\scalebox{0.9}{$\symup{\Delta}$}}}%
  \def\nabla{\mathord{\raisebox{\depth}{\scalebox{1}[-1]{$\symup{\Delta}$}}}}%
  \def\checkmark{\popcheck{popdark}}%
  \def\star{\mathbin{\ast}}\def\bigstar{\mathord{\ast}}%
  \def\diamond{\mathbin{\scalebox{0.6}{$\Diamond$}}}%
  \def\bigcup{\mathop{\mathchoice{\raisebox{-0.3em}{\scalebox{1.7}{$\cup$}}}{\scalebox{1.25}{$\cup$}}{\cup}{\cup}}\displaylimits}%
  \def\bigcap{\mathop{\mathchoice{\raisebox{-0.3em}{\scalebox{1.7}{$\cap$}}}{\scalebox{1.25}{$\cap$}}{\cap}{\cap}}\displaylimits}%
  \def\lesseqqgtr{\mathrel{\lessgtr}}\def\bigoplus{\mathop{\oplus}\displaylimits}%
}
\providecommand{\ssi}{si et seulement si} % abréviation fréquente
\AtBeginDocument{\providecommand{\wideparen}{\overparen}} % arc de cercle

%% FIGFIX-BEGIN v5 — correctifs globaux des figures (docs/onbuch-generation/tools/figfix)
\usepackage{adjustbox}\usepackage{etoolbox}\tcbuselibrary{raster}
% toute figure TikZ/pgfplots plus large que la ligne est réduite à la largeur disponible
\BeforeBeginEnvironment{tikzpicture}{\begin{adjustbox}{max width=\linewidth}}
\AfterEndEnvironment{tikzpicture}{\end{adjustbox}}
% légendes pgfplots lisibles par-dessus les courbes
\pgfplotsset{every axis/.append style={legend style={fill=white,fill opacity=0.92,text opacity=1,draw=black!15,inner sep=3pt}}}
% étiquettes d'arêtes (syntaxe edge["..."]) avec fond blanc
\tikzset{every edge quotes/.append style={fill=white,inner sep=1.2pt}}
%% FIGFIX-END
\begin{document}

\pagegarde{Lutte contre les problèmes liés à la régulation de la glycémie et à la pression artérielle}{SVT}{Tle D}{Module II : L'Éducation à la Santé — La santé nutritionnelle}{N13}{12 h}{Cette leçon étudie les mécanismes de régulation de la glycémie et de la pression artérielle, établis à partir de faits expérimentaux, puis les pathologies liées à leur dérèglement (diabète, hypertension et hypotension artérielles) et les moyens de lutte. Elle constitue le cœur du module d'Éducation à la Santé consacré à la santé nutritionnelle en Terminale D.
\vspace{4pt}
\begin{multicols}{2}\small
\begin{itemize}
\item À la fin de cette leçon, je sais expliquer les techniques de mesure de la glycémie et de la pression artérielle
\item À la fin de cette leçon, je sais analyser des courbes de variation de la glycémie avant et après un repas (hyperglycémie provoquée)
\item À la fin de cette leçon, je sais interpréter des résultats expérimentaux mettant en évidence les rôles du foie et du pancréas dans la régulation de la glycémie
\item À la fin de cette leçon, je sais réaliser un schéma fonctionnel de la régulation hormonale et nerveuse de la glycémie
\item À la fin de cette leçon, je sais exploiter des résultats d'analyses médicales pour identifier un diabète et préciser son type
\item À la fin de cette leçon, je sais analyser des données et courbes illustrant les variations de la pression artérielle
\end{itemize}\end{multicols}}

\begin{sommaire}
\begin{multicols}{2}
\begin{enumerate}
\item La glycémie : mesure et variations
\item Régulation de la glycémie : rôles du foie et du pancréas
\item Le diabète : typologie, causes et moyens de lutte
\item La pression artérielle : mesure et variations
\item Régulation de la pression artérielle
\item Hypertension et hypotension artérielles : lutte et sensibilisation
\item Activité d'intégration
\item Méthodes et exercices résolus
\item Exercices
\item Corrigés des exercices
\item Fiche bilan
\end{enumerate}
\end{multicols}
\end{sommaire}

\begin{objectifs}
\begin{itemize}
\item À la fin de cette leçon, je sais expliquer les techniques de mesure de la glycémie et de la pression artérielle
\item À la fin de cette leçon, je sais analyser des courbes de variation de la glycémie avant et après un repas (hyperglycémie provoquée)
\item À la fin de cette leçon, je sais interpréter des résultats expérimentaux mettant en évidence les rôles du foie et du pancréas dans la régulation de la glycémie
\item À la fin de cette leçon, je sais réaliser un schéma fonctionnel de la régulation hormonale et nerveuse de la glycémie
\item À la fin de cette leçon, je sais exploiter des résultats d'analyses médicales pour identifier un diabète et préciser son type
\item À la fin de cette leçon, je sais analyser des données et courbes illustrant les variations de la pression artérielle
\item À la fin de cette leçon, je sais identifier les causes, symptômes et complications de l'hypertension et de l'hypotension artérielles
\item À la fin de cette leçon, je sais élaborer des outils de sensibilisation à la lutte contre le diabète et l'hypertension
\end{itemize}
\end{objectifs}

\begin{prerequis}
\begin{itemize}
\item Notion d'homéostasie et de milieu intérieur (constance des paramètres du sang)
\item Nutrition et digestion des glucides (Première : devenir du glucose absorbé)
\item Notion d'hormone et de glande endocrine (exemple de l'adrénaline, Première)
\item Rôle du système nerveux dans la régulation des fonctions vitales (leçons précédentes de Terminale)
\item Circulation sanguine : cœur, artères, veines et notion de débit cardiaque
\end{itemize}
\end{prerequis}

\newpage

\coursec{La glycémie : mesure et variations}

\textbf{Situation d'accroche.} Il est 7 h du matin au marché Mokolo, à Yaoundé. Mama Ngo Bell, 52 ans, vend des légumes depuis vingt ans. Depuis quelques mois, elle se sent épuisée, a soif en permanence et se lève plusieurs fois par nuit pour uriner. Inquiète, elle se rend à l'hôpital de district : on lui mesure une glycémie à jeun de \SI{1,8}{g/L} et une tension artérielle de 16/10. Le diagnostic tombe : diabète et hypertension artérielle. Quelques étals plus loin, son voisin, jeune commerçant de 30 ans, s'évanouit régulièrement en fin de matinée : sa tension est de 9/6. Deux malades, deux dérèglements opposés, mais une même question de fond : comment le corps humain maintient-il normalement le taux de sucre dans le sang et la pression artérielle autour de valeurs constantes ? Pourquoi ces régulations peuvent-elles défaillir ? Et comment lutter contre les maladies qui en résultent, de plus en plus fréquentes au Cameroun avec la sédentarité et les changements alimentaires ? Pour répondre à ces questions, commençons par la plus fondamentale des variables régulées de l'organisme : la \cle{glycémie}.

\courssub{La glycémie : définition et valeur normale}

\begin{definition}[La glycémie]
La \cle{glycémie} est la \textbf{concentration de glucose dans le sang}. Elle s'exprime en grammes par litre de sang (\SI{}{g/L}) ou en millimoles par litre (\SI{}{mmol/L}). Chez un sujet sain, à jeun, elle est voisine de \SI{1}{g/L}, comprise entre \SI{0,7}{g/L} et \SI{1,1}{g/L} (soit environ 3,9 à \SI{6,1}{mmol/L}).
\end{definition}

Pourquoi le glucose est-il si important ? Parce qu'il est le \textbf{carburant énergétique privilégié des cellules}, et le seul utilisable en permanence par certains organes comme le cerveau. La respiration cellulaire le dégrade pour produire de l'ATP selon le bilan :
\[ \ce{C6H12O6 + 6O2 -> 6CO2 + 6H2O} + \text{énergie (ATP)} \]
Une chute prolongée de la glycémie prive donc le cerveau d'énergie : c'est le malaise hypoglycémique (sueurs, tremblements, troubles de la conscience, voire coma). À l'inverse, une élévation chronique endommage les vaisseaux sanguins et les nerfs. L'organisme doit donc maintenir la glycémie dans une fourchette étroite : c'est un impératif de survie.

\begin{exemplebox}[Valeurs de référence à connaître]
\begin{itemize}
\item Glycémie à jeun normale : \SI{0,7}{g/L} à \SI{1,1}{g/L}.
\item Glycémie à jeun comprise entre \SI{1,1}{g/L} et \SI{1,26}{g/L} : hyperglycémie modérée à jeun (état pré-diabétique).
\item Glycémie à jeun supérieure ou égale à \SI{1,26}{g/L}, vérifiée à deux reprises : on parle de \textbf{diabète}.
\item Seuil rénal : environ \SI{1,8}{g/L} ; au-delà, du glucose apparaît dans les urines.
\end{itemize}
\end{exemplebox}

\begin{savaistu}[Un seul morceau de sucre dans tout le sang]
Le sang d'un adulte représente environ 5 litres. À \SI{1}{g/L}, cela correspond à seulement \textbf{5 grammes de glucose} en permanence dans tout l'organisme, soit l'équivalent d'un morceau de sucre ! Pourtant, un repas peut apporter 50 à \SI{100}{g} de glucides. Sans mécanisme de régulation, la glycémie serait multipliée par 10 à 20 après chaque repas. Le maintien de ce « morceau de sucre » constant est un exploit physiologique que nous étudierons dans la section suivante.
\end{savaistu}

\courssub{Les techniques de mesure de la glycémie}

\textbf{1. Le dosage sanguin en laboratoire.} C'est la méthode de référence. On prélève du sang veineux (au pli du coude) chez un sujet \textbf{à jeun depuis au moins 8 heures}, puis on dose le glucose plasmatique par des méthodes enzymatiques (utilisant la glucose-oxydase, enzyme qui oxyde spécifiquement le glucose). C'est ce dosage qui permet le diagnostic officiel du diabète.

\textbf{2. Le lecteur de glycémie (glucomètre).} Appareil portable très répandu dans les pharmacies et les services de diabétologie camerounais. Le principe est simple : on pique le bout du doigt avec une lancette, on dépose une goutte de sang capillaire sur une \textbf{bandelette réactive} imprégnée d'une enzyme (glucose-oxydase), et l'appareil affiche la glycémie en quelques secondes. Il permet l'autosurveillance quotidienne du diabétique, mais ne suffit pas à poser le diagnostic.

\begin{popfigure}
\begin{center}
\begin{tikzpicture}[font=\small]
% Doigt
\draw[fill=poppinkL, rounded corners=6pt] (0,0) rectangle (1.6,3.2);
\draw[fill=poppink!40] (0.8,3.2) circle (0.8);
\draw[popdark] (0,0) rectangle (1.6,3.2);
% goutte de sang
\draw[fill=popink] (0.8,4.35) circle (0.14);
\node[above, popink] at (0.8,4.5) {goutte de sang};
% lancette
\draw[fill=popmuted!60] (-1.6,4.6) rectangle (-0.4,4.9);
\draw[fill=popdark] (-0.4,4.62) -- (0.15,4.75) -- (-0.4,4.88) -- cycle;
\node[above] at (-1.0,4.95) {lancette};
% flèche
\draw[->, very thick, popblue] (2.2,2.4) -- (3.6,2.4);
% bandelette
\draw[fill=popgoldL, rounded corners=2pt] (3.8,1.9) rectangle (6.6,2.9);
\draw[fill=popink!70] (3.8,1.9) rectangle (4.3,2.9);
\node[below, align=center] at (5.2,1.8) {bandelette réactive\\ (enzyme : glucose-oxydase)};
% flèche
\draw[->, very thick, popblue] (6.9,2.4) -- (8.3,2.4);
% glucomètre
\draw[fill=popblueL, rounded corners=8pt] (8.6,0.9) rectangle (11.4,4.1);
\draw[fill=white, rounded corners=2pt] (9.0,2.9) rectangle (11.0,3.7);
\node[popdark, font=\bfseries] at (10.0,3.3) {1,8 g/L};
\draw[fill=popmuted!50] (9.4,1.3) circle (0.22);
\draw[fill=popmuted!50] (10.6,1.3) circle (0.22);
\node[below, align=center] at (10.0,0.8) {lecteur de glycémie\\ (glucomètre)};
\end{tikzpicture}
\end{center}
\legende{Principe de la mesure de la glycémie au glucomètre : (1) piqûre du bout du doigt avec une lancette et recueil d'une goutte de sang capillaire ; (2) dépôt du sang sur la zone réactive de la bandelette, où le glucose réagit avec une enzyme ; (3) insertion de la bandelette dans le lecteur, qui affiche la glycémie en quelques secondes.}
\end{popfigure}

\courssub{Les variations de la glycémie : l'expérience de l'hyperglycémie provoquée}

La glycémie est-elle vraiment constante ? Pour le vérifier, réalisons l'expérience classique dite de l'\textbf{hyperglycémie provoquée par voie orale} (HGPO), aussi utilisée en médecine comme test de tolérance au glucose.

\begin{experience}[L'hyperglycémie provoquée par voie orale (HGPO)]
\textbf{Protocole.} Un sujet sain, à jeun depuis 12 heures, ingère \SI{75}{g} de glucose dissous dans de l'eau. On dose sa glycémie toutes les 30 minutes pendant 2 heures.

\textbf{Résultats obtenus chez un sujet sain :}
\begin{center}
\begin{tabularx}{\linewidth}{|l|X|X|X|X|X|}
\hline
\rowcolor{popblueL} Temps (min) & 0 (à jeun) & 30 & 60 & 90 & 120 \\
\hline
Glycémie (g/L) & 0,95 & 1,40 & 1,20 & 1,00 & 0,90 \\
\hline
\end{tabularx}
\end{center}

\textbf{Observations.} La glycémie s'élève rapidement après l'ingestion (pic vers 30 min, à environ \SI{1,4}{g/L}), puis diminue et revient à sa valeur initiale en environ 2 heures.

\textbf{Interprétation.} Pendant la première demi-heure, le glucose absorbé passe de l'intestin grêle dans le sang plus vite qu'il n'est capté par les cellules : la glycémie augmente. Ensuite, le glucose est rapidement capté, stocké ou utilisé par les cellules : la glycémie diminue. L'organisme possède donc un mécanisme capable de faire redescendre la glycémie après un apport : c'est la preuve expérimentale d'une \textbf{régulation}.
\end{experience}

\begin{popfigure}
\begin{center}
\begin{tikzpicture}
\begin{axis}[
width=0.85\linewidth, height=6.5cm,
xlabel={Temps (min)}, ylabel={Glycémie (g/L)},
xmin=0, xmax=150, ymin=0, ymax=2,
xtick={0,30,60,90,120,150}, ytick={0,0.5,1,1.5,2},
grid=both, grid style={popline!40},
legend style={at={(0.97,0.97)}, anchor=north east, font=\small}
]
\addplot[popcurve, mark=*, mark size=1.5pt] coordinates {(0,0.95) (30,1.4) (60,1.2) (90,1.0) (120,0.9) (150,0.95)};
\addlegendentry{Sujet sain}
\addplot[domain=0:150, samples=2, dashed, popgreen, thick] {1};
\addlegendentry{Valeur normale à jeun (\SI{1}{g/L})}
\addplot[domain=0:150, samples=2, dotted, popink, thick] {1.8};
\addlegendentry{Seuil rénal (\SI{1,8}{g/L})}
\end{axis}
\end{tikzpicture}
\end{center}
\legende{Courbe d'hyperglycémie provoquée (HGPO) chez un sujet sain : après ingestion de \SI{75}{g} de glucose, la glycémie culmine vers 30 min sans jamais atteindre le seuil rénal, puis revient à la normale en 2 heures environ.}
\end{popfigure}

\begin{methode}[Exploiter une courbe d'hyperglycémie provoquée]
Face à une courbe de glycémie en fonction du temps, procède toujours en quatre étapes :
\begin{enumerate}
\item \textbf{Repère la valeur à jeun} (temps 0) : est-elle normale (0,7 à \SI{1,1}{g/L}) ?
\item \textbf{Repère le pic} : sa valeur et le temps auquel il survient (chez le sujet sain, vers 30 min et moins de \SI{1,8}{g/L}).
\item \textbf{Mesure la durée du retour à la normale} : environ 2 h chez le sujet sain ; plus longue ou incomplète chez le diabétique.
\item \textbf{Conclus} en comparant à la courbe de référence : une glycémie qui reste élevée après 2 h traduit un défaut de régulation.
\end{enumerate}
\end{methode}

\courssub{Les causes des variations de la glycémie}

La glycémie résulte d'un équilibre dynamique entre deux flux opposés :

\begin{itemize}
\item \textbf{Les apports}, qui font \emph{monter} la glycémie : la digestion des glucides alimentaires (amidon du foutou, du riz, du manioc, sucres des boissons sucrées) libère du glucose qui passe dans le sang au niveau de l'intestin grêle.
\item \textbf{Les consommations et stockages}, qui font \emph{descendre} la glycémie : la respiration cellulaire (toutes les cellules), l'effort musculaire (forte consommation de glucose par les muscles), et le stockage du glucose sous forme de glycogène dans le foie et les muscles.
\end{itemize}

Pourtant, malgré des repas discontinus (apports par « à-coups ») et des efforts variables, la glycémie reste remarquablement stable autour de \SI{1}{g/L}. Cette \textbf{constance d'une variable soumise à des perturbations permanentes} est la signature d'un mécanisme de régulation : la glycémie est une \cle{variable régulée}. Nous verrons dans la section suivante que le foie et le pancréas en sont les organes clés.

\courssub{Conséquence d'une élévation excessive : la glycosurie}

Le rein filtre le sang au niveau des glomérules : le glucose du plasma passe intégralement dans l'urine primitive, puis il est normalement \textbf{réabsorbé en totalité} au niveau des tubules rénaux. Mais cette réabsorption a une capacité maximale : le \cle{seuil rénal}, situé vers \SI{1,8}{g/L}. Si la glycémie dépasse durablement ce seuil, les reins ne peuvent plus tout réabsorber : du glucose apparaît dans les urines définitives. C'est la \cle{glycosurie}. C'est précisément le cas de Mama Ngo Bell : sa glycémie à jeun de \SI{1,8}{g/L} atteint le seuil rénal, ce qui explique ses urines abondantes (le glucose non réabsorbé entraîne de l'eau avec lui : c'est la polyurie) et sa soif permanente (polydipsie, conséquence de la déshydratation).

\begin{attention}[Ne pas confondre glycémie et glycosurie]
Erreur très fréquente au Baccalauréat ! La \textbf{glycémie} est le \emph{taux de glucose dans le sang} (en g/L) : elle existe chez tout le monde, en permanence. La \textbf{glycosurie} est la \emph{présence de glucose dans les urines} : elle est \emph{anormale} et n'apparaît que lorsque la glycémie dépasse le seuil rénal (environ \SI{1,8}{g/L}). Une glycémie élevée n'entraîne une glycosurie qu'au-delà de ce seuil.
\end{attention}

\begin{exoresolu}[Exploitation d'une courbe d'hyperglycémie provoquée]
\textbf{Énoncé.} Chez deux sujets A et B à jeun, on mesure la glycémie après ingestion de \SI{75}{g} de glucose :

\begin{center}
\begin{tabularx}{\linewidth}{|l|X|X|X|X|X|}
\hline
\rowcolor{popblueL} Temps (min) & 0 & 30 & 60 & 90 & 120 \\
\hline
Sujet A (g/L) & 0,90 & 1,35 & 1,15 & 1,00 & 0,90 \\
\hline
Sujet B (g/L) & 1,50 & 2,10 & 2,00 & 1,90 & 1,80 \\
\hline
\end{tabularx}
\end{center}

1. Identifier le sujet sain et le sujet diabétique. Justifier. \quad 2. Le sujet B présente-t-il une glycosurie ? Justifier.
\tcblower
\textbf{Solution détaillée.}

\textbf{1.} Le sujet A a une glycémie à jeun de \SI{0,90}{g/L}, dans la fourchette normale (0,7 à \SI{1,1}{g/L}) ; sa glycémie culmine à \SI{1,35}{g/L} à 30 min puis revient à la normale en 2 h : c'est le \textbf{sujet sain}. Le sujet B a une glycémie à jeun de \SI{1,50}{g/L}, supérieure au seuil diagnostique du diabète (\SI{1,26}{g/L}) ; de plus, sa glycémie reste supérieure ou égale à \SI{1,8}{g/L} au bout de 2 h : le retour à la normale ne s'effectue pas. C'est le \textbf{sujet diabétique}.

\textbf{2.} La glycémie du sujet B dépasse le seuil rénal (\SI{1,8}{g/L}) entre 30 et 120 min : les reins ne peuvent plus réabsorber tout le glucose filtré. Le sujet B présente donc une \textbf{glycosurie} pendant cette période.
\end{exoresolu}

\begin{aretenir}[L'essentiel de la section]
\begin{itemize}
\item La glycémie est la concentration de glucose dans le sang ; elle vaut environ \SI{1}{g/L} à jeun (0,7 à \SI{1,1}{g/L}).
\item Elle se mesure en laboratoire (sang veineux à jeun) ou au glucomètre (goutte de sang capillaire sur bandelette).
\item L'hyperglycémie provoquée montre que, chez le sujet sain, la glycémie monte après un apport de glucose puis revient à la normale en 2 h : preuve d'une régulation.
\item La glycémie varie selon les apports (digestion des glucides) et les consommations (respiration cellulaire, effort, stockage), mais reste constante : c'est une variable régulée.
\item Au-delà du seuil rénal (\SI{1,8}{g/L}), le glucose passe dans les urines : c'est la glycosurie, signe d'alerte du diabète.
\end{itemize}
\end{aretenir}

\begin{exercice}[Pour t'entraîner]
Un patient se présente au centre de santé avec une glycémie capillaire de \SI{1,3}{g/L} mesurée deux heures après un repas riche en foutou et en jus sucré. Peut-on conclure qu'il est diabétique ? Justifie ta réponse en précisant les conditions rigoureuses du diagnostic.
\end{exercice}

\begin{corrige}[Exercice]
Non, on ne peut pas conclure au diabète. Deux heures après un repas riche en glucides, une glycémie de \SI{1,3}{g/L} peut correspondre à une hyperglycémie post-prandiale normale (le pic physiologique peut atteindre \SI{1,4}{g/L} chez le sujet sain). Le diagnostic de diabète exige une glycémie \textbf{à jeun} (au moins 8 h sans apport) supérieure ou égale à \SI{1,26}{g/L}, vérifiée \textbf{à deux reprises}, idéalement par dosage en laboratoire sur sang veineux. Il faut donc refaire la mesure dans les conditions réglementaires avant toute conclusion.
\end{corrige}

\coursec{Régulation de la glycémie : rôles du foie et du pancréas}

Dans la section précédente, nous avons vu que la glycémie d'un individu sain reste remarquablement stable, autour de \SI{1}{g/L} (soit environ \SI{5,5}{mmol/L}) à jeun, malgré les apports discontinus de glucose lors des repas et les dépenses énergétiques variables de l'organisme. Après un repas riche en glucides, la glycémie s'élève transitoirement (jusqu'à \SI{1,4}{g/L} environ), puis revient à sa valeur de référence en moins de deux heures. Cette constance n'est pas un hasard : elle est le résultat d'une \cle{régulation} active, qui fait intervenir deux organes vedettes — le \cle{pancréas} et le \cle{foie} — sous le contrôle du système nerveux. Comment la science a-t-elle découvert ces rôles ? C'est l'une des plus belles histoires de la physiologie, et c'est elle que nous allons revivre ici, en véritables investigateurs.

\courssub{Mise en évidence expérimentale du rôle du pancréas}

\textbf{Observation de départ.} Au XIX\textsuperscript{e} siècle, les médecins connaissent le diabète sucré : les malades ont une soif intense, urinent abondamment une urine sucrée (\cle{glycosurie}), maigrissent et meurent jeunes. Mais quel organe est en cause ?

\begin{experience}[La pancréatectomie de von Mering et Minkowski (1889)]
\textbf{Hypothèse :} le pancréas est impliqué dans le maintien d'une glycémie normale.

\textbf{Protocole :} les deux physiologistes pratiquent chez des chiens une \cle{pancréatectomie} totale (ablation chirurgicale du pancréas), puis suivent la glycémie et la présence de glucose dans les urines. Des chiens témoins subissent une opération fictive (on ouvre l'abdomen sans toucher le pancréas).

\textbf{Résultats :}
\begin{itemize}
\item les chiens pancréatectomisés développent une \cle{hyperglycémie} permanente (glycémie supérieure à \SI{2}{g/L}), une glycosurie importante, une soif intense (\cle{polydipsie}), une émission abondante d'urines (\cle{polyurie}), un amaigrissement, puis meurent en quelques semaines ;
\item les chiens témoins gardent une glycémie normale (environ \SI{1}{g/L}).
\end{itemize}

\textbf{Interprétation :} la seule variable expérimentale est la présence ou l'absence du pancréas. Le pancréas est donc \emph{nécessaire} au maintien d'une glycémie normale : il doit produire une substance hypoglycémiante.

\textbf{Expérience complémentaire :} la greffe d'un fragment de pancréas sous la peau d'un chien pancréatectomisé fait disparaître les symptômes. En 1921, Banting et Best extraient du pancréas une substance — l'\cle{insuline} — dont l'injection à des chiens rendus diabétiques normalise leur glycémie en quelques heures.
\end{experience}

\begin{center}
\begin{tabularx}{\linewidth}{|l|X|X|X|}
\hline
\rowcolor{popblueL} \textbf{Expérience} & \textbf{Glycémie avant} & \textbf{Glycémie après} & \textbf{Conclusion} \\
\hline
Pancréatectomie totale (chien) & \SI{1,0}{g/L} & $>$ \SI{2,0}{g/L} (hyperglycémie, glycosurie, mort) & Le pancréas est nécessaire à la régulation \\
\hline
Greffe de fragment pancréatique & $>$ \SI{2,0}{g/L} & $\approx$ \SI{1,0}{g/L} & Le pancréas suffit à rétablir la normale \\
\hline
Injection d'extrait pancréatique (insuline, 1921) & \SI{2,5}{g/L} & \SI{0,9}{g/L} en quelques heures & La substance active est l'insuline \\
\hline
\end{tabularx}
\end{center}

\begin{savaistu}[Un prix Nobel qui a sauvé des millions de vies]
Frederick Banting et John Macleod ont reçu le prix Nobel de physiologie ou médecine en 1923 pour l'isolement de l'insuline (Banting partagea sa récompense avec Charles Best). Avant 1922, un enfant diabétique était condamné à mourir en quelques mois. Le premier patient traité, Leonard Thompson, un adolescent de 14 ans mourant, reprit du poids et survécut grâce aux injections d'insuline. Aujourd'hui, l'insuline injectable transforme le diabète de type 1 en maladie contrôlable — y compris au Cameroun, où sa disponibilité et sa conservation au froid dans les hôpitaux de district restent un enjeu majeur de santé publique.
\end{savaistu}

\courssub{Le pancréas endocrine : les îlots de Langerhans}

Le pancréas est une glande \cle{mixte} : sa partie \emph{exocrine} sécrète le suc pancréatique (enzymes digestives) dans le duodénum ; sa partie \emph{endocrine} est constituée d'amas de cellules dispersés dans la glande, les \cle{îlots de Langerhans} (environ un million chez l'Homme), qui déversent leurs hormones directement dans le sang.

Deux types cellulaires nous intéressent :
\begin{itemize}
\item les \cle{cellules bêta} (environ 65 à 80 \% des cellules des îlots, situées au centre) sécrètent l'\cle{insuline} ;
\item les \cle{cellules alpha} (environ 20 \%, en périphérie des îlots) sécrètent le \cle{glucagon}.
\end{itemize}

\begin{definition}[L'insuline]
L'\cle{insuline} est une hormone protéique \cle{hypoglycémiante} (elle fait baisser la glycémie), sécrétée par les cellules bêta des îlots de Langerhans du pancréas, en réponse à une élévation de la glycémie.
\end{definition}

\begin{definition}[Le glucagon]
Le \cle{glucagon} est une hormone protéique \cle{hyperglycémiante} (elle fait monter la glycémie), sécrétée par les cellules alpha des îlots de Langerhans du pancréas, en réponse à une baisse de la glycémie.
\end{definition}

\begin{popfigure}
\begin{tikzpicture}[scale=0.9, every node/.style={font=\small}]
% Pancréas (forme allongée)
\draw[fill=poporangeL, draw=poporange, thick] (0,0) .. controls (1,1.2) and (4,1.6) .. (7,0.8)
   .. controls (8.2,0.5) and (8.2,-0.6) .. (7,-0.9)
   .. controls (4,-1.6) and (1,-1.2) .. (0,0);
\node[popdark] at (4.2,0.2) {\textbf{Pancréas}};
% Canal excréteur
\draw[popmuted, thick, dashed] (0.3,-0.1) .. controls (3,0.1) and (6,0) .. (7.8,-0.2);
\node[popmuted, font=\footnotesize] at (7.9,-0.55) {canal excréteur};
% Îlots
\foreach \x/\y in {1.5/0.3, 3.2/-0.3, 5.2/0.4, 6.5/-0.2}{
  \draw[fill=popgreenL, draw=popgreen, thick] (\x,\y) circle (0.28);
}
\node[popgreen, font=\footnotesize] at (5.2,0.95) {îlots de Langerhans};
\draw[->, popgreen] (5.2,0.75) -- (5.2,0.68);
% Zoom îlot
\draw[fill=popcream, draw=popdark, thick] (10.2,0) circle (1.9);
\foreach \a in {0,60,120,180,240,300}{
  \draw[fill=popblueL, draw=popblue] ({10.2+1.45*cos(\a)},{1.45*sin(\a)}) circle (0.22);
}
\foreach \a in {30,90,150,210,270,330}{
  \draw[fill=poppinkL, draw=poppink] ({10.2+0.75*cos(\a)},{0.75*sin(\a)}) circle (0.22);
}
\draw[fill=poppinkL, draw=poppink] (10.2,0) circle (0.22);
\node[font=\footnotesize, popblue] at (12.9,1.3) {cellules $\alpha$ (glucagon)};
\draw[->, popblue] (12.6,1.15) -- (11.6,1.05);
\node[font=\footnotesize, poppink] at (12.9,-1.3) {cellules $\beta$ (insuline)};
\draw[->, poppink] (12.6,-1.15) -- (10.9,-0.55);
\draw[->, popdark, thick] (6.8,0.4) -- (8.1,0.4);
\end{tikzpicture}
\legende{Le pancréas, glande mixte. Les îlots de Langerhans (en vert), disséminés dans le tissu exocrine, constituent la partie endocrine. À droite, coupe d'un îlot : les cellules alpha (périphériques) sécrètent le glucagon, les cellules bêta (centrales) sécrètent l'insuline, directement dans les capillaires sanguins.}
\end{popfigure}

\courssub{Mise en évidence du rôle du foie}

\textbf{Observation et hypothèse.} Le foie est traversé par un double réseau sanguin : la \cle{veine porte} hépatique lui apporte le sang venant de l'intestin (riche en nutriments après un repas), et les veines \cle{sus-hépatiques} ramènent le sang qui a traversé le foie vers la circulation générale. Si le foie intervient dans la régulation de la glycémie, la concentration en glucose devrait différer entre le sang entrant et le sang sortant.

\begin{experience}[Dosages comparés du glucose dans le sang hépatique]
On dose le glucose sanguin chez un chien, simultanément dans la veine porte (sang entrant dans le foie) et dans une veine sus-hépatique (sang sortant du foie), dans deux situations :

\begin{center}
\begin{tabularx}{\linewidth}{|l|X|X|X|}
\hline
\rowcolor{popblueL} \textbf{Situation} & \textbf{Veine porte} & \textbf{Veine sus-hépatique} & \textbf{Interprétation} \\
\hline
Après un repas glucidique & \SI{1,8}{g/L} & \SI{1,2}{g/L} & Le foie \emph{prélève} du glucose \\
\hline
À jeun prolongé & \SI{0,9}{g/L} & \SI{1,1}{g/L} & Le foie \emph{libère} du glucose \\
\hline
\end{tabularx}
\end{center}

\textbf{Conclusion :} le foie est un \cle{régulateur-tampon} : il stocke le glucose quand il est abondant et le restitue quand il manque, maintenant ainsi la glycémie constante.
\end{experience}

\textbf{Mécanismes hépatiques.} Le glucose n'est pas stocké tel quel : il est polymérisé en \cle{glycogène}, une macromolécule de réserve (le foie d'un adulte peut en stocker environ \SI{100}{g}).

\begin{itemize}
\item \textbf{Glycogénogenèse} (après un repas, glycémie élevée) : sous l'action de l'\cle{insuline}, les hépatocytes transforment le glucose en glycogène :
\[ n\;\text{glucose} \xrightarrow{\text{insuline}} \text{glycogène} + n\;\ce{H2O} \]
\item \textbf{Glycogénolyse} (à jeun, glycémie basse) : sous l'action du \cle{glucagon} et de l'\cle{adrénaline}, le glycogène hépatique est hydrolysé et le glucose libéré passe dans le sang :
\[ \text{glycogène} + n\;\ce{H2O} \xrightarrow{\text{glucagon, adrénaline}} n\;\text{glucose} \]
\end{itemize}

Le foie peut aussi fabriquer du glucose neuf à partir d'acides aminés ou de lactate : c'est la \cle{néoglucogenèse}, stimulée par le glucagon lors des jeûnes prolongés.

\begin{propriete}[Savoir admis : bilan des hormones glucorégulatrices]
L'\cle{insuline} est la \emph{seule} hormone hypoglycémiante de l'organisme. Elle agit en favorisant l'entrée du glucose dans les cellules (muscles, tissu adipeux) et sa transformation en glycogène hépatique et musculaire. Le \cle{glucagon} et l'\cle{adrénaline} (sécrétée par les glandes surrénales) sont hyperglycémiants : ils stimulent la glycogénolyse hépatique. Cette propriété est admise ; aucune démonstration n'est exigible au Baccalauréat.
\end{propriete}

\begin{attention}[Erreur fréquente]
N'attribue \emph{jamais} au pancréas la production ou le stockage du glycogène ! Le pancréas \emph{sécrète les hormones} (insuline, glucagon) ; c'est le \textbf{foie} (et, dans une moindre mesure, les \textbf{muscles}) qui \emph{stockent} le glucose sous forme de glycogène. De même, le glycogène musculaire ne peut pas être restitué au sang : il sert exclusivement à l'effort musculaire. Seul le foie libère du glucose dans la circulation.
\end{attention}

\courssub{La régulation nerveuse de la glycémie}

La régulation hormonale est doublée d'un contrôle nerveux, plus rapide. Des \cle{glucorécepteurs} (cellules sensibles à la concentration en glucose du sang, situées notamment dans l'hypothalamus et au niveau hépatique) informent en permanence un \cle{centre glucorégulateur} situé dans l'\cle{hypothalamus}.

Selon la valeur de la glycémie, ce centre envoie des messages via les \cle{nerfs végétatifs} (autonomes) :
\begin{itemize}
\item en cas d'hyperglycémie : stimulation des cellules bêta du pancréas (sécrétion d'insuline) et inhibition des cellules alpha ;
\item en cas d'hypoglycémie : stimulation des cellules alpha (glucagon) et des \cle{glandes surrénales}, qui libèrent de l'\cle{adrénaline}, hormone hyperglycémiante qui accélère la glycogénolyse hépatique.
\end{itemize}

\begin{exemplebox}[L'adrénaline, hormone du stress]
Un élève stressé juste avant l'épreuve de SVT, ou un cultivateur qui fuit un serpent dans sa plantation de cacao, sécrète de l'adrénaline : sa glycémie monte rapidement pour fournir du glucose aux muscles et au cerveau. C'est la réponse d'urgence, médiée par le système nerveux sympathique.
\end{exemplebox}

\courssub{Schéma fonctionnel de la régulation de la glycémie}

\begin{methode}[Construire un schéma fonctionnel de régulation]
Pour schématiser toute régulation physiologique, procède en cinq étapes :
\begin{enumerate}
\item \textbf{Identifier la variable régulée} : ici, la glycémie (valeur de référence $\approx$ \SI{1}{g/L}).
\item \textbf{Identifier les récepteurs} : les glucorécepteurs, qui détectent les écarts.
\item \textbf{Identifier le centre de commande} : le centre glucorégulateur hypothalamique.
\item \textbf{Identifier les voies de commande} : nerfs végétatifs (voie nerveuse) et hormones transportées par le sang (voie hormonale).
\item \textbf{Identifier les effecteurs et le rétrocontrôle} : pancréas (insuline, glucagon), foie (glycogénogenèse/glycogénolyse), surrénales (adrénaline) ; l'effet obtenu ramène la variable vers sa valeur de référence et \emph{freine} la réponse : c'est le \cle{rétrocontrôle négatif}.
\end{enumerate}
\end{methode}

\begin{popfigure}
\begin{tikzpicture}[font=\small,
  box/.style={draw=popdark, thick, rounded corners, fill=popcream, align=center, minimum height=0.9cm, inner sep=4pt},
  eff/.style={draw=popblue, thick, rounded corners, fill=popblueL, align=center, minimum height=0.9cm, inner sep=4pt},
  fleche/.style={->, thick, popdark}]
% Variable régulée
\node[box, fill=popgoldL, draw=popgold, align=center] (gly) at (0,0){\textbf{Glycémie}\\ $\approx$ \SI{1}{g/L}};
% Récepteurs
\node[box, align=center] (rec) at (0,-2.2){Glucorécepteurs\\ (hypothalamus, foie)};
% Centre
\node[box, fill=poppurpleL, draw=poppurple, align=center] (centre) at (0,-4.4){\textbf{Centre glucorégulateur}\\ (hypothalamus)};
% Effecteurs
\node[eff, align=center] (panc) at (-4.6,-6.8){\textbf{Pancréas}\\ insuline ($\beta$)\\ glucagon ($\alpha$)};
\node[eff, align=center] (foie) at (0,-6.8){\textbf{Foie}\\ glycogénogenèse\\ glycogénolyse};
\node[eff, align=center] (surr) at (4.6,-6.8){\textbf{Surrénales}\\ adrénaline};
% Flèches
\draw[fleche] (gly) -- (rec) node[midway, right, font=\footnotesize]{informe};
\draw[fleche] (rec) -- (centre) node[midway, right, font=\footnotesize]{influx nerveux};
\draw[fleche, poppurple] (centre) -- (panc) node[midway, above, sloped, font=\footnotesize]{nerfs végétatifs};
\draw[fleche, poppurple] (centre) -- (foie);
\draw[fleche, poppurple] (centre) -- (surr);
% Retour vers glycémie
\draw[fleche, popgreen] (panc.west) .. controls (-7.5,-4) and (-7.5,-1) .. (gly.west)
   node[midway, left, font=\footnotesize, align=center]{hormones\\ via le sang};
\draw[fleche, popgreen] (foie.north) .. controls (0.6,-5.5) and (1.4,-3) .. (gly.east)
   node[midway, right, font=\footnotesize, align=center]{glucose stocké\\ ou libéré};
\draw[fleche, popgreen] (surr.east) .. controls (7.5,-4) and (7.5,-1) .. (gly.east);
% Rétrocontrôle
\node[font=\footnotesize, popgreen, align=center] at (0,-8.3) {Rétrocontrôle négatif : le retour de la glycémie à la normale stoppe la réponse};
\end{tikzpicture}
\legende{Schéma fonctionnel de la régulation de la glycémie. La variable régulée (glycémie) est surveillée par des glucorécepteurs qui informent le centre hypothalamique ; celui-ci commande, par voie nerveuse et hormonale, trois effecteurs (pancréas, foie, surrénales) dont les actions antagonistes ramènent la glycémie à sa valeur de référence.}
\end{popfigure}

\begin{aretenir}[L'essentiel de la section]
\begin{itemize}
\item Le \cle{pancréas endocrine} (îlots de Langerhans) sécrète deux hormones antagonistes : l'\cle{insuline} (cellules bêta, hypoglycémiante, seule hormone qui fait baisser la glycémie) et le \cle{glucagon} (cellules alpha, hyperglycémiant).
\item Le \cle{foie} stocke le glucose en glycogène (glycogénogenèse, sous l'effet de l'insuline) et le libère (glycogénolyse, sous l'effet du glucagon et de l'adrénaline).
\item Le \cle{centre glucorégulateur hypothalamique}, informé par des glucorécepteurs, coordonne le tout via les nerfs végétatifs.
\item L'ensemble fonctionne en \cle{boucle de rétrocontrôle négatif} : toute déviation de la glycémie déclenche une réponse qui la corrige.
\end{itemize}
\end{aretenir}

\begin{exoresolu}[Interpréter des dosages sanguins]
\textbf{Énoncé.} Chez deux chiens A et B, on dose le glucose (en \SI{}{g/L}) dans la veine porte (VP) et la veine sus-hépatique (VSH), une heure après un repas riche en amidon :

\begin{center}
\begin{tabular}{|l|c|c|}
\hline
\rowcolor{popblueL} & VP & VSH \\
\hline
Chien A (sain) & 1,7 & 1,1 \\
\hline
Chien B (pancréatectomisé) & 2,4 & 2,3 \\
\hline
\end{tabular}
\end{center}

Interprète ces résultats et explique la différence observée entre les deux chiens.
\tcblower
\textbf{Solution détaillée.}

\textbf{1. Chien A.} La glycémie de la veine porte (\SI{1,7}{g/L}) est supérieure à celle de la veine sus-hépatique (\SI{1,1}{g/L}) : le sang qui traverse le foie y laisse du glucose. Le foie du chien A prélève donc du glucose sanguin après le repas. Sous l'action de l'insuline, sécrétée par les cellules bêta des îlots de Langerhans en réponse à l'hyperglycémie post-prandiale, les hépatocytes réalisent la glycogénogenèse : le glucose est stocké sous forme de glycogène. La glycémie du sang quittant le foie est ainsi ramenée vers la normale.

\textbf{2. Chien B.} Les glycémies de la veine porte (\SI{2,4}{g/L}) et de la veine sus-hépatique (\SI{2,3}{g/L}) sont presque identiques : le foie ne prélève pratiquement pas de glucose. Privé de pancréas, le chien B ne sécrète plus d'insuline ; or l'insuline est indispensable à la glycogénogenèse hépatique et à l'entrée du glucose dans les cellules. Le glucose absorbé au niveau intestinal s'accumule donc dans le sang : c'est l'hyperglycémie permanente caractéristique du diabète expérimental.

\textbf{3. Conclusion.} La comparaison confirme le double rôle du couple pancréas-foie : le pancréas, par l'insuline, commande ; le foie, par la glycogénogenèse, exécute le stockage du glucose.
\end{exoresolu}

\begin{exercice}[À toi de jouer]
Un sujet à jeun depuis 18 heures présente une glycémie de \SI{0,95}{g/L} dans la veine porte et de \SI{1,15}{g/L} dans la veine sus-hépatique.
\begin{enumerate}
\item Compare ces deux valeurs et déduis-en le sens des échanges de glucose au niveau du foie.
\item Nomme la réaction biochimique hépatique en cause et les deux hormones qui la stimulent.
\item Explique pourquoi cette libération de glucose est indispensable au cerveau.
\end{enumerate}
\end{exercice}

\begin{corrige}[Exercice « À toi de jouer »]
\begin{enumerate}
\item La glycémie est plus élevée dans la veine sus-hépatique (\SI{1,15}{g/L}) que dans la veine porte (\SI{0,95}{g/L}) : le sang s'enrichit en glucose en traversant le foie. À jeun, le foie \emph{libère} donc du glucose dans la circulation.
\item Il s'agit de la \cle{glycogénolyse} (hydrolyse du glycogène hépatique en glucose), stimulée par le \cle{glucagon} (cellules alpha du pancréas) et l'\cle{adrénaline} (glandes surrénales). Si le jeûne se prolonge, la néoglucogenèse prend le relais.
\item Le glucose est le carburant quasi exclusif des neurones, qui ne possèdent presque pas de réserves. Sans libération hépatique de glucose, la glycémie chuterait (hypoglycémie), entraînant troubles de la conscience, convulsions, voire coma. Le foie garantit ainsi l'approvisionnement permanent du cerveau entre les repas.
\end{enumerate}
\end{corrige}

\coursec{Le diabète : typologie, causes et moyens de lutte}

Dans la section précédente, nous avons montré que la glycémie d'un sujet sain est maintenue autour de \SI{1,00}{g/L} grâce à une régulation fine où le pancréas (insuline hypoglycémiante, glucagon hyperglycémiant) et le foie (stockage et libération du glucose) jouent des rôles complémentaires. Mais que se passe-t-il lorsque cette régulation déraille de façon permanente ? On observe alors une \cle{hyperglycémie chronique} : c'est le \cle{diabète}, l'une des maladies métaboliques les plus répandues au monde, y compris dans nos villes camerounaises. Cette section a pour but de définir le diabète, d'en distinguer les deux grands types, d'en expliquer les causes et les complications, et enfin de dégager les moyens de lutte et de dépistage.

\courssub{Définition et symptômes du diabète}

\begin{definition}[Le diabète sucré]
Le \cle{diabète} (ou diabète sucré) est une maladie chronique caractérisée par une \textbf{hyperglycémie permanente}, c'est-à-dire une \cle{glycémie à jeun} supérieure ou égale à \SI{1,26}{g/L} (soit \SI{7,0}{mmol/L}), constatée à deux reprises. Elle traduit un défaut de la régulation hormonale de la glycémie : soit l'insuline est absente (déficit de sécrétion), soit elle est inefficace (insulinorésistance).
\end{definition}

\textbf{Observation clinique introductive.} Dans un centre de santé de Douala, un patient de 45 ans se plaint de boire énormément d'eau, d'uriner très souvent (y compris la nuit), d'avoir faim en permanence et pourtant de maigrir. Sa glycémie à jeun est de \SI{2,10}{g/L}. Ces signes sont typiques du diabète et s'expliquent directement par l'hyperglycémie :

\begin{itemize}
\item \textbf{Polyurie} : le glucose excédentaire dépasse le seuil de réabsorption rénale (environ \SI{1,80}{g/L}) ; il entraîne de l'eau dans les urines par osmose, dont le volume augmente considérablement.
\item \textbf{Glycosurie} : présence de glucose dans les urines (normalement nulle), conséquence directe du dépassement du seuil rénal de réabsorption du glucose.
\item \textbf{Polydipsie} : soif intense et permanente, liée à la déshydratation cellulaire et extracellulaire provoquée par la polyurie osmotique.
\item \textbf{Polyphagie} : faim excessive, car le glucose sanguin n'entre pas correctement dans les cellules (manque d'insuline ou résistance à l'insuline) qui restent « affamées » malgré l'abondance de sucre dans le sang.
\item \textbf{Amaigrissement} : les cellules, privées de glucose utilisable, puisent dans les réserves de lipides (lipolyse) et de protéines (protéolyse) pour couvrir leurs besoins énergétiques.
\end{itemize}

\begin{aretenir}[Les signes cardinaux du diabète]
Le diabète se reconnaît classiquement par le quatuor : \textbf{polyurie -- polydipsie -- polyphagie -- amaigrissement}, auquel s'ajoute la \textbf{glycosurie}. Le diagnostic biologique repose sur la glycémie à jeun $\geq \SI{1,26}{g/L}$ (confirmée deux fois) ou une glycémie $\geq \SI{2,00}{g/L}$ deux heures après une HGPO.
\end{aretenir}

\courssub{Les deux grands types de diabète}

Tous les diabètes ne se ressemblent pas : on distingue principalement le diabète de type 1 et le diabète de type 2, qui diffèrent par l'âge d'apparition, la cause physiopathologique et le traitement.

\textbf{Diabète de type 1 (diabète insulinodépendant ou diabète juvénile).}

Il apparaît le plus souvent chez l'\textbf{enfant ou le jeune adulte} (pic de fréquence vers 10--14 ans). Sa cause est une \textbf{destruction auto-immune des cellules $\beta$} des îlots de Langerhans du pancréas : le système immunitaire du patient fabrique des auto-anticorps qui détruisent ses propres cellules productrices d'insuline. Conséquence : le pancréas ne sécrète \textbf{plus du tout d'insuline} (\cle{insulinémie} quasi nulle à jeun et après stimulation). Le glucose ne peut plus entrer dans les cellules insulinodépendantes (muscle, tissu adipeux) ni être stocké par le foie sous forme de glycogène : la glycémie s'élève fortement. Le seul traitement vital est l'apport quotidien d'insuline par \textbf{injections sous-cutanées} (stylo injecteur ou pompe) ; l'insuline, étant une protéine, serait digérée si elle était avalée. On parle donc de diabète \cle{insulinodépendant} (ou \cle{insulino-requérant}).

\textbf{Diabète de type 2 (diabète non insulinodépendant ou diabète de l'adulte).}

Il apparaît généralement après 40 ans, chez des sujets souvent en \textbf{surpoids ou obèses} (obésité androïde), sédentaires, et fréquemment dans des familles où le diabète est fréquent (composante \textbf{héréditaire} polygénique). Ici, le pancréas sécrète de l'insuline de façon \textbf{normale, voire excessive} (hyperinsulinémie compensatrice initiale), mais les cellules cibles (muscles, foie, tissu adipeux) répondent mal à son message : c'est l'\cle{insulinorésistance}. Les récepteurs membranaires à insuline sont en nombre insuffisant ou leur transduction du signal est défaillante. L'insuline est donc \textbf{présente mais inefficace}, et la glycémie reste élevée. À long terme, le pancréas s'épuise (insulinopénie secondaire). Le traitement repose d'abord sur les \textbf{mesures hygiéno-diététiques} (régime pauvre en sucres rapides et en graisses saturées, activité physique, perte de poids), puis sur des \textbf{antidiabétiques oraux} (biguanides, sulfamides hypoglycémiants, inhibiteurs de la DPP-4, etc.) qui améliorent la sensibilité à l'insuline ou stimulent sa sécrétion. L'insulinothérapie peut devenir nécessaire en cas d'échec des traitements oraux (insulinopénie installée).

\begin{attention}[Erreur fréquente à éviter]
Ne dis jamais que le diabète de type 2 est dû à un \emph{manque d'insuline} initial ! Dans le diabète de type 2, l'insuline est \textbf{présente en quantité normale ou élevée} (hyperinsulinémie), mais les cellules cibles y sont \textbf{résistantes}. Seul le diabète de type 1 résulte d'une \textbf{absence de sécrétion d'insuline} (insulinopénie) dès le diagnostic. Cette distinction physiopathologique est systématiquement attendue au Baccalauréat.
\end{attention}

\courssub{Mise en évidence expérimentale de la typologie}

\begin{experience}[Dosage de l'insulinémie et HGPO chez trois sujets]
\textbf{Protocole.} On prélève le sang à jeun chez trois sujets A, B et C, puis on leur fait ingérer \SI{75}{g} de glucose dissous dans \SI{250}{mL} d'eau (Hyperglycémie Provoquée par le Glucose par voie Orale, \cle{HGPO}). On mesure la glycémie toutes les 30 minutes pendant 2 heures, et on dose l'insuline sanguine (\cle{insulinémie}) à jeun.

\textbf{Résultats.}
\begin{itemize}
\item Sujet A (témoin sain) : glycémie à jeun \SI{0,90}{g/L} ; pic à \SI{1,40}{g/L} vers 30 min ; retour à la normale ($< \SI{1,10}{g/L}$) en 2 h ; insulinémie à jeun normale (\SI{10}{\micro UI/mL}), augmentation post-prandiale normale.
\item Sujet B : glycémie à jeun \SI{2,00}{g/L} ; pic très élevé (\SI{3,50}{g/L}) et persistant (\SI{3,20}{g/L} à 2 h) ; insulinémie \textbf{quasi nulle} ($< \SI{2}{\micro UI/mL}$) à jeun et après charge.
\item Sujet C : glycémie à jeun \SI{1,50}{g/L} ; pic élevé (\SI{2,60}{g/L}) et retour lent (\SI{2,20}{g/L} à 2 h) ; insulinémie \textbf{normale voire élevée} (\SI{25}{\micro UI/mL}) à jeun, avec pic insulinémique retardé.
\end{itemize}

\textbf{Interprétation.} Le sujet A est sain : son insuline ramène la glycémie à la normale. Le sujet B est diabétique de \textbf{type 1} : l'absence d'insuline explique l'hyperglycémie massive et durable malgré la charge de glucose. Le sujet C est diabétique de \textbf{type 2} : son insuline est présente (voire élevée) mais inefficace (insulinorésistance), d'où une hyperglycémie modérée mais persistante.

\textbf{Conclusion.} Le dosage de l'insulinémie (à jeun et/ou post-charge) permet de \textbf{typer} un diabète : insuline absente ou très faible $\Rightarrow$ type 1 ; insuline normale ou élevée avec hyperglycémie $\Rightarrow$ type 2.
\end{experience}

\begin{popfigure}
\begin{center}
\begin{tikzpicture}
\begin{axis}[
width=0.95\linewidth, height=7cm,
xlabel={Temps (min)}, ylabel={Glyc\'emie (g/L)},
xmin=0, xmax=130, ymin=0, ymax=4.0,
xtick={0,30,60,90,120},
ytick={0,1,2,3,4},
legend style={at={(0.98,0.97)}, anchor=north east, font=\small, cells={anchor=west}},
grid=major, grid style={popline!40},
]
% Sujet Sain : pic 1.4 a 30min, retour 1.0 a 120min
\addplot[popcurve, popblue, domain=0:120, samples=200] {0.9 + 0.6*exp(-(x-30)*(x-30)/800) - 0.1*exp(-(x-60)*(x-60)/2000)};
\addlegendentry{Sujet sain (A)}
% Sujet Type 1 : depart 2.0, monte a 3.5, reste haut
\addplot[popcurve, poppink, domain=0:120, samples=200] {2.0 + 1.5*(1-exp(-x/15))*exp(-x/300)};
\addlegendentry{Diab\'etique type 1 (B)}
% Sujet Type 2 : depart 1.5, monte a 2.6, descend tres lentement
\addplot[popcurve, poporange, domain=0:120, samples=200] {1.5 + 1.2*(1-exp(-x/20))*exp(-x/500)};
\addlegendentry{Diab\'etique type 2 (C)}
% Seuil diabete
\addplot[dashed, popdark, thick, domain=0:120] {1.26};
\addlegendentry{Seuil diag. diab\`ete (\SI{1,26}{g/L})}
% Seuil renal
\addplot[dotted, popmuted, thick, domain=0:120] {1.80};
\addlegendentry{Seuil r\'enal (\SI{1,80}{g/L})}
\end{axis}
\end{tikzpicture}
\end{center}
\legende{Courbes d'hyperglycémie provoquée par le glucose par voie orale (HGPO, ingestion de \SI{75}{g} de glucose à $t = 0$) chez trois sujets. Chez le sujet sain, la glycémie revient sous \SI{1,10}{g/L} en 2 h ; chez le diabétique de type 1, l'hyperglycémie est massive et durable (insuline absente) ; chez le diabétique de type 2, elle est intermédiaire et persistante (insuline présente mais inefficace). Le seuil rénal de réabsorption (\SI{1,80}{g/L}) est dépassé chez les deux diabétiques, expliquant la glycosurie.}
\end{popfigure}

\begin{popfigure}
\begin{center}
\begin{tabularx}{\linewidth}{|l|X|X|}\hline
\rowcolor{popblueL} \textbf{Critère} & \textbf{Diabète de type 1} & \textbf{Diabète de type 2} \\ \hline
\textbf{Âge d'apparition} & Enfant, adolescent, jeune adulte (souvent $< 30$ ans) & Adulte (généralement $> 40$ ans), de plus en plus jeune (adolescents obèses) \\ \hline
\textbf{Corpulence habituelle} & Mince, amaigrissement récent & Surpoids ou obésité (souvent androïde) \\ \hline
\textbf{Cause physiopathologique} & Destruction auto-immune des cellules $\beta$ pancréatiques $\rightarrow$ \textbf{insulinopénie} absolue & \textbf{Insulinorésistance} périphérique (récepteurs) + épuisement secondaire des cellules $\beta$ \\ \hline
\textbf{Insulinémie (à jeun / post-HGPO)} & Nulle ou très faible ($< \SI{5}{\micro UI/mL}$) & Normale, élevée (hyperinsulinémie) ou modérément diminuée (stade évolué) \\ \hline
\textbf{Facteurs favorisants} & Prédisposition génétique (HLA) + facteur environnemental (virus ?) & Obésité, sédentarité, hérédité forte, âge, alimentation déséquilibrée \\ \hline
\textbf{Traitement de référence} & \textbf{Insulinothérapie} vitale (injections multiples ou pompe) + éducation thérapeutique & \textbf{Mesures hygiéno-diététiques} (régime, sport, perte de poids) $\pm$ \textbf{Antidiabétiques oraux} $\pm$ Insuline (échec oral) \\ \hline
\textbf{Urgences métaboliques} & Acidocétose diabétique (fréquente au dévoilement) & État hyperosmolaire (sujet âgé, déshydraté), acidocétose rare \\ \hline
\end{tabularx}
\end{center}
\legende{Tableau comparatif des deux grands types de diabète sucré.}
\end{popfigure}

\courssub{Complications du diabète non contrôlé}

Un diabète mal équilibré (hyperglycémie chronique) expose, à long terme, à des \textbf{lésions des vaisseaux sanguins (angiopathies) et des nerfs (neuropathies)}, dues à la toxicité chronique du glucose (glycation non enzymatique des protéines, stress oxydatif, activation de la voie des polyols) :

\begin{itemize}
\item \textbf{Angiopathies microvasculaires (spécifiques du diabète) :}
      \begin{itemize}
      \item \textbf{Rétinopathie diabétique} : atteinte des capillaires rétiniens (micro-anévrismes, hémorragies, néovaisseaux), première cause de cécité avant 60 ans.
      \item \textbf{Néphropathie diabétique} : glomérulosclérose nodulaire (Kimmelstiel-Wilson), proteinurie puis insuffisance rénale terminale (dialyse, greffe).
      \item \textbf{Neuropathie diabétique} : atteinte des nerfs périphériques (polyneuropathie sensitive symétrique, douleurs, perte de sensibilité aux pieds) et autonomes (troubles érectiles, gastroparésie, hypotension orthostatique, anhidrose).
      \end{itemize}
\item \textbf{Angiopathies macrovasculaires (athérosclérose accélérée) :}
      \begin{itemize}
      \item Coronaires : angor, infarctus du myocarde (silencieux possible par neuropathie).
      \item Cérébraux : accidents vasculaires cérébraux (AVC ischémiques).
      \item Artères des membres inférieurs : artérite des membres inférieurs (risque d'ischémie critique, gangrène).
      \end{itemize}
\item \textbf{Pied diabétique} : association de neuropathie (plaies non ressenties, déformations charcot) et d'artériopathie (mauvaise cicatrisation) $\pm$ infection. Première cause d'amputation non traumatique.
\end{itemize}

\courssub{Moyens de lutte contre le diabète}

La lutte repose sur la prévention primaire (éviter l'apparition), le dépistage précoce (diagnostiquer avant complications) et la prise en charge globale (équilibrer pour prévenir les complications).

\begin{itemize}
\item \textbf{Prévention primaire (surtout Type 2) :} promotion d'une \textbf{alimentation équilibrée} (réduction des sucres ajoutés, boissons sucrées, graisses saturées ; augmentation fibres, légumes, céréales complètes) et d'une \textbf{activité physique régulière} (au moins 150 min/semaine d'activité modérée, ex. marche rapide 30 min/jour). Lutte contre la sédentarité et l'obésité dès l'enfance.
\item \textbf{Dépistage systématique :} mesure de la \textbf{glycémie à jeun} (ou HbA1c) chez tout adulte $> 40$ ans, ou plus jeune en cas de facteurs de risque (obésité, antécédents familiaux, grossesse antérieure macrosomique, HTA, dyslipidémie). Au Cameroun, les campagnes de la Journée Mondiale du Diabète (14 novembre) offrent des dépistages gratuits.
\item \textbf{Diagnostic et Typage :} Glycémie à jeun $\geq \SI{1,26}{g/L}$ (2 fois) \textbf{OU} Glycémie 2h post-HGPO $\geq \SI{2,00}{g/L}$ \textbf{OU} HbA1c $\geq 6,5\%$ \textbf{OU} Glycémie quelconque $\geq \SI{2,00}{g/L}$ + signes cliniques. Typage par insulinémie / peptide C / auto-anticorps anti-îlots / anti-GAD / anti-IA2.
\item \textbf{Prise en charge (Éducation Thérapeutique du Patient - ETP) :}
      \begin{itemize}
      \item \textbf{Type 1 :} Insulinothérapie basale-bolus (ou pompe) + comptage de glucides + autosurveillance glycémique (lecteur ou capteur glucose continu) + régime libre adapté.
      \item \textbf{Type 2 :} Règles hygiéno-diététiques (premier traitement) + Antidiabétiques oraux (Metformine en 1ère ligne) $\pm$ Agonistes GLP-1 (injectables, perte de poids) $\pm$ Insuline (échec). Autosurveillance glycémique adaptée.
      \item \textbf{Objectifs d'équilibre (cibles) :} HbA1c $< 7\%$ (individualisé), Glycémie à jeun $0,70$--$1,10$ g/L, Glycémie post-prandiale $< \SI{1,60}{g/L}$, PA $< 130/80$ mmHg, LDL-c $< 1,00$ g/L (très haut risque).
      \item \textbf{Surveillance des complications :} Fond d'œil annuel, microalbuminurie / clairance rénale annuelle, examen des pieds à chaque consultation, bilan lipidique, ECG.
      \end{itemize}
\end{itemize}

\begin{methode}[Exploiter un bilan médical pour diagnostiquer et typer un diabète]
\begin{enumerate}
\item \textbf{Comparer la glycémie à jeun} à la valeur de référence : si elle est $\geq \SI{1,26}{g/L}$ (à deux reprises), il y a diabète. Si elle est entre \SI{1,10}{g/L} et \SI{1,25}{g/L} $\rightarrow$ hyperglycémie à jeun modifiée (prédiabète).
\item \textbf{Examiner la courbe d'HGPO} (si réalisée) : une glycémie $\geq \SI{2,00}{g/L}$ deux heures après l'ingestion de \SI{75}{g} de glucose confirme le diabète. Entre \SI{1,40}{g/L} et \SI{1,99}{g/L} $\rightarrow$ intolérance au glucose (prédiabète).
\item \textbf{Rechercher une glycosurie} (bandelette urinaire) : sa présence indique que le seuil rénal de réabsorption du glucose (\SI{1,80}{g/L}) est dépassé ; elle soutient le diagnostic mais ne suffit pas seule.
\item \textbf{Doser l'insulinémie} (et/ou le peptide C) : insuline absente ou très faible $\Rightarrow$ diabète de type 1 ; insuline normale ou élevée $\Rightarrow$ diabète de type 2. La recherche d'auto-anticorps (anti-GAD, anti-IA2, anti-îlots) confirme l'origine auto-immune du type 1.
\item \textbf{Conclure} en croisant tous les paramètres (âge, corpulence, antécédents familiaux, cétose, auto-anticorps) et en proposant le traitement adapté (insulinothérapie vitale pour type 1 ; hygiène de vie $\pm$ antidiabétiques oraux pour type 2).
\end{enumerate}
\end{methode}

\begin{exemplebox}[Un exemple chiffré]
Un jeune patient de 16 ans, amaigri de 5 kg en 1 mois, polyuro-polydipsique, présente une glycémie à jeun de \SI{2,00}{g/L}, une glycosurie +++, une céturie + et une insulinémie quasi nulle (\SI{1,5}{\micro UI/mL}). Conclusion : diabète de \textbf{type 1} (insulinodépendant) au dévoilement ; le traitement sera l'injection quotidienne d'insuline (schéma basale-bolus) associée à une éducation diététique et à l'autosurveillance glycémique. L'HbA1c initiale sera probablement $> 10\%$.
\end{exemplebox}

\begin{exoresolu}[Diagnostic et typage d'un diabète à partir d'analyses médicales]
\textbf{Énoncé.} Les résultats d'analyses de deux patients sont consignés ci-dessous.

\begin{center}
\begin{tabularx}{\linewidth}{|l|X|X|}\hline
\rowcolor{popblueL} \textbf{Paramètre} & \textbf{Patient M.} (17 ans, mince) & \textbf{Patient N.} (52 ans, obèse, IMC 32) \\ \hline
Glycémie à jeun & \SI{2,10}{g/L} & \SI{1,60}{g/L} \\ \hline
Glycémie 2 h après HGPO (\SI{75}{g}) & \SI{3,40}{g/L} & \SI{2,40}{g/L} \\ \hline
Glycosurie & présente (+++) & présente (+) \\ \hline
Insulinémie à jeun & \SI{1,2}{\micro UI/mL} (quasi nulle) & \SI{28}{\micro UI/mL} (élevée) \\ \hline
Peptide C & très bas & normal / élevé \\ \hline
Auto-anticorps anti-GAD & positifs & négatifs \\ \hline
\end{tabularx}
\end{center}

\begin{enumerate}
\item Ces deux patients sont-ils diabétiques ? Justifier votre réponse par les valeurs seuils.
\item Déterminer le type de diabète de chaque patient en vous basant sur l'insulinémie et les auto-anticorps. Proposer le traitement de première intention pour chacun.
\end{enumerate}
\tcblower
\textbf{Solution détaillée.}
\begin{enumerate}
\item \textbf{Oui, les deux patients sont diabétiques.}
      \begin{itemize}
      \item \textbf{Patient M :} Glycémie à jeun \SI{2,10}{g/L} $> \SI{1,26}{g/L}$ \textbf{ET} Glycémie 2h post-HGPO \SI{3,40}{g/L} $> \SI{2,00}{g/L}$. La glycosurie confirme le dépassement du seuil rénal (\SI{1,80}{g/L}).
      \item \textbf{Patient N :} Glycémie à jeun \SI{1,60}{g/L} $> \SI{1,26}{g/L}$ \textbf{ET} Glycémie 2h post-HGPO \SI{2,40}{g/L} $> \SI{2,00}{g/L}$. La glycosurie est présente.
      \end{itemize}
      Les deux critères diagnostiques (glycémie à jeun et/ou HGPO) sont validés pour chacun.
\item \textbf{Patient M (17 ans, mince) :} Insulinémie quasi nulle (\SI{1,2}{\micro UI/mL}), Peptide C très bas, Auto-anticorps anti-GAD positifs. $\Rightarrow$ \textbf{Diabète de type 1 auto-immun (insulinodépendant).} Destruction des cellules $\beta$. Traitement de première intention : \textbf{Insulinothérapie basale-bolus (injections multiples quotidiennes) ou pompe à insuline}, éducation thérapeutique (comptage glucides, autosurveillance), régime équilibré.
      \textbf{Patient N (52 ans, obèse) :} Insulinémie élevée (\SI{28}{\micro UI/mL} : hyperinsulinémie), Peptide C normal/élevé, Auto-anticorps négatifs. $\Rightarrow$ \textbf{Diabète de type 2 (insulinorésistance).} L'obésité et l'hyperinsulinémie confirment la résistance à l'insuline. Traitement de première intention : \textbf{Mesures hygiéno-diététiques rigoureuses (régime hypocalorique si surpoids, activité physique 30 min/jour) + Metformine (biguanide) per os}, surveillance glycémique, prise en charge des facteurs de risque cardiovasculaires (HTA, dyslipidémie).
\end{enumerate}
\end{exoresolu}

\begin{savaistu}[Le diabète au Cameroun : un défi de santé publique]
Au Cameroun, la prévalence du diabète en milieu urbain dépasse \textbf{6 \% de la population adulte} (études récentes), et elle augmente régulièrement avec la transition nutritionnelle (consommation croissante de sodas, jus industriels, plats riches en sucres rapides et graisses saturées, type \emph{fast-food}), la \textbf{sédentarité} (transports motorisés, écrans) et l'urbanisation. Une part importante (\textbf{50 à 60 \%}) des diabétiques ignore sa maladie (diabète « silencieux ») : le dépistage gratuit organisé lors des campagnes de santé (Journée Mondiale du Diabète le 14 novembre, Journée Mondiale de la Santé le 7 avril) est donc un enjeu majeur. Un diabète dépisté tôt et bien équilibré (HbA1c $< 7\%$) permet d'éviter ou de retarder les complications graves (cécité, amputations, insuffisance rénale, AVC) qui coûtent très cher aux familles et au système de santé. L'éducation thérapeutique du patient (ETP) est la clé de la réussite : « Mieux connaître sa maladie pour mieux la vivre ».
\end{savaistu}

\begin{aretenir}[L'essentiel de la section : Le diabète]
\begin{itemize}
\item \textbf{Définition :} Hyperglycémie chronique : Glycémie à jeun $\geq \SI{1,26}{g/L}$ (2 fois) ou 2h post-HGPO $\geq \SI{2,00}{g/L}$ ou HbA1c $\geq 6,5\%$.
\item \textbf{Signes cardinaux :} Polyurie, Polydipsie, Polyphagie, Amaigrissement, Glycosurie (seuil rénal \SI{1,80}{g/L} dépassé).
\item \textbf{Type 1 (Insulinodépendant) :} Jeune, mince. Cause : Auto-immunité anti-cellules $\beta$ $\rightarrow$ \textbf{Insulinopénie absolue}. Insulinémie nulle. Traitement : \textbf{Insulinothérapie vitale} (injections/pompe). Urgence : Acidocétose.
\item \textbf{Type 2 (Non insulinodépendant initial) :} Adulte ($>40$ ans), obèse/surpoids, sédentaire, hérédité. Cause : \textbf{Insulinorésistance} $\rightarrow$ Hyperinsulinémie compensatrice puis insulinopénie relative. Insulinémie normale/élevée. Traitement : \textbf{Hygiène de vie (régime, sport) + Antidiabétiques oraux (Metformine 1ère ligne)} $\pm$ Insuline ultérieure. Urgence : État hyperosmolaire.
\item \textbf{Typage :} Dosage insulinémie / peptide C / auto-anticorps (anti-GAD, IA2, îlots).
\item \textbf{Complications :} Microvasculaires (Rétinopathie, Néphropathie, Neuropathie) + Macrovasculaires (IDM, AVC, Artérite) + Pied diabétique.
\item \textbf{Lutte :} Prévention (alimentation, sport), Dépistage précoce (glycémie à jeun / HbA1c), Équilibre strict (HbA1c $< 7\%$), Éducation thérapeutique, Surveillance annuelle complications.
\end{itemize}
\end{aretenir}

\coursec{La pression artérielle : mesure et variations}

Après avoir analysé les désordres du métabolisme glucidique, notamment le diabète, nous abordons maintenant un autre pilier de la santé cardiovasculaire : la pression artérielle. Tout comme la glycémie, la pression artérielle doit être maintenue dans des intervalles étroits pour assurer la perfusion des organes. Une élévation chronique (hypertension artérielle) ou une chute brutale (hypotension) mettent en jeu le pronostic vital. Comprendre comment elle se mesure, quels paramètres la font varier et comment l'organisme la régule est indispensable pour prévenir les maladies cardiovasculaires, première cause de mortalité dans nos hôpitaux au Cameroun.

\courssub{Définition et valeurs de référence}

\begin{definition}[Pression artérielle]
La \cle{pression artérielle} (PA) est la force exercée par le sang sur la paroi des artères. Elle s'exprime en \si{cmHg} (centimètres de mercure) ou en \si{mmHg} (millimètres de mercure) et se caractérise par deux valeurs :
\begin{itemize}
    \item La \cle{pression systolique} (ou maximale, $P_{\text{sys}}$) : pression maximale atteinte lors de la contraction ventriculaire (systole), correspondant à l'éjection du sang dans l'aorte.
    \item La \cle{pression diastolique} (ou minimale, $P_{\text{dias}}$) : pression minimale maintenue lors du relâchement ventriculaire (diastole), assurée par l'élasticité des grosses artères (effet Windkessel).
\end{itemize}
On note classiquement $PA = P_{\text{sys}} / P_{\text{dias}}$.
\end{definition}

\begin{exemplebox}[Valeurs normales au repos]
Chez l'adulte sain, au repos, en position assise, les valeurs de référence sont :
\[
P_{\text{sys}} \approx \SI{12}{cmHg} \text{ à } \SI{13}{cmHg} \quad \text{et} \quad P_{\text{dias}} \approx \SI{8}{cmHg}
\]
On note couramment \textbf{12/8} ou \textbf{13/8}.
\begin{itemize}
    \item La \cle{pression différentielle} (ou pression de pouls) est $P_{\text{sys}} - P_{\text{dias}}$ (environ \SI{4}{cmHg}).
    \item La \cle{pression artérielle moyenne} (PAM) est la pression qui, si elle était constante, assurerait le même débit sanguin. Elle se calcule approximativement par $PAM \approx P_{\text{dias}} + \frac{1}{3}(P_{\text{sys}} - P_{\text{dias}})$. Pour 12/8 : $PAM \approx 8 + \frac{4}{3} \approx \SI{9,3}{cmHg}$.
\end{itemize}
\end{exemplebox}

\begin{attention}[Piège fréquent : Pression artérielle vs Fréquence cardiaque]
Ne confondez jamais \textbf{pression artérielle} (force sur la paroi, en \si{cmHg}) et \textbf{fréquence cardiaque} (nombre de battements par minute, en \si{bpm} ou \si{min^{-1}}). Le pouls que l'on palpe au poignet reflète la fréquence cardiaque, pas l'amplitude de la pression. Un sportif au repos peut avoir une FC de \SI{50}{bpm} et une PA de 11/7 ; un stressé peut avoir une FC de \SI{90}{bpm} et une PA de 14/9. Ce sont deux paramètres physiologiques distincts, quoique liés par le débit cardiaque.
\end{attention}

\courssub{Technique de mesure : le sphygmomanomètre}

La mesure clinique de routine utilise un \cle{sphygmomanomètre} (tensiomètre) à brassard, selon la méthode auscultatoire de Riva-Rocci/Korotkoff ou par oscillométrie (appareils électroniques).

\begin{methode}[Mesure correcte de la tension artérielle (Recommandations OMS/HTA Cameroun)]
Pour obtenir une valeur fiable, le protocole suivant est impératif :
\begin{enumerate}
    \item \textbf{Repos :} Le sujet est assis, dos appuyé, bras nu soutenu à la hauteur du cœur (niveau du 4e espace intercostal), jambes non croisées, après \textbf{5 minutes de repos} absolu (pas de café, tabac, effort, stress immédiat).
    \item \textbf{Brassard :} Choisir un brassard adapté à la circonférence du bras (largeur $\approx$ 40\% de la circonférence, longueur $\approx$ 80\%). Le bord inférieur du brassard est placé \SI{2}{cm} au-dessus du pli du coude. Le tuyau sort du côté du petit doigt (sur l'artère brachiale).
    \item \textbf{Gonflage :} Gonfler rapidement jusqu'à \SI{20}{cmHg} au-dessus de la disparition du pouls radial (ou \SI{30}{cmHg} au-dessus de la systolique estimée).
    \item \textbf{Dégonflage :} Dégonfler lentement (\SI{2}{cmHg} à \SI{3}{cmHg} par seconde).
    \item \textbf{Lecture (Méthode auscultatoire) :} Placer le pavillon du stéthoscope sur l'artère brachiale (dans le pli du coude, face médiale).
    \begin{itemize}
        \item Apparition des \textbf{bruits de Korotkoff} (phase I) $\rightarrow$ \cle{Pression Systolique}.
        \item Disparition des bruits (phase V) $\rightarrow$ \cle{Pression Diastolique}.
    \end{itemize}
    \item \textbf{Répétition :} Effectuer \textbf{deux mesures} espacées de 1 à 2 minutes. La moyenne des deux est retenue. Si écart $> \SI{5}{cmHg}$, faire une 3e mesure.
\end{enumerate}
\end{methode}

\begin{popfigure}
\begin{tikzpicture}[scale=0.9, every node/.style={font=\small}]
    % Bras
    \draw[fill=popcream!50, draw=popdark, thick] (0,0) ellipse (2.5cm and 5cm);
    % Artère brachiale (schématique)
    \draw[popred, very thick] (0.2,-4.5) .. controls (0.5,-3) and (0.5,1) .. (0.2,3);
    \draw[popred, very thick] (-0.2,-4.5) .. controls (-0.5,-3) and (-0.5,1) .. (-0.2,3);
    % Brassard
    \draw[fill=popblue!30, draw=popblue, thick] (-2.5,-1.5) rectangle (2.5,1.5);
    \draw[fill=popblue!50, draw=popblue] (-2.5,-1.5) rectangle (-2.0,1.5); % Velcro
    \draw[fill=popblue!50, draw=popblue] (2.0,-1.5) rectangle (2.5,1.5);
    % Tuyau brassard -> Manomètre
    \draw[popdark, thick] (2.5,0) -- (4,0) -- (4,2.5);
    % Manomètre
    \draw[fill=popgray!20, draw=popdark, thick] (3.5,2.5) circle (1.5cm);
    \draw[->, thick, popdark] (3.5,2.5) -- ++(90:1.2cm) node[above] {Aiguille};
    \foreach \angle/\val in {180/0, 150/2, 120/4, 90/6, 60/8, 30/10, 0/12, -30/14, -60/16, -90/18, -120/20, -150/22, -180/24, -210/26, -240/28} {
        \draw[popdark] (3.5,2.5) ++(\angle:1.1cm) -- ++(\angle:0.3cm);
        \node at ($(3.5,2.5)+(\angle:1.5cm)$) {\tiny \val};
    }
    \node[popdark, align=center] at (3.5, 4.5) {Manomètre\\(\si{cmHg})};
    % Poire de gonflage
    \draw[fill=poporange!30, draw=poporange, thick] (5.5,2.5) ellipse (1cm and 0.6cm);
    \node[poporange] at (5.5, 3.5) {Poire};
    \draw[popdark, thick] (5,0) -- (5.5,2.5);
    % Valve
    \draw[fill=poppurple!30, draw=poppurple] (4.5,0) rectangle (5,0.5);
    \node[poppurple, align=center] at (4.75, -0.5) {Vis de\\dégonflage};
    % Stéthoscope
    \draw[popgreen, thick] (-4, -2) -- (-4, 2) -- (-3, 2);
    \draw[fill=popgreen!30, draw=popgreen, thick] (-4, -2) circle (0.5cm);
    \node[popgreen, align=center] at (-4, -3) {Pavillon\\sur artère\\brachiale};
    % Flèches légendes
    \draw[->, thick, popdark] (-2.5, 1.5) -- ++(-1, 0.5) node[left, align=right] {Brassard\\niveau de cœur};
    \draw[->, thick, popred] (0.5, 2) -- ++(1, 0.5) node[right] {Artère brachiale};
    \draw[->, thick, popblue] (2.5, 1.5) -- ++(1, 0.5) node[right] {Tuyau};
\end{tikzpicture}
\legende{Schéma du dispositif de mesure de la pression artérielle par la méthode auscultatoire (sphygmomanomètre à mercure ou anéroïde + stéthoscope). Le brassard est positionné au niveau du cœur. Le stéthoscope écoute les bruits de Korotkoff au niveau de l'artère brachiale.}
\end{popfigure}

\begin{savaistu}[Les bruits de Korotkoff : une signature biophysique]
Lorsque le brassard est gonflé au-dessus de la pression systolique, l'artère est écoulée : \textbf{silence}. En dégonflant :
\begin{itemize}
    \item \textbf{Phase I (Systolique) :} La pression du brassard devient inférieure à la pression systolique. Le sang jaillit par à-coups dans l'artère partiellement comprimée $\rightarrow$ \textbf{bruits clairs, tapotants} (premier bruit = $P_{\text{sys}}$).
    \item \textbf{Phases II-III :} Bruits de souffle, puis plus forts (coup de marteau).
    \item \textbf{Phase IV :} Bruits sourds, étouffés (artère presque ouverte).
    \item \textbf{Phase V (Diastolique) :} \textbf{Silence total}. L'artère est complètement ouverte, flux laminaire non turbulent. Ce silence marque la \textbf{pression diastolique}.
\end{itemize}
Chez l'enfant ou l'insuffisant aortique, on utilise parfois la phase IV (étouffement) pour la diastole. Au Cameroun, la phase V est la référence standard.
\end{savaistu}

\courssub{Paramètres physiologiques déterminant la pression artérielle}

La pression artérielle moyenne (PAM) est le produit du débit cardiaque ($Q_c$) par les résistances vasculaires périphériques totales (RVP) :
\[
PAM = Q_c \times RVP
\]
Or $Q_c = V_{\text{es}} \times FC$ (Volume d'éjection systolique $\times$ Fréquence Cardiaque).
Les quatre piliers de la PA sont donc :

\begin{propriete}[Les 4 déterminants hémodynamiques de la PA]
\begin{enumerate}
    \item \textbf{Le Débit Cardiaque ($Q_c$) :} Volume de sang éjecté par minute. $Q_c = V_{\text{es}} \times FC$.
        \begin{itemize}
            \item $\uparrow FC$ ou $\uparrow V_{\text{es}}$ (ex: effort, catécholamines) $\rightarrow$ $\uparrow P_{\text{sys}}$ surtout.
        \end{itemize}
    \item \textbf{Le Volume Sanguin Circulant (Volémie) :} Quantité totale de sang ($\approx \SI{5}{L}$ chez l'adulte).
        \begin{itemize}
            \item $\uparrow$ Volémie (rétention $\text{Na}^+$/eau, excès sel) $\rightarrow$ $\uparrow$ Pression de remplissage $\rightarrow$ $\uparrow V_{\text{es}}$ (loi de Frank-Starling) $\rightarrow$ $\uparrow PA$.
            \item $\downarrow$ Volémie (hémorragie, déshydratation sévère) $\rightarrow$ $\downarrow PA$ (choc hypovolémique).
        \end{itemize}
    \item \textbf{Les Résistances Périphériques (RVP) :} Dépendent principalement du \textbf{diamètre des artérioles} (loi de Poiseuille : $R \propto 1/r^4$).
        \begin{itemize}
            \item Vasoconstriction artériolaire (noradrénaline, angiotensine II, endothéline) $\rightarrow$ $\uparrow \uparrow RVP$ $\rightarrow$ $\uparrow \uparrow P_{\text{dias}}$ et $\uparrow P_{\text{sys}}$.
            \item Vasodilatation (NO, prostaglandines, bradykinine, chaleur) $\rightarrow$ $\downarrow RVP$ $\rightarrow$ $\downarrow PA$.
        \end{itemize}
    \item \textbf{L'Élasticité Artérielle (Compliance) :} Capacité des grosses artères (aorte, carotides) à se distendre en systole et se recoiler en diastole (effet Windkessel).
        \begin{itemize}
            \item Artères rigides (athérosclérose, vieillissement, diabète) $\rightarrow$ $\downarrow$ Amortissement $\rightarrow$ $\uparrow \uparrow P_{\text{sys}}$ et $\downarrow P_{\text{dias}}$ $\rightarrow$ $\uparrow$ Pression différentielle (signe de gravité).
        \end{itemize}
\end{enumerate}
\end{propriete}

\begin{experience}[Mise en évidence des facteurs de variation (TP type)]
\textbf{Objectif :} Observer l'effet de l'effort, de la posture et de l'apnée sur la PA.
\textbf{Matériel :} Tensiomètre électronique validé, stepper ou marche, chronomètre.
\textbf{Protocole :}
\begin{enumerate}
    \item Mesure PA au repos assis (M0).
    \item Effort standardisé : monter/descendre un step (\SI{30}{cm}) pendant \SI{3}{min} (ou course sur place).
    \item Mesure immédiate à l'arrêt (M1), puis à \SI{1}{min} (M2), \SI{3}{min} (M3), \SI{5}{min} (M4) de récupération assis.
    \item Série 2 : Mesure assis $\rightarrow$ debout immédiat $\rightarrow$ \SI{3}{min} debout (effet orthostatique).
    \item Série 3 : Apnée inspiratoire forcée (manœuvre de Valsalva) pendant mesure continue (si tensiomètre continu type Finapres) ou observation clinique.
\end{enumerate}
\textbf{Résultats attendus (exemple élève 18 ans) :}
\begin{center}
\begin{tabularx}{\linewidth}{|l|X|X|X|X|X|}\hline
\textbf{Moment} & \textbf{Repos (M0)} & \textbf{Effort (M1)} & \textbf{Récup 1 min (M2)} & \textbf{Récup 3 min (M3)} & \textbf{Récup 5 min (M4)} \\ \hline
PA Systolique (\si{cmHg}) & 12 & 16 & 14 & 12 & 12 \\ \hline
PA Diastolique (\si{cmHg}) & 8 & 8 & 9 & 8 & 8 \\ \hline
FC (\si{bpm}) & 70 & 160 & 110 & 85 & 72 \\ \hline
\end{tabularx}
\end{center}
\textbf{Interprétation :}
\begin{itemize}
    \item \textbf{Effort :} $\uparrow \uparrow FC$ et $\uparrow V_{\text{es}}$ $\rightarrow$ $\uparrow \uparrow Q_c$ $\rightarrow$ $\uparrow \uparrow P_{\text{sys}}$ (12 $\rightarrow$ 16). $P_{\text{dias}}$ stable (vasodilatation musculaire compense $\uparrow Q_c$).
    \item \textbf{Récupération :} $P_{\text{sys}}$ redescend exponentiellement. $P_{\text{dias}}$ peut monter transitoirement (vasoconstriction réflexe post-effort, thermorégulation).
    \item \textbf{Posture (Debout) :} Chute transitoire PA (gravité $\rightarrow$ $\downarrow$ retour veineux $\rightarrow$ $\downarrow Q_c$), puis normalisation par réflexe barorécepteur (tachycardie, vasoconstriction).
\end{itemize}
\end{experience}

\courssub{Variations physiologiques de la pression artérielle}

La PA n'est pas une constante ; elle s'adapte en permanence aux besoins de l'organisme. Toute variation transitoire justifiée qui revient à la normale au repos témoigne d'une \textbf{régulation efficace}.

\begin{exemplebox}[Principales causes de variations physiologiques]
\begin{center}
\begin{tabularx}{\linewidth}{|l|X|X|}\hline
\textbf{Facteur} & \textbf{Mécanisme principal} & \textbf{Effet typique sur PA} \\ \hline
\textbf{Effort physique dynamique} & $\uparrow \uparrow FC$, $\uparrow V_{\text{es}}$ ($\uparrow Q_c$) ; Vasodilatation musculaire ($\downarrow RVP$) & $\uparrow \uparrow P_{\text{sys}}$ ; $P_{\text{dias}}$ stable ou $\downarrow$ légère \\ \hline
\textbf{Effort statique (port de charge)} & $\uparrow$ Pression intrathoracique/abdominale ; $\uparrow$ Tonus sympathique ($\uparrow RVP$) & $\uparrow P_{\text{sys}}$ ET $\uparrow P_{\text{dias}}$ \\ \hline
\textbf{Émotion / Stress aigu} & Décharge adrénergique (adrénaline, noradrénaline) : $\uparrow FC$, $\uparrow Contractilité$, $\uparrow RVP$ (cutané/splanchnique) & $\uparrow P_{\text{sys}}$ et $\uparrow P_{\text{dias}}$ \\ \hline
\textbf{Posture : Passage couché $\rightarrow$ debout} & $\downarrow$ Retour veineux (gravité) $\rightarrow$ $\downarrow Q_c$ transitoire & Chute transitoire (orthostasie) puis normalisation (réflexe barorécepteur) \\ \hline
\textbf{Ingestion de sel (NaCl) chronique} & $\uparrow$ Osmolarité plasmatique $\rightarrow$ Soif + $\uparrow$ ADH + $\uparrow$ Aldostérone $\rightarrow$ $\uparrow$ Volémie & $\uparrow PA$ à long terme (facteur majeur d'HTA au Cameroun : sel de cuisine, cubes bouillon, poisson séché) \\ \hline
\textbf{Âge} & $\downarrow$ Élasticité artérielle (rigidité aortique) & $\uparrow P_{\text{sys}}$, $\downarrow P_{\text{dias}}$ $\rightarrow$ $\uparrow$ Pression de pouls (HTA systolique isolée du sujet âgé) \\ \hline
\textbf{Nycthémère} & Rythme circadien (cortisol, système nerveux autonome) & Minimum la nuit (\textit{dipper} : -10 à -20\%) ; Pic matinal (réveil) \\ \hline
\end{tabularx}
\end{center}
\end{exemplebox}

\courssub{Analyse de courbes : Pression artérielle au cours d'un effort et récupération}

L'analyse de courbes physiologiques est un savoir-faire clé du Bac. Il faut savoir décrire, quantifier (valeurs, temps) et interpréter mécanistiquement.

\begin{popfigure}
\begin{tikzpicture}
\begin{axis}[
    width=\linewidth,
    height=9cm,
    axis lines=middle,
    xlabel={Temps (min)},
    ylabel={Pression artérielle (\si{cmHg})},
    xmin=-0.5, xmax=12.5,
    ymin=5, ymax=18,
    xtick={0,3,6,7,8,9,10,11,12},
    xticklabels={Repos, D\'ebut effort, Fin effort (3 min), R\'ecup 1 min, R\'ecup 2 min, R\'ecup 3 min, R\'ecup 4 min, R\'ecup 5 min, R\'ecup 6 min},
    ytick={6,8,10,12,14,16,18},
    legend pos=north east,
    grid=major,
    grid style={popline},
    popcurve/.style={thick, mark=*, mark size=2},
]
% Systolique : Repos 0-3 (12), Effort 3-6 (montée à 16), Récup 6-12 (descente à 12)
\addplot[popred, popcurve] coordinates {
    (0,12) (3,12) (4,13.5) (5,15) (6,16) (7,14.5) (8,13.5) (9,12.5) (10,12.2) (11,12) (12,12)
};
% Diastolique : Repos 0-3 (8), Effort 3-6 (légère montée à 9), Récup 6-12 (retour à 8)
\addplot[popblue, popcurve] coordinates {
    (0,8) (3,8) (4,8.2) (5,8.5) (6,9) (7,8.8) (8,8.5) (9,8.2) (10,8) (11,8) (12,8)
};
\legend{$P_{\text{sys}}$ (Systolique), $P_{\text{dias}}$ (Diastolique)}
% Annotations
\draw[popred, dashed] (3,5) -- (3,12);
\draw[popred, dashed] (6,5) -- (6,16);
\node[popred, anchor=south] at (4.5, 17) {Effort dynamique (3 min)};
\node[align=center, anchor=north] at (0, 11.5) {Repos\\12/8};
\node[align=center, anchor=south] at (6, 16.5) {Pic effort\\16/9};
\node[align=center, anchor=north] at (12, 11.5) {Récupération\\12/8};
\end{axis}
\end{tikzpicture}
\legende{Évolution type de la pression artérielle systolique (rouge) et diastolique (bleu) au cours d'un effort dynamique modéré (3 min) et de la récupération (6 min) chez un sujet jeune sain. On note l'augmentation nette de la systolique liée à l'augmentation du débit cardiaque, la stabilité relative de la diastolique (vasodilatation musculaire), et le retour progressif aux valeurs de repos.}
\end{popfigure}

\begin{exoresolu}[Analyse quantitative d'une courbe PA/Effort]
\textbf{Énoncé :} Un élève de Terminale D, âgé de 17 ans, réalise un test d'effort sur cyclo-ergomètre. La figure ci-dessus (ou un tableau de valeurs) donne l'évolution de sa PA. 
\begin{enumerate}
    \item Relever les valeurs de PA au repos, au pic d'effort (fin de l'effort, 3 min) et à la 5e minute de récupération.
    \item Calculer la pression différentielle (pression de pouls) à ces trois instants.
    \item Expliquer les variations observées de $P_{\text{sys}}$ et $P_{\text{dias}}$ pendant l'effort en lien avec les déterminants hémodynamiques.
    \item Conclure sur l'adaptation cardiovasculaire de ce sujet.
\end{enumerate}

\textbf{Solution détaillée :}
\begin{enumerate}
    \item \textbf{Lecture :}
    \begin{itemize}
        \item Repos (t=0) : $P_{\text{sys}} = \SI{12}{cmHg}$, $P_{\text{dias}} = \SI{8}{cmHg}$ $\rightarrow$ PA = \textbf{12/8}.
        \item Pic effort (t=6 min, fin effort) : $P_{\text{sys}} = \SI{16}{cmHg}$, $P_{\text{dias}} = \SI{9}{cmHg}$ $\rightarrow$ PA = \textbf{16/9}.
        \item Récup 5 min (t=11 min) : $P_{\text{sys}} \approx \SI{12}{cmHg}$, $P_{\text{dias}} \approx \SI{8}{cmHg}$ $\rightarrow$ PA $\approx$ \textbf{12/8}.
    \end{itemize}
    \item \textbf{Pression différentielle ($\Delta P = P_{\text{sys}} - P_{\text{dias}}$) :}
    \begin{align*}
        \text{Repos : } & 12 - 8 = \SI{4}{cmHg} \\
        \text{Pic effort : } & 16 - 9 = \SI{7}{cmHg} \quad (\text{augmentation marquée}) \\
        \text{Récup 5 min : } & 12 - 8 = \SI{4}{cmHg} \quad (\text{retour à la normale})
    \end{align*}
    \item \textbf{Interprétation mécanistique :}
    \begin{itemize}
        \item \textbf{$\uparrow P_{\text{sys}}$ (12 $\rightarrow$ 16) :} Due à l'augmentation massive du \textbf{débit cardiaque} ($Q_c = FC \times V_{\text{es}}$). La FC passe de $\sim 70$ à $\sim 160$ \si{bpm} ; le $V_{\text{es}}$ augmente (loi de Frank-Starling + contractilité accrue par catécholamines). L'éjection rapide et volumineuse dans l'aorte élève le pic de pression.
        \item \textbf{$P_{\text{dias}}$ stable/légèrement $\uparrow$ (8 $\rightarrow$ 9) :} Deux effets opposés : (1) L'augmentation du $Q_c$ tend à monter la diastole. (2) La \textbf{vasodilatation artériolaire massive dans les muscles actifs} (métabolites : adénosine, $\text{K}^+$, $\text{H}^+$, NO) fait chuter les \textbf{résistances périphériques totales (RVP)}. L'effet (2) compense presque l'effet (1), stabilisant la diastole. C'est la signature de l'effort \textbf{dynamique} (aérobie).
    \end{itemize}
    \item \textbf{Conclusion :} Le sujet présente une \textbf{réponse hémodynamique normale à l'effort dynamique} : augmentation adaptée de la systolique pour perfuser les muscles, stabilité diastolique traduisant une bonne capacité de vasodilatation périphérique, et récupération rapide (retour aux valeurs de repos en $< 5$ min), signe d'un bon conditionnement cardiovasculaire et d'une régulation barorécepteur efficace.
\end{enumerate}
\end{exoresolu}

\begin{attention}[Erreurs d'interprétation à ne pas commettre au Bac]
\begin{itemize}
    \item \textbf{Dire "la tension monte" sans distinguer systolique et diastolique.} Précisez toujours laquelle varie et dans quel sens.
    \item \textbf{Confondre effort dynamique et statique.} À l'effort statique (haltérophilie, port de charge), $P_{\text{dias}}$ monte fortement (compression vaisseaux + réflexe pressor) $\rightarrow$ PA 18/12 par ex. À l'effort dynamique (course, vélo), $P_{\text{dias}}$ ne monte pas (ou baisse).
    \item \textbf{Oublier l'unité (\si{cmHg} ou \si{mmHg}).} Une valeur sans unité vaut 0 point.
    \item \textbf{Ne pas relier la variation au déterminant ($Q_c$, RVP, Volémie, Élasticité).} L'explication doit être causale : "L'augmentation de la FC entraîne une augmentation du $Q_c$ qui explique la hausse de $P_{\text{sys}}$".
\end{itemize}
\end{attention}

\courssub{Synthèse : de la mesure à la pathologie}

La mesure de la PA est un acte médical banal en apparence, mais qui requiert rigueur (repos, brassard, technique) pour éviter le "syndrome de la blouse blanche" (élévation transitoire due au stress du cabinet médical). Les valeurs normales (12/8) masquent une dynamique complexe pilotée par le débit cardiaque, la volémie, le calibre artériolaire et la rigidité artérielle. 

L'effort physique est le test physiologique par excellence : il révèle la réserve cardiaque (montée de $P_{\text{sys}}$) et la réserve vasculaire (stabilité de $P_{\text{dias}}$). Une $P_{\text{sys}}$ qui ne monte pas (incompétence cardiaque), une $P_{\text{dias}}$ qui monte à l'effort dynamique (raideur artérielle précoce, obésité, début d'HTA), ou une récupération lente ($> \SI{10}{min}$) sont des signaux d'alerte.

Dans la section suivante, nous verrons comment les \textbf{barorécepteurs} (carotide, aorte) et les \textbf{systèmes hormonaux} (rénine-angiotensine-aldostérone, ADH, péptides natriurétiques) orchestrent la régulation minute par minute, et pourquoi leur défaillance conduit à l'hypertension artérielle (HTA) ou à l'hypotension orthostatique, deux fléaux de santé publique au Cameroun.

\coursec{Régulation de la pression artérielle}

Dans la section précédente, tu as appris à mesurer la pression artérielle avec le sphygmomanomètre et tu as constaté qu'elle n'est pas constante : elle augmente à l'effort, au stress, et diminue au repos ou pendant le sommeil. Pourtant, malgré ces variations, la pression artérielle d'un individu sain reste remarquablement stable autour de valeurs moyennes (environ \SI{12}{cm} de mercure, soit \SI{120}{mmHg}, pour la pression systolique ; \SI{80}{mmHg} pour la diastolique). Comment l'organisme parvient-il à corriger si vite les écarts ? C'est le problème que nous abordons maintenant : la \cle{régulation de la pression artérielle}. Nous verrons qu'elle repose sur deux mécanismes complémentaires : une régulation \cle{nerveuse}, très rapide (secondes), et une régulation \cle{hormonale}, plus lente (minutes à heures) mais durable.

\courssub{Mise en évidence expérimentale de la régulation nerveuse}

\textbf{Observation et question.} Lorsqu'un animal (ou un être humain) passe brutalement de la position couchée à la position debout, le sang tend à s'accumuler dans les membres inférieurs sous l'effet de la pesanteur : le retour veineux diminue, la pression artérielle au niveau de la tête devrait chuter. Or, en quelques secondes, elle est rétablie. Quel mécanisme assure ce rétablissement ?

\begin{experience}[Stimulation et section du nerf vague et des nerfs sympathiques chez l'animal]
\textbf{Matériel :} un chien anesthésié, un capteur de pression artérielle (artère fémorale ou carotidienne) relié à un enregistreur (kymographe ou acquisition numérique), électrodes de stimulation, dispositif de mesure de la fréquence cardiaque (ECG).

\textbf{Protocole et résultats :}
\begin{enumerate}
\item On enregistre la pression artérielle moyenne (PAM) et la fréquence cardiaque (FC) de l'animal au repos : valeurs stables (par exemple PAM = \SI{110}{mmHg}, FC = 90 battements/min).
\item On \textbf{stimule électriquement le nerf vague} (nerf pneumogastrique, nerf crânien X, voie parasympathique) : on observe un \textbf{ralentissement marqué du cœur} (bradycardie, FC \textdownarrow) et une \textbf{chute de la pression artérielle} (PAM \textdownarrow).
\item On \textbf{sectionne les deux nerfs vagues} : le cœur s'accélère nettement (tachycardie, FC \textuparrow) et la pression artérielle s'élève durablement (PAM \textuparrow).
\item On \textbf{stimule un nerf cardiaque accélérateur} (voie orthosympathique) : le cœur s'accélère (tachycardie, inotropisme positif) et la pression artérielle augmente (vasoconstriction artériolaire).
\end{enumerate}

\textbf{Interprétation :} Le cœur et les vaisseaux (artérioles) reçoivent en permanence deux influx nerveux antagonistes : le nerf vague (parasympathique) freine le cœur (effet chronotrope négatif), tandis que les nerfs orthosympathiques l'accélèrent (chronotrope positif), augmentent sa force de contraction (inotropisme positif) et provoquent la constriction des artérioles (augmentation des résistances périphériques totales). La pression artérielle au repos résulte de l'équilibre dynamique entre ces deux influences toniques.

\textbf{Conclusion :} Il existe une régulation nerveuse de l'activité cardiaque (fréquence et force) et du calibre des vaisseaux (résistances périphériques), donc de la pression artérielle.
\end{experience}

Mais comment le système nerveux « sait-il » que la pression varie ? Une seconde série d'expériences apporte la réponse : si l'on sectionne les nerfs afférents issus de la \textbf{crosse de l'aorte} (nerfs vagues X) et des \textbf{sinus carotidiens} (nerfs glossopharyngiens IX), l'animal devient incapable de corriger les variations de pression : celle-ci fluctue de façon anarchique (labilité tensionnelle). Inversement, la distension expérimentale d'un sinus carotidien (par compression externe en amont) provoque immédiatement un ralentissement du cœur et une baisse de pression. Il existe donc, dans ces parois artérielles, des \cle{récepteurs sensibles à la pression}.

\begin{definition}[Les barorécepteurs]
Les \cle{barorécepteurs} (ou récepteurs barosensibles) sont des terminaisons nerveuses sensitives (mécanorécepteurs) sensibles à l'\textbf{étirement des parois artérielles} induit par la pression sanguine. Ils sont situés principalement dans la paroi de la \textbf{crosse aortique} (innervés par le nerf vague X) et dans les \textbf{sinus carotidiens} (renflements à la bifurcation des artères carotides communes, innervés par le nerf glossopharyngien IX). Plus la pression artérielle est élevée, plus la paroi est étirée, plus la fréquence des potentiels d'action émis par les barorécepteurs vers les centres nerveux est élevée.
\end{definition}

\begin{popfigure}
\begin{center}
\begin{tikzpicture}[scale=1.0, every node/.style={font=\small}, >=Stealth]
% Coeur
\draw[fill=poppinkL, draw=poppink, thick] (0,0) ellipse (1.1 and 1.4);
\node at (0,0) {\textbf{C\oe ur}};
% Aorte ascendante et crosse
\draw[draw=popink, line width=2.2pt] (0.2,1.3) -- (0.2,2.2) 
    arc[start angle=0, end angle=150, radius=1.1] (-1.1,2.9) -- (-1.1,3.6);
\draw[draw=popink, line width=2.2pt] (0.2,2.2) arc[start angle=0, end angle=110, radius=1.1];
% Tronc brachiocéphalique et carotides
\draw[draw=popink, line width=1.6pt] (-0.9,2.75) -- (-1.6,4.8);
\draw[draw=popink, line width=1.6pt] (-0.55,2.95) -- (0.4,4.8);
% Sinus carotidiens (renflements)
\draw[fill=poporangeL, draw=poporange, thick] (-1.6,3.8) ellipse (0.22 and 0.35);
\draw[fill=poporangeL, draw=poporange, thick] (0.4,3.8) ellipse (0.22 and 0.35);
% Barorécepteurs crosse aortique (zone)
\draw[fill=poporange, draw=poporange] (-0.75,3.15) circle (0.1);
% Nerfs afférents vers le bulbe (bleu)
\draw[draw=popblue, thick, ->] (-1.6,4.15) -- (-1.6,5.5) node[midway, left, font=\footnotesize] {IX};
\draw[draw=popblue, thick, ->] (0.4,4.15) -- (0.4,5.5) node[midway, right, font=\footnotesize] {IX};
\draw[draw=popblue, thick, ->] (-0.75,3.35) -- (-0.75,4.5) -- (-0.6,5.5) node[midway, left, font=\footnotesize] {X};
% Bulbe rachidien
\draw[fill=popblueL, draw=popblue, thick] (-0.6,5.5) ellipse (1.8 and 0.6);
\node at (-0.6,5.5) {\textbf{Bulbe rachidien}};
\node[font=\footnotesize, popdark] at (-0.6,5.15) {(Centres vasomoteurs)};
% Nerfs efférents
\draw[draw=popgreen, thick, ->] (-1.9,5.3) .. controls (-3.2,4.0) and (-3.2,1.5) .. (-1.4,0.4);
\node[font=\footnotesize, popgreen, align=center] at (-3.5,2.6) {Nerf vague (X)\\(Parasympathique)};
\draw[draw=popgold, thick, ->] (0.7,5.3) .. controls (2.5,4.0) and (2.7,1.5) .. (1.2,0.6);
\node[font=\footnotesize, popgold, align=center] at (3.1,2.6) {Nerfs orthosympathiques};
% Légendes pointées
\node[font=\footnotesize, align=left, anchor=west] at (1.9,4.5) {\textcircled{1} Sinus carotidien};
\draw[popmuted, thin] (1.9,4.4) -- (0.62,3.9);
\node[font=\footnotesize, align=left, anchor=east] at (-4.8,3.5) {\textcircled{2} Crosse aortique};
\draw[popmuted, thin] (-2.4,3.5) -- (-0.85,3.2);
\end{tikzpicture}
\end{center}
\legende{Localisation des barorécepteurs et voies nerveuses de la régulation. \textcircled{1} Les sinus carotidiens (renflements à l'origine des carotides internes) et \textcircled{2} la crosse aortique contiennent les barorécepteurs. Les influx afférents (flèches bleues) gagnent les centres vasomoteurs du bulbe rachidien via les nerfs glossopharyngiens (IX) et vagues (X). Les réponses efférentes empruntent le nerf vague (freinateur cardiaque, vert) et les nerfs orthosympathiques (excitateurs cardiaques et vasoconstricteurs, doré).}
\end{popfigure}

\courssub{La régulation nerveuse rapide : le réflexe barorécepteur (baroréflexe)}

Le mécanisme mis en évidence est un \cle{réflexe} : on parle de \textbf{réflexe barorécepteur} (ou baroréflexe). C'est un arc réflexe complet : récepteur $\rightarrow$ voie afférente $\rightarrow$ centre $\rightarrow$ voie efférente $\rightarrow$ effecteur. Décrivons son fonctionnement dans les deux situations possibles.

\textbf{Cas d'une élévation de la pression artérielle} (par exemple lors d'un effort musculaire intense, stress aigu) :
\begin{enumerate}
\item L'élévation de la pression étire davantage les parois de la crosse aortique et des sinus carotidiens ;
\item les barorécepteurs, davantage stimulés (seuil de déclenchement \textasciitilde \SI{60}{mmHg}), augmentent la fréquence des influx nerveux envoyés aux \textbf{centres vasomoteurs du bulbe rachidien} (partie caudale du tronc cérébral) ;
\item le bulbe répond en \textbf{activant la sortie parasympathique} (noyau ambigu $\rightarrow$ nerf vague) et en \textbf{inhibant la sortie orthosympathique} (centre vasomoteur latéral) ;
\item conséquences effectrices : \textbf{ralentissement cardiaque} (bradycardie réflexe, effet chronotrope négatif), légère diminution de la force de contraction, et \textbf{vasodilatation} des artérioles (diminution des résistances périphériques totales) par inhibition du tonus sympathique vasomoteur ;
\item le débit cardiaque et les résistances périphériques diminuent : la \textbf{pression artérielle baisse} et revient vers sa valeur de consigne (rétrocontrôle négatif).
\end{enumerate}

\textbf{Cas d'une chute de la pression artérielle} (par exemple hémorragie, lever brutal, déshydratation) : la pression diminue, les parois artérielles sont moins étirées, la fréquence de décharge des barorécepteurs \textbf{diminue} (ils sont toniquement actifs au repos) ; le bulbe rachidien \textbf{inhibe la sortie vagale} (levée du frein) et \textbf{active massivement la sortie orthosympathique} ; il en résulte une \textbf{accélération cardiaque} (tachycardie réflexe), une augmentation de la contractilité (inotropisme positif) et une \textbf{vasoconstriction} artériolaire généralisée (augmentation des résistances périphériques) : la pression remonte vers la normale.

\begin{popfigure}
\begin{center}
\begin{tikzpicture}[scale=0.9]
\begin{axis}[
    width=\linewidth, height=7cm,
    axis lines=middle,
    xlabel={Pression artérielle moyenne (mmHg)}, ylabel={Fréquence de décharge (imp/s) ou FC (bpm)},
    xmin=40, xmax=180, ymin=0, ymax=180,
    xtick={60,100,140}, ytick={0,50,100,150},
    grid=major, grid style={popline},
    legend style={at={(0.98,0.98)}, anchor=north east, fill=popcream, draw=popmuted, font=\footnotesize},
    clip=false
]
% Courbe barorécepteurs (sigmoïde)
\addplot[popblue, thick, domain=40:180, samples=200] {150/(1+exp(-(x-100)/15))};
\addlegendentry{Décharge barorécepteurs (afférente)}
% Courbe FC (inverse)
\addplot[popink, thick, domain=40:180, samples=200] {180 - 150/(1+exp(-(x-100)/15))};
\addlegendentry{Fréquence cardiaque (efférente)}
% Zone physiologique
\addplot[popgreenL!50, fill opacity=0.3, domain=80:130] {0} \closedcycle;
\node[font=\footnotesize, popgreen, align=center] at (axis cs:105,20) {Zone de\\régulation physiologique};
% Point de consigne
\draw[popmuted, dashed] (axis cs:100,0) -- (axis cs:100,150);
\node[font=\footnotesize, popdark] at (axis cs:100,165) {Consigne \textasciitilde 100 mmHg};
\end{axis}
\end{tikzpicture}
\end{center}
\legende{Relation sigmoïde entre la pression artérielle moyenne et l'activité des barorécepteurs (bleu) / fréquence cardiaque (rose). Autour de la consigne (\textasciitilde 100 mmHg), une faible variation de pression entraîne une grande variation de la décharge afférente et donc une correction efficace de la FC (gain du baroréflexe maximal). En dehors de cette zone (hypertension sévère ou choc), le réflexe sature.}
\end{popfigure}

\begin{attention}[Pièges classiques au Baccalauréat]
\begin{itemize}
\item \textbf{Localisation du centre :} Erreur très fréquente : situer le centre de régulation de la PA dans l'\textbf{hypothalamus}. Le centre nerveux \emph{directement} responsable (intégration du réflexe) est le \textbf{bulbe rachidien} (moelle allongée), qui abrite les \textbf{centres vasomoteurs} (centre cardio-accélérateur sympathique, centre cardio-inhibiteur parasympathique / noyau ambigu, centre vasomoteur sympathique). L'hypothalamus module l'activité bulbaire (émotions, thermorégulation) mais n'est pas le centre du baroréflexe.
\item \textbf{Barorécepteurs vs Chimiorécepteurs :} Ne pas confondre. Les barorécepteurs (crosse aortique, sinus carotidiens) sont sensibles à la \textbf{pression} (étirement). Les chimiorécepteurs (corps aortique, corps carotidien, même topographie) sont sensibles aux variations de \ce{O2}, \ce{CO2} et pH (régulation respiratoire principalement).
\item \textbf{Sens de la variation :} Une \textbf{hausse} de PA $\rightarrow$ \textbf{augmentation} de l'influx afférent $\rightarrow$ activation vagale / inhibition sympathique $\rightarrow$ \textbf{baisse} de la PA. Une \textbf{baisse} de PA $\rightarrow$ \textbf{diminution} de l'influx afférent (levée de l'inhibition tonique) $\rightarrow$ inhibition vagale / activation sympathique $\rightarrow$ \textbf{hausse} de la PA. C'est un rétrocontrôle \textbf{négatif}.
\item \textbf{Compression « en aval » :} Dans l'exercice type (compression carotidienne en aval du sinus), on \textbf{diminue} la pression dans le sinus $\rightarrow$ les barorécepteurs sont \textbf{moins} étirés $\rightarrow$ l'organisme « croit » à une hypotension $\rightarrow$ il \textbf{monte} la PA. Bien distinguer « en amont » (pression monte dans le sinus) et « en aval » (pression baisse dans le sinus).
\end{itemize}
\end{attention}

\courssub{La régulation hormonale lente (maintien à long terme)}

La régulation nerveuse corrige les variations brutales en quelques secondes (baroréflexe). Mais les barorécepteurs s'adaptent (phénomène d'\textbf{adaptation} ou \textbf{réinitialisation} du seuil) en 1 à 2 jours : en cas d'hypertension chronique, ils considèrent la nouvelle valeur comme « normale » et ne corrigent plus. Pour maintenir la pression sur des durées longues (minutes, heures, jours, vie), l'organisme fait appel à des \cle{hormones} qui agissent sur le \textbf{volume sanguin} (volémie) et le \textbf{tonus vasculaire} de fond.

\textbf{1. L'adrénaline et la noradrénaline (médullosurrénale).} En situation de stress aigu, peur, effort intense, douleur, les centres nerveux (hypothalamus, bulbe) stimulent, par la voie orthosympathique (fibres pré-ganglionnaires), la \textbf{moelle des glandes surrénales} (ganglion sympathique modifié), qui libèrent dans le sang \textbf{adrénaline} (\textasciitilde 80\%) et \textbf{noradrénaline} (\textasciitilde 20\%). Ces catécholamines potentialisent et prolongent l'action sympathique : $\beta_1$ cardiaque (FC \textuparrow, force \textuparrow), $\alpha_1$ vasculaire (vasoconstriction cutanée, splanchnique, rénale), $\beta_2$ musculaire (vasodilatation modérée). Leur action, transmise par la circulation générale, est plus lente (latence \textasciitilde 10-30 s) mais plus durable que l'influx nerveux direct.

\textbf{2. Le système rénine-angiotensine-aldostérone (SRAA).} Le rein joue le rôle central dans la régulation à long terme car il contrôle le \textbf{volume sanguin circulant} (volémie) via l'excrétion urinaire.

\begin{propriete}[Système rénine-angiotensine-aldostérone (SRAA) — Admis]
Lorsque la pression de perfusion rénale baisse (détection par les cellules juxtaglomérulaires de l'appareil juxtaglomérulaire) ou sous l'effet de la stimulation sympathique ($\beta_1$), le rein sécrète l'enzyme \textbf{rénine}. La rénine clive l'\textbf{angiotensinogène} (protéine hépatique, $\alpha_2$-globuline) en \textbf{angiotensine I} (inactive). L'\textbf{enzyme de conversion (ECA)}, présente principalement à la surface des cellules endothéliales pulmonaires, transforme l'angiotensine I en \textbf{angiotensine II} (peptide actif puissant).
L'angiotensine II a un double effet majeur :
\begin{itemize}
\item \textbf{Vasoconstriction} artériolaire intense et rapide (récepteurs AT1) $\rightarrow$ augmentation immédiate des résistances périphériques totales (RPT).
\item \textbf{Stimulation de la zone glomérulée du cortex surrénalien} $\rightarrow$ sécrétion d'\textbf{aldostérone} (minéralocorticoïde).
\end{itemize}
L'aldostérone agit sur les cellules principales du tube distal et du canal collecteur rénal : elle augmente la \textbf{réabsorption active du sodium} (\ce{Na+}) et, par suite osmotique, la \textbf{réabsorption d'eau}. Le volume urinaire diminue, le \textbf{volume sanguin (volémie) augmente}, le débit cardiaque augmente $\rightarrow$ élévation durable de la pression artérielle.
\textit{Note : La connaissance détaillée de la cascade enzymatique (ECA) est admise ; sa justification expérimentale n'est pas exigible au Bac.}
\end{propriete}

\textbf{3. L'hormone antidiurétique (ADH ou Vasopressine).} Synthétisée par les noyaux supra-optique et paraventriculaire de l'\textbf{hypothalamus}, transportée le long de l'axe neuro-hypophysaire, stockée et libérée par la \textbf{post-hypophyse} (neurohypophyse). Sa libération est stimulée par :
\begin{itemize}
\item Une \textbf{baisse de la pression artérielle} (détectée par les barorécepteurs $\rightarrow$ bulbe $\rightarrow$ hypothalamus) : c'est le stimulus principal en cas d'hémorragie.
\item Une \textbf{augmentation de l'osmolarité plasmatique} (osmorécepteurs hypothalamiques) : déshydratation.
\end{itemize}
L'ADH agit sur les récepteurs V2 des cellules du canal collecteur rénal : insertion de canaux à eau (aquaporines-2) $\rightarrow$ \textbf{réabsorption d'eau libre} $\rightarrow$ urine concentrée, volume sanguin maintenu/augmenté $\rightarrow$ PA maintenue. À fortes doses, l'ADH a aussi un effet vasoconstricteur (récepteurs V1 sur muscles lisses vasculaires), d'où son nom de vasopressine.

\begin{methode}[Distinguer et articuler les deux régulations]
Pour analyser un document ou répondre à une question, pose-toi systématiquement ces trois questions :
\begin{enumerate}
\item \textbf{Quelle est la latence / vitesse de la réponse ?}
      \begin{itemize}
      \item Effet en \textbf{secondes} (lever brutal, effort soudain, choc émotionnel) $\rightarrow$ Régulation \textbf{nerveuse} (baroréflexe : c\oe ur + vaisseaux).
      \item Effet en \textbf{minutes à heures} (compensation hémorragie, déshydratation, régime sans sel, stress chronique) $\rightarrow$ Régulation \textbf{hormonale} (surrénales, reins, hypophyse).
      \end{itemize}
\item \textbf{Quels sont les effecteurs ciblés ?}
      \begin{itemize}
      \item C\oe ur (FC, force) + Calibre artérioles (RPT) $\rightarrow$ Principalement nerveuse (+ adrénaline).
      \item Volume sanguin (volémie) + Excrétion urinaire + Tonus vasculaire de fond $\rightarrow$ Hormonale (Rénine/Angio II/Aldo, ADH, Adrénaline).
      \end{itemize}
\item \textbf{Quelle est la durée d'action et l'adaptation ?}
      \begin{itemize}
      \item Courte, réversible instantanément, sujette à adaptation (réinitialisation barorécepteurs en 24-48h) $\rightarrow$ Nerveuse.
      \item Longue, soutenue, non adaptative (intègre la masse hydrosodée) $\rightarrow$ Hormonale.
      \end{itemize}
\end{enumerate}
\textbf{Articulation :} Les deux régulations coopèrent en permanence. La nerveuse assure les ajustements « beat-to-beat » et la stabilité minute par minute. L'hormonale assure l'étanchéité du volume circulant et le tonus vasculaire basal, sans lesquels la régulation nerveuse serait inefficace (ex. : en cas d'hypovolémie massive, le baroréflexe maximal ne peut pas maintenir la PA si le volume manque).
\end{methode}

\courssub{Schéma fonctionnel global : une boucle à rétrocontrôle négatif}

L'ensemble forme une \textbf{boucle de régulation à rétrocontrôle négatif} : toute variation de la pression artérielle par rapport à une valeur de consigne (point de fonctionnement) déclenche des réponses effectrices qui s'opposent à cette variation pour la ramener vers la consigne.

\begin{popfigure}
\begin{center}
\begin{tikzpicture}[
  box/.style={draw=popblue, thick, fill=popblueL, rounded corners=4pt, align=center, font=\small, minimum height=1.0cm, minimum width=2.8cm},
  eff/.style={draw=popgreen, thick, fill=popgreenL, rounded corners=4pt, align=center, font=\small, minimum height=1.0cm, minimum width=3.0cm},
  stim/.style={-{Stealth}, thick, popgreen},
  inhib/.style={-{Stealth}, thick, popink, dashed},
  every node/.style={font=\small}]
% Noeuds principaux (ligne haute)
\node[box, align=center] (pa) at (0,0){\textbf{Variable régulée :}\\\textbf{Pression Artérielle}\\(PA = DC $\times$ RPT)};
\node[box, align=center] (baro) at (5.5,0){\textbf{Capteurs :}\\\textbf{Barorécepteurs}\\(Crosse aortique, Sinus carotidiens)\\\textit{Étirement paroi $\propto$ PA}};
\node[box, align=center] (bulbe) at (11,0){\textbf{Centre Intégrateur :}\\\textbf{Bulbe Rachidien}\\(Centres Vasomoteurs)};
% Effecteurs (ligne basse)
\node[eff, align=center] (coeur) at (11,-3.5){\textbf{Effecteurs Rapides (Nerveux) :}\\\textbf{C\oe ur \& Artérioles}\\\textit{Vague (X) : FC$\downarrow$, Vasodilat.}\\\textit{Sympathique : FC$\uparrow$, Force$\uparrow$, Vasoconstr.}};
\node[eff, align=center] (horm) at (5.5,-3.5){\textbf{Effecteurs Lents (Hormonaux) :}\\\textbf{Reins \& Surrénales}\\\textit{Adrénaline / Noradrénaline}\\\textit{Rénine $\rightarrow$ Angio II $\rightarrow$ Aldo}\\\textit{ADH (Vasopressine)}};
% Flèches afférentes
\draw[stim] (pa) -- node[above, font=\footnotesize, align=center]{Variation PA $\rightarrow$\\Étirement paroi} (baro);
\draw[stim] (baro) -- node[above, font=\footnotesize, align=center]{Influx afférents (IX, X)\\\textit{Fréquence $\propto$ PA}} (bulbe);
% Flèches efférentes rapides
\draw[stim] (bulbe) -- node[right, font=\footnotesize, align=left]{Voie efférente nerveuse\\(Secondes)} (coeur);
% Flèches efférentes lentes (depuis bulbe/hypothalamus)
\draw[stim] (bulbe) -- node[left, font=\footnotesize, align=center]{Voie efférente hormonale\\(Minutes -- Heures)} (horm);
% Boucle de rétrocontrôle (retour vers PA)
\draw[stim, popblue] (coeur.south) .. controls (11,-5.5) and (0,-5.5) .. (pa.south) 
    node[pos=0.5, below, font=\footnotesize, align=center, popblue]{\textbf{Rétrocontrôle Négatif :}\\\textit{Correction PA $\rightarrow$ Consigne}};
\draw[stim, poppurple] (horm.south) .. controls (5.5,-5.0) and (2.5,-5.2) .. (1.5,-5.5);
% Signes +/-
\node[font=\footnotesize\bfseries, popgreen] at (9.5,-1.8) {$+$ si PA $\downarrow$};
\node[font=\footnotesize\bfseries, popink] at (12.5,-1.8) {$-$ si PA $\uparrow$};
\end{tikzpicture}
\end{center}
\legende{Schéma fonctionnel synthétique de la régulation de la pression artérielle. La variable régulée (PA) est détectée par les barorécepteurs. L'erreur (écart à la consigne) est traitée par le bulbe rachidien. La réponse nerveuse rapide (secondes) corrige via le c\oe ur et le calibre vasculaire. La réponse hormonale lente (minutes-heures) corrige via le volume sanguin (rénine-angiotensine-aldostérone, ADH) et le tonus vasculaire (adrénaline). Le retour à la consigne caractérise le rétrocontrôle négatif.}
\end{popfigure}

\begin{savaistu}[Le malaise du lever brutal : l'hypotension orthostatique]
Quand tu te lèves brusquement après être resté longtemps allongé ou accroupi, \SI{500}{ml} à \SI{800}{ml} de sang s'accumulent momentanément dans les veines des jambes (capacité veineuse importante, pesanteur) : le retour veineux diminue brutalement, le débit cardiaque chute, la pression artérielle s'effondre au niveau du cœur et de la tête. Le cerveau est transitoirement moins bien irrigué (ischémie cérébrale brève) — d'où une sensation de vertige, des « étoiles » devant les yeux (scotomes), parfois un malaise syncopal (perte de connaissance brève). C'est l'\textbf{hypotension orthostatique} (ou hypotension de décubitus).
En une à deux secondes, les barorécepteurs des sinus carotidiens (au niveau du cou, au niveau du cœur) et de la crosse aortique, moins étirés par la chute de pression, voient leur décharge afférente diminuer. Le baroréflexe se déclenche : levée du frein vagal + activation sympathique massive $\rightarrow$ tachycardie (\textasciitilde +20 à 30 bpm) et vasoconstriction artériolaire et veineuse (remobilisation du sang veineux) $\rightarrow$ la pression remonte et le vertige disparaît.
Ce malaise banal illustre concrètement le \textbf{temps de réponse de la régulation nerveuse} : très court (\textasciitilde 1-3 s), mais pas nul ! Chez les personnes âgées (artères rigides, barorécepteurs moins sensibles), les diabétiques (neuropathie autonome), ou sous certains traitements (bêta-bloquants, antihypertenseurs, diurétiques, antidépresseurs), ce réflexe est ralenti ou atténué, ce qui explique la fréquence des chutes graves ; il est vivement conseillé de se lever \textbf{progressivement} (s'asseoir 30 s, puis se lever).
\end{savaistu}

\begin{exoresolu}[Analyse d'une expérience classique sur le baroréflexe]
\textbf{Énoncé.} Chez un chien anesthésié, on enregistre la pression artérielle moyenne (PAM). On obtient les résultats suivants :

\begin{center}
\begin{tabular}{|l|c|}
\hline
\rowcolor{popblueL} \textbf{Condition expérimentale} & \textbf{PAM (mmHg)} \\
\hline
Témoin (animal intact, au repos) & 110 \\
\hline
Compression des artères carotides \textbf{en aval} des sinus carotidiens & 150 \\
\hline
Après section bilatérale des nerfs glossopharyngiens (nerfs des sinus), même compression & 112 \\
\hline
\end{tabular}
\end{center}

\begin{enumerate}
\item Interpréter les résultats de la compression des carotides chez l'animal intact.
\item Que montre la comparaison avec l'animal dont les nerfs des sinus ont été sectionnés ?
\end{enumerate}

\tcblower
\textbf{Solution détaillée.}

\textbf{1. Interprétation de la compression chez l'animal intact.}
La compression des artères carotides \textbf{en aval} (vers le cerveau) des sinus carotidiens crée un obstacle à l'écoulement sanguin. La pression artérielle \textbf{en amont} de la compression (donc \textbf{dans les sinus carotidiens}) \textbf{diminue} (chute de charge). Les parois des sinus sont \textbf{moins étirées}. Les barorécepteurs, moins stimulés, diminuent leur fréquence de décharge (influx afférents vers le bulbe \textdownarrow).
Le centre vasomoteur bulbaire « perçoit » une hypotension (erreur négative). Il déclenche la réponse effectrice orthosympathique (tachycardie, vasoconstriction) et inhibe le vagal. La pression artérielle systémique \textbf{s'élève} (de \SI{110}{mmHg} à \SI{150}{mmHg}) pour tenter de rétablir la perfusion cérébrale. L'organisme compense une hypotension \emph{artificiellement créée localement} au niveau des capteurs.

\textbf{2. Interprétation de la section des nerfs glossopharyngiens.}
Après section des nerfs afférents (IX) issus des sinus carotidiens, la même compression \textbf{ne provoque plus aucune variation significative} de la PAM (\SI{112}{mmHg} $\approx$ témoin \SI{110}{mmHg}). L'information sur la diminution de l'étirement des parois des sinus ne parvient plus au bulbe rachidien : la voie afférente de l'arc réflexe est \textbf{interrompue}. Le baroréflexe carotidien est aboli.
Cela démontre que :
\begin{itemize}
\item Les sinus carotidiens (et leurs nerfs afférents IX) constituent le \textbf{récepteur} indispensable de ce réflexe de régulation (les barorécepteurs aortiques seuls ne suffisent pas à compenser dans ce protocole, ou leur seuil est plus haut).
\item L'arc réflexe complet est : \textbf{Récepteur} (barorécepteurs sinus/crosse) $\rightarrow$ \textbf{Voie afférente} (IX, X) $\rightarrow$ \textbf{Centre} (Bulbe rachidien) $\rightarrow$ \textbf{Voies efférentes} (Vague X / Sympathique) $\rightarrow$ \textbf{Effecteurs} (C\oe ur, Vaisseaux).
\end{itemize}
\end{exoresolu}

\begin{aretenir}[L'essentiel de la section — Régulation de la PA]
\begin{itemize}
\item La pression artérielle est maintenue stable par une \textbf{boucle à rétrocontrôle négatif}.
\item \textbf{Détection :} \textbf{Barorécepteurs} (mécanorécepteurs d'étirement) situés dans la \textbf{crosse aortique} (nerf X) et les \textbf{sinus carotidiens} (nerf IX).
\item \textbf{Intégration :} \textbf{Bulbe rachidien} (centres vasomoteurs : cardio-inhibiteur, cardio-accélérateur, vasomoteur). \textit{Pas l'hypothalamus.}
\item \textbf{Régulation nerveuse rapide (Baroréflexe, secondes) :}
      \begin{itemize}
      \item PA $\uparrow$ $\rightarrow$ Afférences $\uparrow$ $\rightarrow$ Vague $\uparrow$ (Bradycardie) + Sympathique $\downarrow$ (Vasodilatation) $\rightarrow$ PA $\downarrow$.
      \item PA $\downarrow$ $\rightarrow$ Afférences $\downarrow$ $\rightarrow$ Vague $\downarrow$ + Sympathique $\uparrow$ (Tachycardie, Vasoconstriction) $\rightarrow$ PA $\uparrow$.
      \item Phénomène d'adaptation (réinitialisation) en 24-48h.
      \end{itemize}
\item \textbf{Régulation hormonale lente (Minutes à Jours) :}
      \begin{itemize}
      \item \textbf{Adrénaline / Noradrénaline} (Médullosurrénale) : relais hormonal du sympathique (FC$\uparrow$, Vasoconstriction).
      \item \textbf{SRAA} (Rein $\rightarrow$ Rénine $\rightarrow$ Angio II $\rightarrow$ Aldo) : Vasoconstriction + Réabsorption \ce{Na+}/\ce{H2O} $\rightarrow$ Volémie $\uparrow$ $\rightarrow$ PA $\uparrow$.
      \item \textbf{ADH / Vasopressine} (Hypothalamus/Post-hypophyse) : Réabsorption \ce{H2O} rénale $\rightarrow$ Volémie $\uparrow$ ; Vasoconstriction (dose forte).
      \end{itemize}
\item \textbf{Coopération :} Nerveuse = stabilité beat-to-beat ; Hormonale = maintien volume/tonus basal. Défaillance $\rightarrow$ Pathologies (HTA, Hypotension).
\end{itemize}
\end{aretenir}

Ces mécanismes de régulation sont d'une efficacité remarquable chez l'individu sain. Mais que se passe-t-il lorsqu'ils sont défaillants, débordés de façon chronique, ou que la « consigne » elle-même dérive vers des valeurs trop élevées ? C'est ce que nous étudierons dans la dernière section, consacrée à l'\textbf{hypertension artérielle (HTA)} et à l'\textbf{hypotension artérielle}, véritables problèmes de santé publique au Cameroun, et aux moyens de lutte, de prévention et de sensibilisation.

\coursec{Hypertension et hypotension artérielles : lutte et sensibilisation}

Dans la section précédente, nous avons vu que la pression artérielle est maintenue autour d'une valeur de référence (environ \SI{12}{cmHg}/\SI{8}{cmHg}) grâce à une régulation nerveuse (barorécepteurs, centre vasomoteur bulbaire) et hormonale (adrénaline, système rénine-angiotensine-aldostérone). Mais que se passe-t-il lorsque cette régulation est durablement dépassée, vers le haut ou vers le bas ? C'est là tout l'enjeu de cette dernière section : comprendre deux affections très répandues au Cameroun — l'\cle{hypertension artérielle} et l'\cle{hypotension artérielle} — et surtout apprendre à les prévenir et à sensibiliser notre entourage.

\courssub{L'hypertension artérielle (HTA) : définition et gravité}

\textbf{Une observation pour commencer.} Lors d'une journée de dépistage organisée dans un quartier de Yaoundé, les infirmiers mesurent la tension des passants. M. Ndi, 52 ans, commerçant en pleine forme, affiche \SI{16}{cmHg}/\SI{10}{cmHg}. Surpris, il déclare : « Je n'ai jamais rien senti, je ne suis pas malade ! » Trois ans plus tard, il est victime d'un accident vasculaire cérébral. Cette histoire, malheureusement fréquente, illustre le caractère insidieux de l'hypertension artérielle.

\begin{definition}[Hypertension artérielle]
On parle d'\cle{hypertension artérielle} (HTA) lorsque la pression artérielle est \textbf{élevée de façon permanente}, c'est-à-dire supérieure ou égale à \textbf{\SI{14}{cmHg}/\SI{9}{cmHg}} (\SI{14}{cmHg} pour la pression systolique, \SI{9}{cmHg} pour la pression diastolique), constatée à plusieurs reprises au repos.
\end{definition}

\begin{attention}[Le « tueur silencieux »]
Erreur fréquente : croire que l'HTA se reconnaît toujours à des symptômes (maux de tête, vertiges...). \textbf{C'est faux !} Dans la majorité des cas, l'hypertension est \textbf{totalement asymptomatique} pendant des années : on peut être hypertendu sans rien ressentir. Pendant ce temps, la pression élevée use silencieusement les parois des artères et les organes. C'est pourquoi on l'appelle le \emph{« tueur silencieux »} et pourquoi le \textbf{dépistage régulier} (mesure de la tension au moins une fois par an chez l'adulte) est la seule façon fiable de la détecter.
\end{attention}

\courssub{Causes et facteurs de risque de l'HTA}

Dans environ $90\,\%$ des cas, l'HTA est dite \emph{essentielle} : elle n'a pas de cause unique, mais résulte de la combinaison de plusieurs \cle{facteurs de risque}. On les classe en deux catégories.

\begin{itemize}
\item \textbf{Facteurs non modifiables} : l'\textbf{hérédité} (des parents hypertendus augmentent le risque), l'\textbf{âge} (les artères perdent leur élasticité avec le temps) et le sexe.
\item \textbf{Facteurs modifiables} (sur lesquels on peut agir !) : l'\textbf{excès de sel} alimentaire, l'\textbf{obésité}, la \textbf{sédentarité} (manque d'activité physique), le \textbf{stress} chronique, la consommation d'\textbf{alcool} et de \textbf{tabac}.
\end{itemize}

Dans environ $10\,\%$ des cas, l'HTA est \emph{secondaire} : elle est la conséquence d'une maladie identifiable, le plus souvent une \textbf{maladie rénale} (le rein, acteur clé de la régulation via la rénine, perturbe la pression lorsqu'il est malade).

\begin{popfigure}
\begin{center}
\begin{tikzpicture}[font=\small]
% Noeud central
\node[draw=popdark, fill=popblueL, rounded corners, thick, minimum width=3.6cm, minimum height=1cm, align=center] (hta) {\textbf{Hypertension}\\ \textbf{artérielle}};
% Facteurs non modifiables (gauche)
\node[draw=popmuted, fill=popcream, rounded corners, thick, align=center, minimum width=3.4cm, above left=1.2cm and 1.6cm of hta] (nm1) {\textbf{Hérédité}};
\node[draw=popmuted, fill=popcream, rounded corners, thick, align=center, minimum width=3.4cm, below=0.3cm of nm1] (nm2) {\textbf{Âge}};
\node[draw=popmuted, fill=popcream, rounded corners, thick, align=center, minimum width=3.4cm, below=0.3cm of nm2] (nm3) {\textbf{Maladies rénales}};
\node[above=0.15cm of nm1, font=\small\bfseries, text=popmuted] {Facteurs NON modifiables};
% Facteurs modifiables (droite)
\node[draw=poporange, fill=poporangeL, rounded corners, thick, align=center, minimum width=3.4cm, above right=1.2cm and 1.6cm of hta] (m1) {\textbf{Excès de sel}};
\node[draw=poporange, fill=poporangeL, rounded corners, thick, align=center, minimum width=3.4cm, below=0.3cm of m1] (m2) {\textbf{Obésité, sédentarité}};
\node[draw=poporange, fill=poporangeL, rounded corners, thick, align=center, minimum width=3.4cm, below=0.3cm of m2] (m3) {\textbf{Stress, alcool, tabac}};
\node[above=0.15cm of m1, font=\small\bfseries, text=poporange] {Facteurs modifiables};
% Flèches
\draw[-{Stealth}, thick, popmuted] (nm1.east) -- (hta.west);
\draw[-{Stealth}, thick, popmuted] (nm2.east) -- (hta.west);
\draw[-{Stealth}, thick, popmuted] (nm3.east) -- (hta.west);
\draw[-{Stealth}, thick, poporange] (m1.west) -- (hta.east);
\draw[-{Stealth}, thick, poporange] (m2.west) -- (hta.east);
\draw[-{Stealth}, thick, poporange] (m3.west) -- (hta.east);
\end{tikzpicture}
\end{center}
\legende{Facteurs de risque de l'hypertension artérielle. À gauche, les facteurs non modifiables (hérédité, âge, maladies rénales) ; à droite, les facteurs modifiables sur lesquels porte la prévention (sel, obésité, sédentarité, stress, alcool, tabac).}
\end{popfigure}

\courssub{Symptômes et complications de l'HTA}

\textbf{Symptômes.} Lorsqu'ils existent, ils sont discrets et non spécifiques : \textbf{maux de tête} (surtout matinaux, à l'arrière du crâne), \textbf{vertiges}, \textbf{bourdonnements d'oreilles}, \textbf{saignements de nez}. Mais rappelle-toi : le plus souvent, il n'y a \emph{aucun} symptôme.

\textbf{Complications.} Une pression élevée en permanence fragilise les parois artérielles (athérosclérose) et épuise les organes. Les organes cibles de l'HTA sont au nombre de quatre :

\begin{popfigure}
\begin{center}
\begin{tikzpicture}[font=\small]
% HTA centrale
\node[draw=popdark, fill=poppinkL, rounded corners, thick, minimum width=3.2cm, minimum height=0.9cm, align=center] (hta) {\textbf{HTA chronique}};
% Cerveau
\node[draw=popblue, fill=popblueL, ellipse, thick, minimum width=2.6cm, minimum height=1.1cm, align=center, above left=1.6cm and 2.2cm of hta] (cerveau) {\textbf{Cerveau}};
\node[below=0.1cm of cerveau, align=center, text=popdark] (avc) {AVC (hémorragie ou\\ infarctus cérébral)};
% Coeur
\node[draw=poppink, fill=poppinkL, ellipse, thick, minimum width=2.6cm, minimum height=1.1cm, align=center, above right=1.6cm and 2.2cm of hta] (coeur) {\textbf{C\oe ur}};
\node[below=0.1cm of coeur, align=center, text=popdark] (inf) {Infarctus du myocarde,\\ insuffisance cardiaque};
% Reins
\node[draw=popgreen, fill=popgreenL, ellipse, thick, minimum width=2.6cm, minimum height=1.1cm, align=center, below left=1.6cm and 2.2cm of hta] (reins) {\textbf{Reins}};
\node[below=0.1cm of reins, align=center, text=popdark] (ir) {Insuffisance rénale\\ chronique};
% Yeux
\node[draw=popgold, fill=popgoldL, ellipse, thick, minimum width=2.6cm, minimum height=1.1cm, align=center, below right=1.6cm and 2.2cm of hta] (yeux) {\textbf{Yeux (rétine)}};
\node[below=0.1cm of yeux, align=center, text=popdark] (ret) {Atteinte de la rétine,\\ perte de vision};
% Flèches
\draw[-{Stealth}, very thick, popdark] (hta) -- (cerveau);
\draw[-{Stealth}, very thick, popdark] (hta) -- (coeur);
\draw[-{Stealth}, very thick, popdark] (hta) -- (reins);
\draw[-{Stealth}, very thick, popdark] (hta) -- (yeux);
\end{tikzpicture}
\end{center}
\legende{Les organes cibles de l'hypertension artérielle et leurs complications : le cerveau (accident vasculaire cérébral), le c\oe ur (infarctus du myocarde, insuffisance cardiaque), les reins (insuffisance rénale) et les yeux (atteinte rétinienne).}
\end{popfigure}

\begin{itemize}
\item \textbf{Cerveau} : l'\cle{accident vasculaire cérébral} (AVC), par rupture ou obstruction d'une artère cérébrale, peut entraîner paralysie, troubles du langage, voire la mort.
\item \textbf{C\oe ur} : l'\cle{infarctus du myocarde} (nécrose d'une partie du muscle cardiaque) et l'\cle{insuffisance cardiaque} (le c\oe ur, fatigué de pomper contre une pression trop élevée, ne remplit plus son rôle).
\item \textbf{Reins} : l'\cle{insuffisance rénale} chronique, pouvant conduire à la dialyse.
\item \textbf{Yeux} : l'atteinte des vaisseaux de la \textbf{rétine}, pouvant aboutir à une perte de vision.
\end{itemize}

\begin{savaistu}[L'HTA au Cameroun]
L'hypertension artérielle touche environ \textbf{un adulte sur trois au Cameroun} et constitue la \textbf{première cause d'AVC} dans notre pays. Beaucoup de Camerounais ignorent leur état : on estime qu'une grande proportion des hypertendus ne sont ni dépistés ni traités. Les services de cardiologie des hôpitaux de Yaoundé et de Douala reçoivent chaque année des milliers de patients au stade de complications, alors qu'un dépistage précoce aurait suffi à les éviter.
\end{savaistu}

\courssub{L'hypotension artérielle}

\begin{definition}[Hypotension artérielle]
On parle d'\cle{hypotension artérielle} lorsque la pression artérielle est \textbf{abaissée de façon durable} en dessous de \textbf{\SI{9}{cmHg}/\SI{6}{cmHg}}.
\end{definition}

\textbf{Causes.} Les plus fréquentes sont la \textbf{déshydratation} (fortes chaleurs, diarrhées, sudation abondante), le \textbf{jeûne prolongé}, une \textbf{hémorragie} (perte importante de sang, par exemple lors d'un accident), et la prise de \textbf{certains médicaments}.

\textbf{Symptômes.} Contrairement à l'HTA, l'hypotension se manifeste : \textbf{fatigue}, \textbf{vertiges} (surtout au passage brusque de la position couchée à la position debout : hypotension orthostatique), \textbf{malaises}, voire \textbf{évanouissement} (perte de connaissance brève, due à une mauvaise irrigation du cerveau).

\textbf{Gravité.} L'hypotension est \textbf{généralement bénigne} et ne nécessite qu'un repos, une réhydratation et une alimentation suffisante. Elle devient grave dans un cas particulier : le \textbf{choc} (choc hémorragique, septique ou allergique), où la chute brutale de pression prive les organes vitaux d'oxygène — c'est alors une \textbf{urgence vitale}.

\begin{exoresolu}[Interpréter des mesures de tension]
\textbf{Énoncé.} Lors d'une campagne de dépistage à Douala, trois adultes présentent les tensions suivantes, mesurées au repos à trois reprises : A : \SI{11}{cmHg}/\SI{7}{cmHg} ; B : \SI{15}{cmHg}/\SI{10}{cmHg} ; C : \SI{8}{cmHg}/\SI{5}{cmHg}. Pour chaque personne, indique son état tensionnel et la conduite à tenir.
\tcblower
\textbf{Solution.}
\begin{itemize}
\item \textbf{A (\SI{11}{cmHg}/\SI{7}{cmHg})} : tension \textbf{normale} (inférieure à \SI{14}{cmHg}/\SI{9}{cmHg} et supérieure à \SI{9}{cmHg}/\SI{6}{cmHg}). Conduite : aucune inquiétude ; rappeler les règles d'hygiène de vie et un contrôle annuel.
\item \textbf{B (\SI{15}{cmHg}/\SI{10}{cmHg})} : les deux valeurs dépassent les seuils de \SI{14}{cmHg} et \SI{9}{cmHg}, et cela de façon répétée au repos : il s'agit d'une \textbf{hypertension artérielle}. Conduite : consultation médicale pour confirmation et mise en route d'un traitement, réduction du sel, activité physique, surveillance régulière, même en l'absence de tout symptôme.
\item \textbf{C (\SI{8}{cmHg}/\SI{5}{cmHg})} : tension durablement inférieure à \SI{9}{cmHg}/\SI{6}{cmHg} : il s'agit d'une \textbf{hypotension artérielle}. Conduite : rechercher une cause (déshydratation, jeûne, médicament) ; en l'absence de signe de choc, repos, hydratation et alimentation suffisent généralement ; consulter si les malaises persistent.
\end{itemize}
\end{exoresolu}

\courssub{Moyens de lutte contre l'HTA : prévention et traitement}

La lutte contre l'hypertension repose sur deux piliers complémentaires : l'\textbf{hygiène de vie} et le \textbf{traitement médical}.

\begin{enumerate}
\item \textbf{Réduire la consommation de sel} : c'est la mesure la plus efficace. Le sel retient l'eau dans le sang, augmente le volume sanguin et donc la pression.
\item \textbf{Adopter une alimentation riche en fruits et légumes} (bananes plantains, légumes feuilles, avocats...), pauvre en graisses saturées.
\item \textbf{Pratiquer une activité physique régulière} : au moins \SI{30}{min} de marche rapide par jour.
\item \textbf{Lutter contre le stress} et \textbf{arrêter le tabac et l'alcool}, qui endommagent directement les parois artérielles.
\item \textbf{Se faire dépister régulièrement} : mesure de la tension au moins une fois par an, davantage en cas d'hérédité.
\item \textbf{Suivre scrupuleusement le traitement antihypertenseur} prescrit : il se prend \emph{à vie}, tous les jours, même quand on se sent bien. L'arrêt du traitement expose aux complications.
\end{enumerate}

\begin{exemplebox}[L'effet du sel, en chiffres]
L'Organisation mondiale de la santé (OMS) recommande de limiter la consommation de sel à \textbf{moins de \SI{5}{g} par jour} (environ une cuillère à café), alors que de nombreux adultes en consomment \SI{10}{g} à \SI{15}{g} par jour (plats très salés, cubes d'assaisonnement, poisson fumé salé). Des études montrent que ramener la consommation sous ce seuil \textbf{diminue significativement le risque d'HTA} et fait baisser la pression systolique de plusieurs mmHg chez les hypertendus.
\end{exemplebox}

\begin{aretenir}[L'essentiel]
\begin{itemize}
\item HTA : pression $\geq$ \SI{14}{cmHg}/\SI{9}{cmHg} de façon permanente ; maladie \textbf{silencieuse} ; complications : AVC, infarctus, insuffisances cardiaque et rénale, atteinte de la rétine.
\item Hypotension : pression $\leq$ \SI{9}{cmHg}/\SI{6}{cmHg} ; généralement bénigne, grave en cas de choc.
\item Prévention de l'HTA : moins de sel, fruits et légumes, activité physique, pas de tabac ni d'alcool, gestion du stress, dépistage et traitement suivi.
\end{itemize}
\end{aretenir}

\courssub{Élaborer des outils de sensibilisation}

Savoir ne suffit pas : il faut \textbf{transmettre}. En tant que citoyen et élève de Terminale, tu peux participer à la lutte contre l'HTA en concevant des outils de sensibilisation adaptés à ta communauté : \textbf{affiches} placardées dans les marchés et les centres de santé, \textbf{sketchs} joués lors des réunions de quartier, \textbf{spots radio en langues locales} (ewondo, douala, fulfuldé...), et participation aux \textbf{journées de dépistage communautaire} organisées dans les quartiers de Yaoundé et de Douala.

\begin{methode}[Construire une affiche de sensibilisation efficace]
\begin{enumerate}
\item \textbf{Un message clair et unique} : une seule idée forte, par exemple « Connais ta tension, sauve ta vie ».
\item \textbf{Un chiffre repère} : afficher le seuil \textbf{\SI{14}{cmHg}/\SI{9}{cmHg}} en grand, pour que chacun sache à partir de quand consulter.
\item \textbf{Trois gestes préventifs} illustrés par des dessins simples : moins de sel, bouger chaque jour, pas de tabac.
\item \textbf{Un appel au dépistage} : indiquer où et quand faire mesurer sa tension (centre de santé le plus proche, journée de dépistage du quartier).
\end{enumerate}
Une bonne affiche se lit en moins de dix secondes et doit être compréhensible par tous, y compris les personnes non scolarisées : privilégie les images et les phrases courtes.
\end{methode}

\begin{exercice}[À toi de jouer]
Ton établissement organise une semaine de la santé. Conçois le texte d'un spot radio de \SI{30}{s} en langue locale de ton choix (ou en français simple) sensibilisant à l'HTA. Ton spot devra contenir : la définition du seuil \SI{14}{cmHg}/\SI{9}{cmHg}, le caractère silencieux de la maladie, deux complications redoutées et un appel au dépistage gratuit au centre de santé de ton quartier.
\end{exercice}

\begin{corrige}[Exercice]
Exemple de spot : « Attention, message de santé ! L'hypertension, c'est quand ta tension dépasse \SI{14}{cmHg}/\SI{9}{cmHg} en permanence. Elle ne fait pas mal, elle ne prévient pas, mais elle détruit en silence le c\oe ur, le cerveau et les reins : infarctus, AVC ! Toi qui m'écoutes, n'attends pas les vertiges. Mange moins de sel, bouge chaque jour, et viens mesurer ta tension gratuitement samedi au centre de santé du quartier. Connais ta tension, sauve ta vie ! » Ce spot respecte les quatre consignes : seuil chiffré, silence de la maladie, complications citées, appel au dépistage avec lieu et date.
\end{corrige}

Ainsi se clôt notre étude des problèmes liés à la régulation de la glycémie et de la pression artérielle. Tu disposes désormais des savoirs et des gestes pour protéger ta santé et celle de ta communauté : mesurer, comprendre, prévenir et sensibiliser. C'est là l'essence même de l'éducation à la santé.

\newpage

\coursec{Activité d'intégration}

\coursec{Activité d'intégration : Analyse de dossiers cliniques — Glycémie et Pression Artérielle}

\courssub{Objectifs de l'activité}
\begin{itemize}
    \item Mobiliser l'ensemble des savoirs et savoir-faire de la leçon (mesure et régulation de la glycémie et de la pression artérielle, pathologies associées) pour analyser des dossiers médicaux réalistes.
    \item Poser un diagnostic étiologique argumenté à partir de données biologiques (glycémie, insulinémie, HGPO) et cliniques (tension artérielle, anamnèse).
    \item Réaliser des schémas fonctionnels illustrant la défaillance des boucles de régulation (insuline/glucagon ; système nerveux autonome/hormonal rénal).
    \item Identifier les facteurs de risque modifiables et non modifiables, ainsi que les complications aiguës et chroniques.
    \item Concevoir un outil de sensibilisation (affiche) adapté au contexte socioculturel camerounais (marché de quartier) pour promouvoir l'hygiène de vie.
\end{itemize}

\courssub{Situation}
\begin{center}
    \textbf{Centre de Santé Intégré (CSI) « La Paix » — Arrondissement de Yaoundé III, Cameroun}
\end{center}
\vspace{0.5em}
Vous êtes une équipe de médecins généralistes et d'infirmiers en stage de fin d'année au CSI « La Paix ». Ce mardi matin, la salle d'attente est pleine. Le chef de service vous confie deux dossiers complexes nécessitant une analyse collégiale avant la consultation de l'après-midi. Votre mission : établir un diagnostic précis, expliquer les mécanismes physiopathologiques aux patients et proposer un plan de prise en charge et de prévention communautaire.

\vspace{1em}
\noindent\textbf{Dossier n° 1 : Mme A., 48 ans, commerçante au marché Mfoundi.}
\begin{itemize}
    \item \textbf{Motif de consultation :} Fatigue intense depuis 3 mois, soif excessive (polydipsie), envies fréquentes d'uriner (polyurie), perte de poids inexpliquée malgré un appétit conservé.
    \item \textbf{Antécédents :} Pas de diabète connu dans la famille au premier degré. Mode de vie sédentaire (activité assise toute la journée). Alimentation riche en glucides raffinés (pain, pâtes, beignets, boissons sucrées) et en graisses (fritures). IMC calculé à \SI{32,5}{kg.m^{-2}} (Obésité classe I). Tour de taille : \SI{98}{cm}.
    \item \textbf{Examens complémentaires demandés la veille (résultats du jour) :}
    \begin{itemize}
        \item Glycémie à jeun (GàJ) : \SI{1,60}{g.L^{-1}} (norme : 0,70 -- 1,10 g/L).
        \item Hyperglycémie Provoquée Orale (HGPO) : 75 g de glucose per os.
        \begin{center}
            \begin{tabular}{|c|c|c|c|c|}\hline
                Temps (min) & 0 (à jeun) & 30 & 60 & 120 \\ \hline
                Glycémie (g/L) & 1,60 & 2,10 & \textbf{2,40} & 2,20 \\ \hline
            \end{tabular}
        \end{center}
        \item Insulinémie à jeun : \SI{18}{\mu UI.mL^{-1}} (norme : 5 -- 25 \mu UI/mL).
        \item HbA1c (Hémoglobine glyquée) : 8,2 \% (norme < 6,5 \%).
        \item Profil lipidique : Cholestérol total \SI{2,45}{g.L^{-1}}, LDL-c \SI{1,70}{g.L^{-1}}, Triglycérides \SI{2,10}{g.L^{-1}}, HDL-c \SI{0,35}{g.L^{-1}}.
    \end{itemize}
\end{itemize}

\vspace{1em}
\noindent\textbf{Dossier n° 2 : M. B., 55 ans, chauffeur de taxi-brousse (ligne Yaoundé-Douala).}
\begin{itemize}
    \item \textbf{Motif de consultation :} Céphalées matinales récurrentes, vertiges, bourdonnements d'oreilles (acouphènes) depuis 6 mois. Ne consulte que sur insistance de sa femme.
    \item \textbf{Antécédents :} Père décédé d'un AVC à 60 ans. Tabagisme : 15 cigarettes/jour depuis 30 ans. Alcool : consommation quotidienne de \textit{vin de palme} (environ 1 L/jour) et bière locale. Alimentation : plats très salés (sauces « pistache », « ndolé » avec beaucoup de sel de cuisine et cubes de bouillon), faible consommation de fruits/légumes. Activité physique : nulle (conduite assise > 12 h/jour). IMC : \SI{28,8}{kg.m^{-2}} (Surpoids).
    \item \textbf{Examens cliniques (trois mesures espacées de 2 min, au repos, bras au niveau du cœur, brassard adapté) sur 3 jours consécutifs :}
    \begin{center}
        \begin{tabular}{|c|c|c|c|}\hline
            Jour & PAS (mmHg) & PAD (mmHg) & FC (bpm) \\ \hline
            J1 & 160 & 100 & 88 \\
            J2 & 150 & 95 & 92 \\
            J3 & 170 & 105 & 90 \\ \hline
        \end{tabular}
    \end{center}
    \item \textbf{Bilan biologique rapide :} Glycémie à jeun \SI{1,05}{g.L^{-1}}, Créatininémie \SI{14}{\mg.L^{-1}} (norme 6--12 mg/L), Protéinurie aux bandelettes : + (traces), Ionogramme sanguin : Na\textsuperscript{+} \SI{145}{mmol.L^{-1}}, K\textsuperscript{+} \SI{3,8}{mmol.L^{-1}}.
\end{itemize}

\courssub{Consignes}
\emph{Travail en groupe (4--5 élèves). Durée : 2 h. Restitution orale de 10 min par groupe.}
\begin{enumerate}
    \item \textbf{Analyse du Dossier 1 (Mme A.) :}
    \begin{enumerate}
        \item Interpréter la glycémie à jeun et la courbe d'HGPO au regard des seuils diagnostiques OMS. Calculer l'amplitude de la variation glycémique (pic - à jeun).
        \item Analyser le couple Glycémie/Insulinémie à jeun. Que suggère une insulinémie « normale » face à une hyperglycémie marquée ? En déduire le type de diabète (typologie étiologique).
        \item Expliquer la valeur de l'HbA1c. Quel est l'intérêt de ce marqueur par rapport à la glycémie instantanée ?
        \item Réaliser le \textbf{schéma fonctionnel de la régulation de la glycémie chez Mme A.} en faisant apparaître : le stimulus (hyperglycémie), les capteurs (cellules $\beta$ pancréatiques), le message chimique (insuline), les effecteurs (foie, muscle, tissu adipeux), la réponse (entrée du glucose) et la \textbf{défaillance} (flèche barrée ou symbole $\varnothing$ sur la sensibilité à l'insuline). Légender précisément.
        \item Lister les complications chroniques probables au vu du bilan lipidique et de l'HbA1c (macro et microangiopathies).
    \end{enumerate}
    \item \textbf{Analyse du Dossier 2 (M. B.) :}
    \begin{enumerate}
        \item Valider le diagnostic d'Hypertension Artérielle (HTA) selon les classifications actuelles (seuils PAS/PAD). Calculer la Pression Artérielle Moyenne (PAM) pour le Jour 3. Rappeler la formule.
        \item Identifier les facteurs de risque \textbf{modifiables} et \textbf{non modifiables} présents chez M. B. Hiérarchiser leur impact physiopathologique (volume sanguin, résistances périphériques, rigidité artérielle).
        \item Expliquer le mécanisme physiologique liant l'excès de sel (NaCl) et d'alcool à l'élévation de la PAM (rôle du volume extracellulaire, du système rénine-angiotensine-aldostérone - SRAA, du système nerveux sympathique).
        \item Réaliser le \textbf{schéma fonctionnel de la régulation de la PA chez M. B.} en montrant la suractivation du SRAA et du système sympathique, et l'insuffisance des mécanismes dépresseurs (barorécepteurs, peptides natriurétiques). Indiquer les cibles thérapeutiques (diurétiques, IEC/ARA II, bêta-bloquants, calciques).
        \item Interpréter la créatininémie et la protéinurie : quelle atteinte d'organe « cible » est suspectée ?
    \end{enumerate}
    \item \textbf{Synthèse et Action Communautaire :}
    \begin{enumerate}
        \item Comparer les deux dossiers : points communs (facteurs de risque métaboliques) et différences (régulation défaillante).
        \item Concevoir une \textbf{affiche de sensibilisation} (format A3, portrait) destinée à être affichée au marché Mfoundi et à la gare routière. L'affiche doit comporter : un titre accrocheur, 3 messages clés illustrés (pictogrammes), les chiffres clés locaux (prévalence HTA/Diabète au Cameroun $\approx$ 30\% / 6-8\% adultes), les numéros d'urgence (SAMU 112, Centre Antipoison) et l'adresse du CSI. Utiliser le français courant et une langue locale (ex: Ewondo, Bassa, Fulfulde) pour un slogan.
    \end{enumerate}
\end{enumerate}

\courssub{Ressources à mobiliser}
\begin{itemize}
    \item \textbf{Seuils diagnostiques OMS (Diabète) :} GàJ $\ge$ 1,26 g/L (7,0 mmol/L) OU Glycémie 2h post-HGPO $\ge$ 2,00 g/L (11,1 mmol/L) OU HbA1c $\ge$ 6,5 \%. Prédiabète : GàJ 1,10--1,25 g/L (altération de la glycémie à jeun) ou 2h-HGPO 1,40--1,99 g/L (intolérance au glucose).
    \item \textbf{Seuils HTA (ESC/ISH 2023 / OMS) :} HTA confirmée au cabinet si PAS $\ge$ 140 mmHg et/ou PAD $\ge$ 90 mmHg. Grade 1 : 140-159 / 90-99. Grade 2 : 160-179 / 100-109. Grade 3 : $\ge$ 180 / $\ge$ 110.
    \item \textbf{Formule PAM :} $PAM = PAD + \frac{1}{3}(PAS - PAD)$.
    \item \textbf{Régulation Glycémie :} Boucle de rétroaction négative. Hyperglycémie $\xrightarrow{\text{cellules }\beta}$ Insuline $\uparrow$ $\xrightarrow{\text{Récepteurs tyrosine kinase}}$ Translocation GLUT4 (muscle/adipose) + Glycogénèse/Glycolyse (foie) $\rightarrow$ Glycémie $\downarrow$. Hypoglycémie $\xrightarrow{\text{cellules }\alpha}$ Glucagon $\uparrow$ $\xrightarrow{\text{Récepteurs GPCR}}$ Glycogénolyse/Gluconéogenèse (foie) $\rightarrow$ Glycémie $\uparrow$. \textbf{Insulinorésistance :} Défaut de transduction du signal (récepteur/IRS/PI3K) $\rightarrow$ GLUT4 non translocé.
    \item \textbf{Régulation PA :} $PA = DC \times RPP$ (Débit Cardiaque $\times$ Résistances Périphériques Totales). Barorécepteurs (sinus carotidien, arc aortique) $\rightarrow$ Bulbe (centre vasomoteur) $\rightarrow$ Système Nerveux Orthosympathique (vasoconstriction, FC$\uparrow$, FC$\downarrow$ via parasympathique). SRAA : Rénine $\rightarrow$ Angiotensine I $\xrightarrow{\text{ECA}}$ Angiotensine II (vasoconstriction puissante, soif, Aldostérone $\rightarrow$ Réabsorption Na\textsuperscript{+}/Eau rein). Peptides Natriurétiques (ANP/BNP) : vasodilatation, natriurèse (antagonistes).
    \item \textbf{Complications :} Diabète : Rétinopathie, Néphropathie, Neuropathie, Pied diabétique, AVC, IDM. HTA : AVC (hémorragique/ischémique), Insuffisance Cardiaque, Néphroangiosclérose, Rétinopathie hypertensive.
\end{itemize}

\begin{methode}[Diagnostic biologique du Diabète (OMS)]
\begin{enumerate}
    \item Mesurer la Glycémie à Jeun (GàJ) : $\ge$ 1,26 g/L (7,0 mmol/L) $\rightarrow$ Diabète. Entre 1,10 et 1,25 g/L $\rightarrow$ Prédiabète (AGJ altérée).
    \item Si GàJ équivoque ou pour confirmation : HGPO (75 g glucose). Glycémie à 120 min $\ge$ 2,00 g/L (11,1 mmol/L) $\rightarrow$ Diabète. Entre 1,40 et 1,99 g/L $\rightarrow$ Prédiabète (Intolérance au glucose).
    \item HbA1c $\ge$ 6,5 \% $\rightarrow$ Diabète (reflet équilibre 3 mois). Attention : hémoglobinopathies, anémies hémolytiques faussent le résultat.
    \item \textbf{Typage étiologique :} Couple Glycémie/Insulinémie. Insulinémie basse/indétectable + Hyperglycémie $\rightarrow$ Type 1 (auto-immun). Insulinémie normale/élevée + Hyperglycémie $\rightarrow$ Insulinorésistance (Type 2).
\end{enumerate}
\end{methode}

\begin{methode}[Diagnostic et Grade de l'HTA]
\begin{enumerate}
    \item Mesurer la PA (3 mesures, 2 min d'intervalle, bras au niveau du cœur, brassard adapté) sur 3 consultations espacées (ou automesure validée).
    \item Moyenne des mesures : PAS $\ge$ 140 \textbf{et/ou} PAD $\ge$ 90 mmHg $\rightarrow$ HTA confirmée.
    \item Grader : Grade 1 (140-159/90-99), Grade 2 (160-179/100-109), Grade 3 ($\ge$180/$\ge$110).
    \item Calculer PAM = PAD + 1/3(PAS-PAD). PAM normale $\approx$ 90-100 mmHg.
    \item Rechercher atteintes d'organes cibles (cœur, rein, vaisseaux, œil, cerveau) et facteurs de risque associés.
\end{enumerate}
\end{methode}

\begin{attention}[Pièges fréquents au Bac]
\begin{itemize}
    \item \textbf{Insulinémie « normale » $\neq$ fonction normale :} Une insulinémie dans les normes (ex: 18 µUI/mL) face à une hyperglycémie (1,60 g/L) traduit une \textbf{insulinorésistance} (échec de la sécrétion compensatrice). Ne pas conclure à un diabète de type 1 (insulinopénie).
    \item \textbf{HbA1c et anémie :} Une anémie hémolytique ou une drépanocytose (fréquente au Cameroun) raccourcit la vie du globule rouge $\rightarrow$ HbA1c faussement basse. Demander un dosage de la fructosamine si doute.
    \item \textbf{HTA « blouse blanche » :} Une seule mesure élevée ne suffit pas. Exiger 3 séries de mesures sur 3 jours (ou MAPA/automesure).
    \item \textbf{Grade vs Stade :} Le \textbf{Grade} dépend des chiffres de PA. Le \textbf{Stade} dépend des atteintes d'organes (Stade 1 : pas d'atteinte, Stade 2 : atteinte asymptomatique -- ex: protéinurie, HVG, Stade 3 : maladie cardiovasculaire établie).
    \item \textbf{Unités :} Glycémie en g/L (ou mmol/L), Insulinémie en µUI/mL (ou pmol/L : 1 µUI/mL = 6,945 pmol/L), HbA1c en \%, PA en mmHg, Créatininémie en mg/L (ou µmol/L : 1 mg/L = 88,4 µmol/L).
\end{itemize}
\end{attention}

\courssub{Démarche et corrigé commenté}
\begin{exoresolu}[Analyse intégrée des deux dossiers cliniques]
\textbf{1. DOSSIER 1 : Mme A. — Diabète de Type 2 (DT2) avec insulinorésistance majeure.}

\textbf{1.1. Interprétation des glycémies.}
\begin{itemize}
    \item GàJ = \SI{1,60}{g.L^{-1}} $>$ \SI{1,26}{g.L^{-1}} $\Rightarrow$ \textbf{Diabète} (seuil OMS franchi).
    \item HGPO : Pic à \SI{2,40}{g.L^{-1}} à 60 min. Valeur à 120 min = \SI{2,20}{g.L^{-1}} $>$ \SI{2,00}{g.L^{-1}} $\Rightarrow$ \textbf{Intolérance au glucose sévère / Diabète confirmé}. Absence de retour à la normale (non résorption du glucose).
    \item Amplitude = Pic - À jeun = 2,40 - 1,60 = \SI{0,80}{g.L^{-1}}. Amplitude physiologique attendue $<$ 1,10 g/L avec retour à la normale à 120 min. Ici, défaut de clairance du glucose.
\end{itemize}

\textbf{1.2. Couple Glycémie/Insulinémie \& Typage.}
\begin{itemize}
    \item Insulinémie à jeun = \SI{18}{\mu UI.mL^{-1}} (norme 5-25). Elle est \textbf{normale en valeur absolue} mais \textbf{inappropriément basse} relativement à l'hyperglycémie (1,60 g/L).
    \item Chez un sujet sain, une glycémie à 1,60 g/L devrait stimuler une sécrétion massive d'insuline ($>$ 50-100 $\mu$UI/mL).
    \item \textbf{Conclusion :} Hyperglycémie + Insulinémie normale/faible pour le niveau de glycémie = \textbf{Insulinorésistance périphérique} (foie, muscle, adipeux) + \textbf{Défaillance relative de la sécrétion insulinique} (épuisement progressif des cellules $\beta$).
    \item \textbf{Diagnostic étiologique : Diabète de Type 2 (DT2).} Non insulino-dépendant au départ. Contexte : Obésité abdominale (adipocytes hypertrophiés sécrétant des adipokines pro-inflammatoires : TNF-$\alpha$, IL-6, résistine $\rightarrow$ inhibition signalisation insuline), sédentarité, âge > 40 ans.
\end{itemize}

\textbf{1.3. HbA1c = 8,2 \%.}
\begin{itemize}
    \item Reflète la glycémie moyenne sur la durée de vie de l'hémoglobine ($\approx$ 120 jours).
    \item Correspond à une glycémie moyenne estimée (eAG) $\approx$ \SI{1,90}{g.L^{-1}} (formule ADAG : $eAG(\text{mg/dL}) = 28,7 \times HbA1c - 46,7$ $\rightarrow$ 189 mg/dL).
    \item Confirme un \textbf{déséquilibre glycémique chronique} (3 mois), excluant un pic de stress aigu. Objectif thérapeutique : $<$ 7\% (individualisé selon âge/comorbidités).
\end{itemize}

\textbf{1.4. Schéma fonctionnel de la régulation défaillante (DT2).}
\begin{popfigure}
\begin{tikzpicture}[
    node distance=1.2cm and 1.5cm,
    box/.style={rectangle, draw=popdark, thick, fill=popcream, rounded corners, minimum height=1.2cm, minimum width=2.5cm, align=center},
    boxorange/.style={box, fill=poporange!20, draw=poporange},
    boxgreen/.style={box, fill=popgreen!20, draw=popgreen},
    arrow/.style={-Stealth, thick, popdark},
    orangearr/.style={-Stealth, thick, poporange},
    cross/.style={poporange, thick, scale=1.5}
]
% Nodes
\node[box, align=center] (stim){Stimulus :\\ Hyperglycémie\\ (1,60 g/L)};
\node[box, right=of stim, align=center] (capteur){Capteurs :\\ Cellules $\beta$ pancréatiques\\ (Îlots de Langerhans)};
\node[box, right=of capteur, align=center] (msg){Message :\\ Insuline libérée\\ (Quantité normale mais\\ inefficace)};
\node[box, below=of msg, xshift=1.5cm, align=center] (eff){Effecteurs cibles :\\ \textbf{Foie}, \textbf{Muscle}, \textbf{Adipose}};
\node[box, right=of eff, align=center] (rep){Réponse attendue :\\ $\downarrow$ Glycogénolyse / $\uparrow$ Glycogénèse (Foie)\\
$\uparrow$ Translocation GLUT4 / $\uparrow$ Captage Glc (Muscle/Adipose)\\
$\downarrow$ Lipolyse / $\uparrow$ Lipogenèse (Adipose)};
\node[boxorange, below=of rep, align=center] (def){DÉFAILLANCE :\\ \textbf{Insulinorésistance}\\ (Récepteur Insuline / IRS-1 / PI3K altérés)\\
$\rightarrow$ GLUT4 non translocé\\
$\rightarrow$ Gluconéogenèse hépatique non freinée\\
$\rightarrow$ Lipolyse non inhibée};
\node[boxgreen, right=of def, align=center] (cons){Conséquence :\\ \textbf{Hyperglycémie entretenue}\\
+ Dyslipidémie (HTG, bas HDL)\\
+ HbA1c $\uparrow$};

% Arrows
\draw[arrow] (stim) -- node[above, font=\small] {Détection} (capteur);
\draw[arrow] (capteur) -- node[above, font=\small] {Sécrétion} (msg);
\draw[arrow] (msg) -- node[right, font=\small] {Transport sanguin} (eff);
\draw[arrow] (eff) -- node[above, font=\small] {Action normale} (rep);
\draw[orangearr] (rep) -- node[right, font=\small] {ÉCHEC} (def);
\draw[orangearr] (def) -- (cons);

% Inhibition feedback (broken)
\draw[arrow, dashed, popmuted] (cons.west) to[out=180, in=-90] node[left, font=\small, popmuted] {Rétroaction (-) INEFFICACE} (stim.north west);

% Cross on insulin arrow
\node[cross] at ($(msg.east)!0.5!(eff.west) + (0,0.4)$) {\textbf{$\times$}};
\node[font=\small, poporange, anchor=south] at ($(msg.east)!0.5!(eff.west) + (0,0.6)$) {Résistance};
\end{tikzpicture}
\legende{Schéma fonctionnel de la régulation de la glycémie chez Mme A. La flèche orange barrée et le symbole $\times$ indiquent le blocage de l'action de l'insuline au niveau de ses récepteurs cellulaires (insulinorésistance), empêchant la baisse de la glycémie malgré une sécrétion pancréatique conservée.}
\end{popfigure}

\textbf{1.5. Complications probables.}
\begin{itemize}
    \item \textbf{Macroangiopathies :} Risque cardiovasculaire très élevé (Score FRAMINGHAM/SCORE2 majoré par diabète). Dyslipidémie athérogène (LDL-c \SI{1,70}{g.L^{-1}} > cible < 1,00 g/L chez diabétique ; HDL-c bas \SI{0,35}{g.L^{-1}} ; TG élevés \SI{2,10}{g.L^{-1}}). Risque d'IDM, AVC ischémique, artériopathie des membres inférieurs.
    \item \textbf{Microangiopathies :} Néphropathie (recherche microalbuminurie urgente), Rétinopathie (fond d'œil urgent), Neuropathie sensitive/motrice (risque pied diabétique).
\end{itemize}

\vspace{1em}
\textbf{2. DOSSIER 2 : M. B. — Hypertension Artérielle Essentielle (Primaire) Grade 2, Stade 2 (atteinte rénale naissante).}

\textbf{2.1. Validation diagnostique \& PAM.}
\begin{itemize}
    \item Mesures J1: 160/100 ; J2: 150/95 ; J3: 170/105. Toutes $\ge$ 140/90 mmHg $\Rightarrow$ \textbf{HTA confirmée}.
    \item Classification : Grade 2 (PAS 160-179 et/ou PAD 100-109).
    \item PAM (Jour 3) = $PAD + \frac{1}{3}(PAS - PAD) = 105 + \frac{1}{3}(170 - 105) = 105 + \frac{65}{3} = 105 + 21,7 = \SI{126,7}{mmHg}$.
    \item PAM normale $\approx$ 90-100 mmHg. PAM élevée $\rightarrow$ surcharge de post-charge ventriculaire gauche $\rightarrow$ hypertrophie ventriculaire gauche (HVG).
\end{itemize}

\textbf{2.2. Facteurs de risque.}
\begin{center}
\begin{tabularx}{\linewidth}{|l|X|}\hline
\textbf{Non Modifiables} & \textbf{Modifiables (Cibles prioritaires)} \\ \hline
Âge (55 ans) & \textbf{Excès de Sel} (NaCl > 12-15 g/j vs reco < 5 g/j) : Rétention hydrosodée $\uparrow$ Volémie $\uparrow$ DC. \\
Sexe (Homme) & \textbf{Alcool} (Vin de palme/Bière chroniques) : Activation Sympathique + Effet direct vasculaire + Apport calorique. \\
Antécédents familiaux (Père AVC) & \textbf{Sédentarité} (Conduite > 12h) : Déconditionnement cardiovasculaire, prise de poids. \\
Ethnie (Noir africain : sensibilité au sel accrue, prévalence HTA plus forte) & \textbf{Tabagisme} (15 cig/j) : Endothélium dysfonctionnel (NO$\downarrow$, Endothéline$\uparrow$), rigidité artérielle, vasoconstriction aiguë. \\
& \textbf{Surpoids} (IMC 28,8) : Hyperinsulinémie/Insulinorésistance $\rightarrow$ SRAA/SNS activation + Compression rénale. \\ \hline
\end{tabularx}
\end{center}

\textbf{2.3. Physiopathologie Sel/Alcool $\rightarrow$ HTA.}
\begin{itemize}
    \item \textbf{Sel (NaCl) :} Augmente l'osmolarité plasmatique $\rightarrow$ Soif + Sécrétion ADH (vasopressine) $\rightarrow$ Rétention d'eau $\rightarrow$ \textbf{Augmentation du Volume Extracellulaire (VEC) et Volémie} $\rightarrow$ Augmentation du Débit Cardiaque (DC) $\rightarrow$ Augmentation PA (Loi de Frank-Starling). À long terme : « Autorégulation » vasculaire $\rightarrow$ Hypertrophie média des artérioles $\rightarrow$ \textbf{Augmentation des Résistances Périphériques (RPP)} (HTA « à RPP élevée »).
    \item \textbf{Alcool :} Activation aigüe du \textbf{Système Nerveux Sympathique} (libération noradrénaline) $\rightarrow$ Vasoconstriction + Tachycardie. Chronique : Lésion endothéliale, inflammation, activation SRAA, prise de poids (calories vides).
    \item \textbf{Synergie :} Sel + Alcool + Obésité + Tabac = « Syndrome métabolique » incomplet (glycémie limite 1,05 g/L) $\rightarrow$ Inflammation bas bruit + Stress oxydant $\rightarrow$ Dysfonction endothéliale (NO$\downarrow$) $\rightarrow$ Rigidité artérielle (PA de pouls élevée : 170-105 = 65 mmHg).
\end{itemize}

\textbf{2.4. Schéma fonctionnel de la régulation défaillante (HTA).}
\begin{popfigure}
\begin{tikzpicture}[
    node distance=1.0cm and 1.2cm,
    box/.style={rectangle, draw=popdark, thick, fill=popcream, rounded corners, minimum height=1.1cm, minimum width=2.4cm, align=center, font=\small},
    boxorange/.style={box, fill=poporange!15, draw=poporange},
    boxgreen/.style={box, fill=popgreen!15, draw=popgreen},
    arrow/.style={-Stealth, thick, popdark},
    orangearr/.style={-Stealth, thick, poporange},
    greenarr/.style={-Stealth, thick, popgreen},
]
% Nodes
\node[box, align=center] (stim){Stimuli chroniques :\\ Sel$\uparrow$, Alcool, Tabac,\\ Stress, Surpoids};
\node[boxorange, right=of stim, xshift=0.5cm, align=center] (sns){Système Nerveux\\ Sympathique \textbf{SURACTIVÉ}\\(Noradrénaline $\uparrow$)};
\node[boxorange, below=of sns, align=center] (sraa){SRAA \textbf{SURACTIVÉ}\\(Angio II $\uparrow$, Aldo $\uparrow$)};
\node[box, above right=of sns, xshift=1cm, align=center] (heart){Cœur :\\ FC $\uparrow$, Contractilité $\uparrow$\\ $\rightarrow$ \textbf{DC $\uparrow$}};
\node[box, right=of heart, align=center] (vessels){Artérioles :\\ Vasoconstriction $\uparrow\uparrow$\\ $\rightarrow$ \textbf{RPP $\uparrow\uparrow$}};
\node[box, below right=of vessels, align=center] (kidney){Rein :\\ Réabs Na\textsuperscript{+}/Eau $\uparrow$\\ (Aldo, SNS)\\ $\rightarrow$ \textbf{Volémie $\uparrow$}};
\node[boxgreen, below=of kidney, yshift=-0.5cm, align=center] (depress){Systèmes Dépresseurs :\\ Barorécepteurs (Réglage\\ à seuil plus haut = \textbf{Réinitialisation})\\ Peptides Natriurétiques (ANP/BNP)\\ \textbf{INSUFFISANTS}};
\node[box, right=of vessels, xshift=0.5cm, fill=poporange!20, draw=poporange, align=center] (pa){RÉSULTAT :\\ \textbf{PA = DC $\times$ RPP $\uparrow\uparrow$}\\ Grade 2 (170/105)\\ PAM = 127 mmHg};

% Arrows
\draw[orangearr] (stim) -- node[above, font=\tiny] {Agression} (sns);
\draw[orangearr] (stim) -- (sraa);
\draw[orangearr] (sns) -- (heart);
\draw[orangearr] (sns) -- (vessels);
\draw[orangearr] (sns) -- (kidney);
\draw[orangearr] (sraa) -- (vessels);
\draw[orangearr] (sraa) -- (kidney);
\draw[orangearr] (heart) -- (pa);
\draw[orangearr] (vessels) -- (pa);
\draw[orangearr] (kidney) -- (pa);
\draw[greenarr, dashed] (depress) -- node[left, font=\tiny] {Tentative frein} (sns);
\draw[greenarr, dashed] (depress) -- (sraa);
\draw[greenarr, dashed] (depress) -- (vessels);
\draw[greenarr, dashed] (depress) -- (kidney);

% Targets
\node[draw=popblue, thick, rounded corners, fill=popblue!10, font=\scriptsize, anchor=north west, align=center] at (pa.north east){CIBLES THÉRAPEUTIQUES :\\
\textbullet\ Diurétiques (Volémie$\downarrow$)\\
\textbullet\ IEC / ARA II (SRAA$\downarrow$)\\
\textbullet\ Bêta-bloquants / Centraux (SNS$\downarrow$)\\
\textbullet\ Calciumiques (RPP$\downarrow$)};
\end{tikzpicture}
\legende{Schéma fonctionnel de la régulation de la PA chez M. B. Les voies pressrices (SNS, SRAA) sont suractivées (flèches orange pleines) par les facteurs de risque. Les systèmes dépresseurs (barorécepteurs, ANP) sont débordés et « réinitialisés » à un seuil plus haut (flèches vertes pointillées). Les cibles médicamenteuses sont indiquées.}
\end{popfigure}

\textbf{2.5. Atteinte d'organe « Cible » (Rénale).}
\begin{itemize}
    \item Créatininémie \SI{14}{\mg.L^{-1}} (légèrement au-dessus norme H \SI{12}{\mg.L^{-1}}). Clairance de Cockcroft-Gault estimée (poids estimé 85 kg pour IMC 28,8 / 1,72 m) : $Cl_{Cr} = \frac{(140-55) \times 85}{0,814 \times 14} \approx \SI{63}{mL.min^{-1}}$.
    \item \textbf{Stade 2 KDIGO (GFR 60-89 mL/min) : Insuffisance rénale légère.} (Correction : le seuil du Stade 3a est 45 mL/min).
    \item Protéinurie (+) : Signe de \textbf{glomérulosclérose hyalinique} (néphroangiosclérose hypertensive) ou atteinte glomérulaire diabétique débutante (glycémie 1,05 g/L limite). Marqueur pronostique majeur (risque CV et rénal).
\end{itemize}

\vspace{1em}
\textbf{3. SYNTHÈSE COMPARATIVE \& AFFICHE DE SENSIBILISATION}

\textbf{3.1. Tableau comparatif.}
\begin{center}
\begin{tabularx}{\linewidth}{|l|X|X|}\hline
\textbf{Paramètre} & \textbf{Mme A. (DT2)} & \textbf{M. B. (HTA Grade 2)} \\ \hline
\textbf{Régulation défaillante} & Boucle insuline/glucose : \textbf{Résistance périphérique} (effecteurs sourds) + $\beta$-cellules épuisées. & Boucle PA : \textbf{Surcharge volémique} + \textbf{Vasoconstriction excessive} (SNS/SRAA) + Barorécepteurs réinitialisés. \\ \hline
\textbf{Facteurs communs} & Obésité abdominale, Sédentarité, Alimentation déséquilibrée (sucre/gras/sel), Stress oxydant, Inflammation bas bruit. & Obésité, Sédentarité, Alimentation (Sel/Alcool/Gras), Tabac, Stress. \\ \hline
\textbf{Risque partagé} & \textbf{Risque Cardiovasculaire Majeur} (AVC, IDM, Insuffisance Cardiaque, Néphropathie). Syndrome métabolique (Mme A. confirmé, M. B. probable). \\ \hline
\textbf{Prise en charge 1ère ligne} & \textbf{Hygiène de vie} (Régime méditerranéen/adapté local, Activité Physique 150 min/sem, Perte de poids 5-10\%). \textbf{Metformine} (sensibilisateur insuline). & \textbf{Hygiène de vie} (Régime DASH : Sel < 5g, K$\uparrow$, Alcool$\downarrow$, Arrêt Tabac, AP). \textbf{Bithérapie} d'emblée (Grade 2) : IEC/ARA II + Diurétique ou Calciumique. \\ \hline
\end{tabularx}
\end{center}

\textbf{3.2. Proposition d'Affiche de Sensibilisation (Contenu textuel pour maquette).}
\begin{center}
\fbox{\begin{minipage}{0.9\linewidth}
\textbf{\large \textsc{VOTRE SANTÉ, VOTRE TRÉSOR !}} \\
\textit{« Nkwa a si nkwa » (La santé, c'est la vie - Ewondo) \quad « Ndjock a mbot » (La force est dans le corps - Bassa)}

\vspace{0.5em}
\noindent\textbf{\textbullet\ MESSAGE 1 : BOUGEZ, MANGEZ MALIN !}
\begin{itemize}
    \item[--] 30 min de marche rapide / jour (aller au marché, descendre du taxi plus tôt).
    \item[--] Assiette : 1/2 Légumes (feuilles : eru, folong, kpwem), 1/4 Protéines (poisson, haricots, œufs), 1/4 Féculents (plantain, manioc, riz brun).
    \item[--] \textbf{STOP} : Beignets, boissons sucrées, cubes maggi excessifs, fritures, vin de palme/bière tous les jours.
\end{itemize}

\textbf{\textbullet\ MESSAGE 2 : CONNAISSEZ VOS CHIFFRES !}
\begin{itemize}
    \item \textbf{Glycémie à jeun :} Normale < 1,10 g/L. Diabète si $\ge$ 1,26 g/L (2 fois).
    \item \textbf{Tension :} Normale < 140/90 mmHg. HTA si $\ge$ 140/90 (3 mesures, 3 jours).
    \item \textbf{Tour de taille :} Homme < 94 cm, Femme < 80 cm.
    \item \textbf{Au Cameroun :} 1 adulte sur 3 a de l'HTA ! 1 sur 12 a du Diabète ! \textbf{La moitié l'ignore.}
\end{itemize}

\textbf{\textbullet\ MESSAGE 3 : DÉPISTEZ-VOUS, SOIGNEZ-VOUS !}
\begin{itemize}
    \item Dépistage GRATUIT / PEU CHER au CSI « La Paix » (Mfoundi) et centres de santé de district.
    \item Traitement = \textbf{À VIE} (souvent). Ne pas arrêter sans avis médecin.
    \item \textbf{Urgences :} Malaise, paralysie brutale, douleur thoracique $\rightarrow$ \textbf{APPELEZ LE 112 (SAMU)}.
\end{itemize}

\vspace{0.5em}
\noindent\textbf{CENTRE DE SANTÉ INTÉGRÉ « LA PAIX »} \\
Quartier Mfoundi, derrière le Marché Central. \\
Tél : 22 20 XX XX \quad | \quad Centre Antipoison : 22 21 XX XX \\
\textit{Venez vous faire dépister cette semaine !}
\end{minipage}}
\end{center}

\end{exoresolu}

\courssub{Bilan de l'activité}
Cette activité d'intégration a permis de :
\begin{enumerate}
    \item \textbf{Valider les compétences diagnostiques :} Différencier un diabète de type 2 (insulinorésistance avec hyperinsulinisme relatif puis échec $\beta$) d'une HTA essentielle multifactorielle (volume-dépendante et vasoconstrictrice) à partir de données brutes (courbes, tableaux).
    \item \textbf{Modéliser la physiopathologie :} La construction des schémas fonctionnels a mis en évidence la \textbf{rupture de la rétroaction négative} dans les deux cas : résistance à l'effet de l'hormone (insuline) pour Mme A. ; réinitialisation du point de consigne des barorécepteurs et suractivation des systèmes pressteurs pour M. B.
    \item \textbf{Relier l'épidémiologie locale à la clinique :} L'obésité abdominale, la sédentarité, l'excès sel/sucre/alcool et le tabagisme sont les piliers communs de la transition épidémiologique au Cameroun (double fardeau maladies infectieuses / MNT).
    \item \textbf{Traduire le savoir savant en savoir agir :} La conception de l'affiche a imposé la simplification du langage médical (vocabulaire accessible, langues locales, pictogrammes) et l'intégration des ressources locales (CSI, numéros d'urgence), compétence clé de l'éducation à la santé (Savoir-être : responsabilité, communication, citoyenneté).
    \item \textbf{Ancrer la rigueur quantitative :} Calcul de l'amplitude glycémique, de la PAM, interprétation des seuils OMS/ESC avec unités correctes (g/L, mmHg, $\mu$UI/mL, \%), lecture de courbes HGPO.
\end{enumerate}
\emph{Cette mise en situation prépare directement aux épreuves pratiques (ECE) et écrites du Baccalauréat (analyse de documents, schémas fonctionnels, argumentation éthique/sociale).}

\begin{aretenir}[Points clés pour le Bac]
\begin{itemize}
    \item \textbf{Diabète Type 2 :} Hyperglycémie + Insulinémie \textbf{inappropriée} (normale ou élevée mais insuffisante pour la glycémie) = Insulinorésistance. Schéma : Flèche insulinique barrée au niveau effecteurs (GLUT4).
    \item \textbf{HbA1c :} Miroir de l'équilibre sur 3 mois. Cible < 7\% (individualisée). Attention aux hémoglobinopathies (drépanocytose fréquente au Cameroun).
    \item \textbf{HTA :} Diagnostic sur 3 séries de mesures (cabinet) ou automesure/MAPA. Grade = chiffres. Stade = atteintes organes cibles (Cœur, Rein, Vaisseaux, Œil, Cerveau).
    \item \textbf{PAM = PAD + 1/3(PAS-PAD).} Normale ~95 mmHg. Élévation = surcharge cardiaque (HVG) + risque vasculaire.
    \item \textbf{Physiopathologie HTA :} Sel $\rightarrow$ Volémie $\uparrow$ (phase initiale) puis RPP $\uparrow$ (phase établie). Alcool/Tabac $\rightarrow$ SNS $\uparrow$ + Endothélium $\downarrow$.
    \item \textbf{Prévention commune :} Activité physique, Régime DASH/Méditerranéen (sel < 5g, fruits/légumes, poisson), Arrêt tabac, Modération alcool, Dépistage régulier.
\end{itemize}
\end{aretenir}

\begin{savaistu}[Contexte camerounais : Double fardeau]
Au Cameroun, la prévalence de l'HTA est estimée à \textbf{29,7\%} chez les adultes (Enquête STEPS 2003, confirmée par études récentes) et celle du Diabète à \textbf{6-8\%} (IDF Atlas). La transition nutritionnelle (urbanisation, alimentation transformée, sédentarité) explique l'explosion des Maladies Non Transmissibles (MNT). Le système de santé (CSI, Hôpitaux de district) intègre désormais la prise en charge des MNT dans le paquet d'activités minimum. Le dépistage communautaire (marchés, gares, églises) est une stratégie prioritaire du Plan Stratégique National de Lutte contre les MNT 2020-2025.
\end{savaistu}

\newpage

\coursec{Méthodes et exercices résolus}

\courssub{Fiche méthode 1 : Expliquer les techniques de mesure de la glycémie et de la pression artérielle}

\begin{methode}[Décrire et justifier une technique de mesure biologique]
Pour expliquer une technique de mesure au Baccalauréat, il faut toujours articuler quatre points dans l'ordre :
\begin{enumerate}
\item \textbf{Nommer l'appareil et le principe physique ou biochimique} : lecteur de glycémie (glucomètre) à bandelettes réactives pour la glycémie ; sphygmomanomètre (à brassard, manuel ou électronique) pour la pression artérielle.
\item \textbf{Décrire le protocole étape par étape} : pour la glycémie, prélèvement d'une goutte de sang capillaire au bout du doigt (après désinfection), dépôt sur la bandelette, lecture en quelques secondes ; pour la pression artérielle, sujet assis au repos depuis 5 minutes, brassard au bras à hauteur du c\oe ur, gonflage au-delà de la pression systolique puis dégonflage progressif avec écoute des bruits de Korotkoff au stéthoscope (méthode manuelle).
\item \textbf{Définir les valeurs mesurées et leurs unités} : glycémie en g/L (ou mmol/L : $\SI{1}{g/L} \approx \SI{5,5}{mmol/L}$, car la masse molaire du glucose est $M = \SI{180}{g/mol}$) ; pression artérielle maximale (systolique) et minimale (diastolique) en cmHg (ou mmHg).
\item \textbf{Citer les valeurs de référence et les conditions de validité} : glycémie à jeun normale entre 0,70 et \SI{1,10}{g/L} ; pression artérielle normale voisine de 12/7 cmHg (soit 120/70 mmHg). Préciser toujours les conditions : mesure de la glycémie à jeun (8 à 12 h sans repas) ou post-prandiale (2 h après le début du repas) ; mesure de la pression au repos, en position assise, sans café, tabac ni effort récent.
\end{enumerate}
\textbf{Réflexe rédactionnel :} ne jamais écrire « on mesure le diabète » ; on mesure la \cle{glycémie}, qui est le taux de glucose sanguin, paramètre dont la perturbation chronique définit le diabète.
\end{methode}

\begin{exoresolu}[Mesures au centre de santé de Bafoussam]
Lors d'une campagne de dépistage au centre de santé intégré de Bafoussam, un infirmier mesure chez M. Kamga, 52 ans, à jeun depuis 10 heures : une glycémie de \SI{1,32}{g/L} au glucomètre et une pression artérielle de 15/9 cmHg.\\
1) Décrire le protocole de chacune de ces deux mesures.\\
2) Comparer chaque résultat aux valeurs de référence et conclure sur l'état de santé du patient.\\
3) Expliquer pourquoi l'infirmier demande à M. Kamga de revenir un autre jour, à jeun, pour une seconde mesure de glycémie en laboratoire.
\tcblower
\textbf{1) Protocoles de mesure.}\\
\emph{Mesure de la glycémie :} l'infirmier désinfecte le bout du doigt du patient à l'alcool, pique la peau avec une lancette stérile, dépose la goutte de sang capillaire sur une bandelette réactive insérée dans le glucomètre. Le glucose réagit avec une enzyme (glucose oxydase) fixée sur la bandelette ; l'appareil convertit l'intensité de la réaction en concentration et affiche la glycémie en g/L en quelques secondes.\\
\emph{Mesure de la pression artérielle :} le patient est assis, au repos depuis au moins 5 minutes, le bras nu posé à hauteur du c\oe ur. Le brassard est placé autour du bras, gonflé jusqu'à interruption du flux sanguin dans l'artère humérale, puis dégonflé lentement. Le premier bruit perçu au stéthoscope (reprise du flux turbulent) correspond à la pression \textbf{maximale} (systolique) ; la disparition des bruits (flux redevenu silencieux) correspond à la pression \textbf{minimale} (diastolique).\\[2mm]
\textbf{2) Comparaison aux valeurs de référence.}\\
Glycémie mesurée : \SI{1,32}{g/L} $>$ \SI{1,10}{g/L} (limite supérieure de la norme à jeun) : il y a \cle{hyperglycémie} à jeun.\\
Pression mesurée : 15/9 cmHg, alors que la norme est voisine de 12/7 cmHg : la maximale (15 $>$ 14 cmHg) est élevée et la minimale (9 cmHg) est à la valeur limite : on suspecte une \cle{hypertension artérielle}.\\[2mm]
\textbf{3) Justification de la seconde mesure.} Une glycémie capillaire unique au glucomètre est un examen de \emph{dépistage}, pas de diagnostic : le résultat peut être faussé par un stress, une prise alimentaire récente mal déclarée ou une imprécision de l'appareil. Le diagnostic de diabète exige une glycémie veineuse de laboratoire, à jeun, supérieure ou égale à \SI{1,26}{g/L}, \textbf{confirmée une seconde fois} un autre jour.\\
\textbf{Conclusion :} M. Kamga présente une hyperglycémie à jeun et une pression artérielle élevée qui doivent être confirmées par des mesures répétées en conditions standardisées avant tout diagnostic de diabète ou d'hypertension artérielle.
\end{exoresolu}

\courssub{Fiche méthode 2 : Analyser les courbes de variation de la glycémie avant et après un repas}

\begin{methode}[Exploiter une courbe de glycémie]
\begin{enumerate}
\item \textbf{Lire les axes avec leurs unités} : en abscisse le temps (en min ou h, l'origine étant souvent le début du repas ou de l'ingestion de glucose) ; en ordonnée la glycémie en g/L.
\item \textbf{Repérer les points remarquables} : valeur initiale (glycémie à jeun), maximum (pic post-prandial, atteint vers 30 min à 1 h), temps de retour à la normale (environ 2 h chez le sujet sain).
\item \textbf{Décrire la courbe en trois phases} avec le vocabulaire imposé : phase d'\emph{hyperglycémie} (absorption intestinale du glucose), phase de \emph{décroissance} (action de l'insuline : stockage du glucose sous forme de glycogène hépatique et musculaire, entrée du glucose dans les cellules), phase de \emph{retour à la glycémie de base}.
\item \textbf{Comparer les courbes} (sujet sain / diabétique) en citant des valeurs chiffrées précises relevées sur le graphique : chez le diabétique, glycémie à jeun plus élevée, pic plus haut et retour à la normale retardé ou absent.
\item \textbf{Conclure en reliant la courbe au mécanisme} : la régulation hormonale (insuline hypoglycémiante, glucagon hyperglycémiant) maintient la glycémie autour d'une valeur à peu près constante (environ \SI{1}{g/L}) : c'est un exemple d'\emph{homéostasie}.
\end{enumerate}
\textbf{Formule utile :} variation relative $= \dfrac{\text{valeur finale} - \text{valeur initiale}}{\text{valeur initiale}} \times 100$ (en \%).
\end{methode}

\begin{exoresolu}[Hyperglycémie provoquée chez deux sujets]
On fait ingérer à jeun \SI{75}{g} de glucose dissous dans de l'eau à deux sujets : A (sain) et B (diabétique). On mesure leur glycémie toutes les 30 minutes :
\begin{center}
\begin{tabularx}{\linewidth}{|l|X|X|X|X|X|X|}\hline
\rowcolor{popblueL} Temps (min) & 0 & 30 & 60 & 90 & 120 & 180 \\ \hline
Sujet A (g/L) & 0,90 & 1,40 & 1,20 & 1,00 & 0,90 & 0,88 \\ \hline
Sujet B (g/L) & 1,45 & 2,30 & 2,60 & 2,40 & 2,10 & 1,80 \\ \hline
\end{tabularx}
\end{center}
1) Décrire l'évolution de la glycémie du sujet A et l'expliquer.\\
2) Comparer avec le sujet B en citant trois différences chiffrées.\\
3) Calculer l'augmentation relative maximale de la glycémie chez chaque sujet et conclure.
\tcblower
\textbf{1) Sujet A (sain).} À jeun, sa glycémie est de \SI{0,90}{g/L} (normale). Après ingestion, elle s'élève jusqu'à un pic de \SI{1,40}{g/L} à $t = 30$ min : c'est l'\emph{hyperglycémie post-prandiale}, due à l'absorption du glucose par l'intestin grêle. Puis elle décroît et revient à \SI{0,90}{g/L} à $t = 120$ min : l'élévation de la glycémie a stimulé les cellules $\beta$ des îlots de Langerhans du pancréas, qui sécrètent de l'\textbf{insuline}. Cette hormone provoque l'entrée du glucose dans les cellules (muscles, tissu adipeux) et sa transformation en glycogène dans le foie (\emph{glycogénèse}), ce qui fait baisser la glycémie.\\[2mm]
\textbf{2) Comparaison avec le sujet B.}
\begin{itemize}
\item Glycémie à jeun : \SI{1,45}{g/L} chez B contre \SI{0,90}{g/L} chez A : B part d'une hyperglycémie chronique ($> \SI{1,26}{g/L}$).
\item Pic : \SI{2,60}{g/L} à 60 min chez B contre \SI{1,40}{g/L} à 30 min chez A : pic plus élevé \emph{et} plus tardif.
\item À 180 min, B est encore à \SI{1,80}{g/L}, loin de sa valeur de base, alors qu'A est revenu à \SI{0,88}{g/L} dès 120 min : le retour à la normale est impossible chez B.
\end{itemize}
\textbf{3) Augmentations relatives maximales.}
\begin{align*}
\text{Sujet A : } & \frac{1,40 - 0,90}{0,90} \times 100 \approx 55,6\ \% \\
\text{Sujet B : } & \frac{2,60 - 1,45}{1,45} \times 100 \approx 79,3\ \%
\end{align*}
\textbf{Conclusion :} chez le sujet sain, la régulation insulinique ramène la glycémie à sa valeur de base en 2 h ; chez le diabétique, l'insuline est absente, insuffisante ou inefficace, d'où une hyperglycémie majeure ($+79,3\ \%$ contre $+55,6\ \%$) et durable : la glycémie n'est plus régulée.
\end{exoresolu}

\courssub{Fiche méthode 3 : Analyser et interpréter les résultats expérimentaux portant sur le foie et le pancréas}

\begin{methode}[Interpréter une expérience de physiologie (ablation, ligature, injection)]
\begin{enumerate}
\item \textbf{Identifier la variable expérimentale} (ce qu'on modifie : ablation du pancréas, ligature du canal pancréatique, injection d'extrait, perfusion du foie) et la \textbf{variable mesurée} (glycémie, glycosurie, taux de glycogène hépatique).
\item \textbf{Comparer systématiquement l'animal témoin et l'animal expérimenté} : sans témoin, aucune conclusion n'est valable.
\item \textbf{Formuler l'interprétation sous forme de lien de cause à effet} : « l'ablation de X entraîne Y, donc X est nécessaire à / produit une substance qui... ».
\item \textbf{Connaître les expériences fondatrices à citer} :
\begin{itemize}
\item Ablation du pancréas chez le chien (von Mering et Minkowski, 1889) : hyperglycémie et glycosurie $\Rightarrow$ le pancréas prévient le diabète.
\item Ligature du canal pancréatique : le suc pancréatique (exocrine) ne s'écoule plus, le pancréas exocrine dégénère, mais \emph{pas de diabète} $\Rightarrow$ le facteur antidiabétique est sécrété par les îlots de Langerhans (sécrétion \emph{endocrine}, libérée dans le sang).
\item Injection d'extrait d'îlots (insuline, Banting et Best, 1921) à l'animal dépancréaté : disparition de l'hyperglycémie $\Rightarrow$ l'insuline est l'hormone hypoglycémiante.
\item Perfusion du foie isolé avec du glucose : formation de glycogène $\Rightarrow$ le foie stocke le glucose (glycogénèse) et le libère (glycogénolyse) sous contrôle hormonal.
\end{itemize}
\item \textbf{Conclure par le rôle de l'organe} dans la régulation, en précisant le sens de l'action (hypoglycémiant ou hyperglycémiant).
\end{enumerate}
\end{methode}

\begin{exoresolu}[Trois chiens, trois protocoles]
On réalise les expériences suivantes sur trois lots de chiens :
\begin{center}
\begin{tabularx}{\linewidth}{|l|X|X|}\hline
\rowcolor{popblueL} Lot & Protocole & Résultat \\ \hline
1 (témoin) & Aucune intervention & Glycémie stable à \SI{1,00}{g/L} \\ \hline
2 & Ablation totale du pancréas & Glycémie : \SI{2,50}{g/L} ; glucose dans les urines ; mort en quelques semaines \\ \hline
3 & Ligature du canal excréteur du pancréas & Atrophie des acini (pancréas exocrine) ; glycémie normale à \SI{1,05}{g/L} \\ \hline
\end{tabularx}
\end{center}
1) Interpréter les résultats des lots 2 et 3.\\
2) Quelle expérience complémentaire permettrait d'identifier la substance responsable ? Préciser le résultat attendu.\\
3) Conclure sur le double rôle du pancréas.
\tcblower
\textbf{1) Interprétation.}\\
\emph{Lot 2 :} l'ablation du pancréas provoque une hyperglycémie majeure (\SI{2,50}{g/L}) et une glycosurie (le seuil rénal de réabsorption du glucose, environ \SI{1,80}{g/L}, est dépassé), signes d'un diabète expérimental. Le pancréas est donc nécessaire au maintien d'une glycémie normale.\\
\emph{Lot 3 :} la ligature du canal pancréatique supprime la fonction exocrine (le suc digestif ne peut plus s'écouler, les acini dégénèrent), pourtant la glycémie reste normale. Le facteur qui régule la glycémie n'est donc \emph{pas} véhiculé par le canal excréteur : il est libéré directement dans le sang par une structure épargnée, les \textbf{îlots de Langerhans} (pancréas endocrine).\\[2mm]
\textbf{2) Expérience complémentaire.} On injecte à un chien du lot 2 (dépancréaté) un extrait aqueux d'îlots de Langerhans (insuline). Résultat attendu : la glycémie redescend vers \SI{1}{g/L} et la glycosurie disparaît. Cela démontre que les îlots sécrètent une hormone hypoglycémiante, l'\textbf{insuline}.\\[2mm]
\textbf{3) Conclusion.} Le pancréas est une \cle{glande mixte} : exocrine (suc pancréatique digestif acheminé par un canal vers le duodénum) et endocrine (insuline sécrétée par les cellules $\beta$ et glucagon par les cellules $\alpha$ des îlots, libérés dans le sang). Seule la fonction endocrine intervient dans la régulation de la glycémie.
\end{exoresolu}

\courssub{Fiche méthode 4 : Réaliser un schéma fonctionnel de la régulation de la glycémie et exploiter des analyses médicales}

\begin{methode}[Construire un schéma fonctionnel de régulation hormonale]
\begin{enumerate}
\item \textbf{Placer la variable régulée au centre} : la glycémie (valeur de consigne $\approx \SI{1}{g/L}$).
\item \textbf{Identifier les deux déviations} : hyperglycémie (après un repas) et hypoglycémie (jeûne, effort).
\item \textbf{Relier chaque déviation à son récepteur et à sa réponse hormonale} : hyperglycémie $\rightarrow$ cellules $\beta$ des îlots $\rightarrow$ insuline $\rightarrow$ effets hypoglycémiants (glycogénèse hépatique, entrée cellulaire du glucose) ; hypoglycémie $\rightarrow$ cellules $\alpha$ $\rightarrow$ glucagon $\rightarrow$ glycogénolyse hépatique (libération de glucose).
\item \textbf{Fermer la boucle par une rétroaction négative} : la correction de la déviation supprime le stimulus (flèche en « $-$ »). C'est le principe de tout schéma de régulation : stimulus $\rightarrow$ capteur $\rightarrow$ effecteur $\rightarrow$ réponse $\rightarrow$ extinction du stimulus.
\item \textbf{Pour exploiter une analyse médicale} : comparer chaque paramètre (glycémie à jeun, glycosurie, hémoglobine glyquée HbA$_{1c}$) aux valeurs de référence, puis rattacher l'anomalie au type de diabète (type 1 : insuline absente ; type 2 : insulino-résistance).
\end{enumerate}
\textbf{Réflexe :} un schéma fonctionnel sans flèche de rétroaction négative est incomplet et coûte des points.
\end{methode}

\begin{exoresolu}[Bilan biologique de deux patientes à l'hôpital de Yaoundé]
\begin{center}
\begin{tabularx}{\linewidth}{|l|X|X|X|}\hline
\rowcolor{popblueL} Paramètre & Valeurs de référence & Mme A (25 ans, maigre) & Mme B (55 ans, obèse) \\ \hline
Glycémie à jeun & 0,70 -- 1,10 g/L & 2,40 g/L & 1,90 g/L \\ \hline
Glycosurie & absente & présente & présente \\ \hline
Insulinémie & 5 -- 25 µUI/mL & quasi nulle & 40 µUI/mL \\ \hline
\end{tabularx}
\end{center}
1) Réaliser le schéma fonctionnel de la régulation de la glycémie chez le sujet sain.\\
2) Identifier le type de diabète de chaque patiente en justifiant.\\
3) Proposer un traitement adapté à chacune.
\tcblower
\textbf{1) Schéma fonctionnel.}
\begin{popfigure}
\begin{tikzpicture}[node distance=1.1cm, >=Stealth, font=\small]
\node[draw=popblue, fill=popblueL, rounded corners, minimum width=3.2cm, align=center] (gly) {Glycémie\\ ($\approx$ \SI{1}{g/L})};
\node[draw=poporange, fill=poporangeL, rounded corners, left=2.6cm of gly, align=center] (hyper) {Hyperglycémie\\ (repas)};
\node[draw=poppurple, fill=poppurpleL, rounded corners, right=2.6cm of gly, align=center] (hypo) {Hypoglycémie\\ (jeûne, effort)};
\node[draw=popgreen, fill=popgreenL, rounded corners, below=1.2cm of hyper, align=center] (beta) {Cellules $\beta$\\ insuline};
\node[draw=popgreen, fill=popgreenL, rounded corners, below=1.2cm of hypo, align=center] (alpha) {Cellules $\alpha$\\ glucagon};
\node[draw=popdark, fill=popcream, rounded corners, below=1.2cm of beta, align=center] (eff1) {Glycogénèse,\\ entrée du glucose\\ dans les cellules};
\node[draw=popdark, fill=popcream, rounded corners, below=1.2cm of alpha, align=center] (eff2) {Glycogénolyse\\ hépatique};
\draw[->, thick, poporange] (hyper) -- (beta);
\draw[->, thick, poppurple] (hypo) -- (alpha);
\draw[->, thick, popgreen] (beta) -- (eff1);
\draw[->, thick, popgreen] (alpha) -- (eff2);
\draw[->, thick, popblue, dashed] (eff1.south) .. controls +(0,-0.8) and +(0,-2.2) .. node[below, align=center]{rétroaction $-$ : baisse de la glycémie} (gly.south);
\draw[->, thick, popblue, dashed] (eff2.south) .. controls +(0,-0.8) and +(0,-2.2) .. node[below, align=center]{rétroaction $-$ : hausse de la glycémie} (gly.south);
\end{tikzpicture}
\legende{Schéma fonctionnel de la régulation de la glycémie : deux boucles antagonistes (insuline hypoglycémiante, glucagon hyperglycémiant) à rétroaction négative.}
\end{popfigure}
\textbf{2) Typologie.} Mme A : hyperglycémie sévère avec \emph{insulinémie quasi nulle} $\Rightarrow$ le pancréas ne produit plus d'insuline (destruction auto-immune des cellules $\beta$) : c'est un \textbf{diabète de type 1} (insulino-dépendant), fréquent chez le sujet jeune et maigre. Mme B : hyperglycémie avec insulinémie \emph{élevée} (40 µUI/mL) $\Rightarrow$ l'insuline est produite mais inefficace : les cellules cibles résistent à son action ; l'obésité en est un facteur favorisant : c'est un \textbf{diabète de type 2} (insulino-résistant).\\[2mm]
\textbf{3) Traitement.} Mme A : injections quotidiennes d'\textbf{insuline} (le seul traitement possible, l'hormone étant absente), associées à une surveillance glycémique. Mme B : mesures hygiéno-diététiques d'abord (régime pauvre en sucres rapides, activité physique régulière, perte de masse grasse), éventuellement complétées par des antidiabétiques oraux qui améliorent la sensibilité à l'insuline ; l'insuline n'est pas le traitement de première intention.\\
\textbf{Conclusion :} la mesure conjointe de la glycémie et de l'insulinémie permet de distinguer un diabète de type 1 (déficit de sécrétion) d'un diabète de type 2 (défaut d'action), et donc d'adapter le traitement.
\end{exoresolu}

\courssub{Fiche méthode 5 : Analyser les données et courbes illustrant les variations de la pression artérielle}

\begin{methode}[Exploiter des données de pression artérielle]
\begin{enumerate}
\item \textbf{Définir les deux paramètres} : pression maximale (systolique, éjection du sang par le ventricule gauche) et pression minimale (diastolique, relâchement cardiaque), exprimées en cmHg ou mmHg.
\item \textbf{Identifier les paramètres physiologiques qui font varier la pression} : débit cardiaque (fréquence cardiaque $\times$ volume d'éjection systolique), volume sanguin, élasticité et diamètre des artères (résistances périphériques), viscosité du sang. Relation qualitative : $P \propto \text{débit cardiaque} \times \text{résistances périphériques}$.
\item \textbf{Interpréter les variations} : à l'effort ou sous l'effet du stress, l'adrénaline et l'accélération cardiaque (nerfs sympathiques) augmentent le débit cardiaque $\Rightarrow$ la pression s'élève ; au repos, le nerf vague (parasympathique) ralentit le c\oe ur $\Rightarrow$ la pression baisse.
\item \textbf{Distinguer variation physiologique et pathologie} : une élévation transitoire à l'effort est normale ; une pression $\geq 14/9$ cmHg mesurée au repos, à plusieurs reprises, définit l'hypertension artérielle.
\item \textbf{Citer des valeurs chiffrées} du document dans la réponse : jamais de description sans appui sur les données.
\end{enumerate}
\end{methode}

\begin{exoresolu}[Pression artérielle au repos et à l'effort]
On mesure la pression artérielle et la fréquence cardiaque d'un élève de Terminale dans trois situations :
\begin{center}
\begin{tabularx}{\linewidth}{|l|X|X|X|}\hline
\rowcolor{popblueL} Situation & Repos assis & Après 10 min de course & Après 15 min de récupération \\ \hline
Pression max (cmHg) & 12,0 & 16,5 & 12,5 \\ \hline
Pression min (cmHg) & 7,5 & 8,0 & 7,5 \\ \hline
Fréquence cardiaque (batt/min) & 70 & 150 & 80 \\ \hline
\end{tabularx}
\end{center}
1) Analyser les variations observées pendant l'effort.\\
2) Expliquer ces variations par les mécanismes de régulation nerveuse et hormonale.\\
3) L'élève présente-t-il une hypertension artérielle ? Justifier.
\tcblower
\textbf{1) Analyse.} Pendant l'effort, la fréquence cardiaque passe de 70 à 150 batt/min (soit $+114\ \%$) et la pression maximale s'élève de 12,0 à 16,5 cmHg ($+4,5$ cmHg), tandis que la minimale varie peu (7,5 à 8,0 cmHg). Après récupération, toutes les valeurs reviennent quasiment à leurs valeurs de repos : la variation est \emph{transitoire et réversible}.\\[2mm]
\textbf{2) Mécanismes.} L'effort augmente les besoins en dioxygène des muscles. Les barorécepteurs des arcs aortique et carotidien et les centres nerveux du bulbe rachidien déclenchent, via les \textbf{nerfs sympathiques}, une accélération du rythme cardiaque et une augmentation de la force des contractions : le débit cardiaque augmente, donc la pression systolique s'élève. Parallèlement, la médullo-surrénale libère de l'\textbf{adrénaline}, qui renforce et prolonge ces effets (accélération cardiaque, vasoconstriction des territoires non prioritaires). Au repos, le \textbf{nerf vague} (parasympathique) ralentit le c\oe ur et la pression redescend : c'est une régulation à rétroaction négative.\\[2mm]
\textbf{3) Hypertension ?} Non. L'hypertension artérielle se définit par une pression $\geq 14/9$ cmHg mesurée \emph{au repos}, de façon \emph{répétée}. Ici, la pression de repos est normale (12,0/7,5 cmHg) et l'élévation n'apparaît qu'à l'effort : il s'agit d'une adaptation physiologique normale, pas d'une pathologie.\\
\textbf{Conclusion :} la pression artérielle est un paramètre variable, ajusté en permanence aux besoins de l'organisme par le système nerveux autonome et l'adrénaline ; seule une élévation permanente au repos est pathologique.
\end{exoresolu}

\courssub{Fiche méthode 6 : Identifier causes, symptômes et complications de l'hypertension et élaborer un outil de sensibilisation}

\begin{methode}[Construire une démarche de prévention et de sensibilisation]
\begin{enumerate}
\item \textbf{Définir la pathologie avec ses seuils chiffrés} : hypertension artérielle (HTA) = pression $\geq 14/9$ cmHg au repos, confirmée sur plusieurs mesures ; hypotension = pression nettement inférieure à la normale (en dessous d'environ 9/6 cmHg).
\item \textbf{Classer les causes} : causes avérées de l'HTA (excès de sel, obésité, sédentarité, stress chronique, alcool, tabac, hérédité, âge, certaines maladies rénales) ; l'hypotension peut résulter d'une déshydratation, d'une hémorragie, d'une insuffisance de sécrétion d'adrénaline ou d'une station debout prolongée.
\item \textbf{Associer symptômes et complications} : l'HTA est souvent \emph{silencieuse} (d'où le dépistage) ; ses complications sont l'infarctus du myocarde, l'accident vasculaire cérébral (AVC), l'insuffisance rénale, la rétinopathie. L'hypotension provoque vertiges, malaises, syncopes.
\item \textbf{Proposer des moyens de lutte hiérarchisés} : prévention (réduire le sel à moins de \SI{5}{g/j}, activité physique régulière, limiter alcool et tabac, gérer le stress, dépistage régulier) puis traitement médical (antihypertenseurs, pris durablement et sans interruption).
\item \textbf{Pour l'outil de sensibilisation} (affiche, sketch, spot radio) : un message clair, un public cible, des chiffres, des gestes concrets, un vocabulaire adapté au contexte local (marchés, lieux de culte, radios communautaires).
\end{enumerate}
\end{methode}

\begin{exoresolu}[Campagne « Octobre rouge » à Douala]
Dans un quartier de Douala, une enquête révèle que 30 \% des adultes de plus de 40 ans ont une pression artérielle supérieure à 14/9 cmHg, et que la majorité l'ignore. La consommation moyenne de sel y est estimée à \SI{12}{g} par jour et par habitant.\\
1) Définir l'hypertension artérielle et expliquer pourquoi elle est dite « tueur silencieux ».\\
2) Citer deux facteurs de risque présents dans ce quartier et expliquer leur rôle.\\
3) Rédiger le contenu d'une affiche de sensibilisation (titre, trois messages clés, un appel à l'action).
\tcblower
\textbf{1) Définition.} L'hypertension artérielle est une élévation permanente de la pression du sang dans les artères ($\geq 14/9$ cmHg au repos, confirmée par des mesures répétées). Elle est dite « tueur silencieux » car elle ne provoque généralement \emph{aucun symptôme} pendant des années tout en endommageant progressivement la paroi des artères, le c\oe ur, le cerveau et les reins : la première manifestation peut être une complication grave (AVC, infarctus, insuffisance rénale). Seule la mesure régulière de la pression permet de la détecter.\\[2mm]
\textbf{2) Facteurs de risque.}
\begin{itemize}
\item \emph{Excès de sel} (\SI{12}{g/j}, soit plus du double de la recommandation de \SI{5}{g/j}) : le sodium retient l'eau dans le sang, ce qui augmente le volume sanguin, donc la pression artérielle.
\item \emph{Absence de dépistage} (la majorité ignore son état) : sans mesure, l'HTA évolue sans traitement jusqu'aux complications. On peut ajouter la sédentarité urbaine et le stress, facteurs favorisants classiques.
\end{itemize}
\textbf{3) Affiche de sensibilisation.}\\
\emph{Titre :} « Connais ta tension, protège ta vie ! »\\
\emph{Trois messages clés :}
\begin{itemize}
\item L'hypertension ne fait pas mal, mais elle détruit le c\oe ur, le cerveau et les reins : près d'un adulte sur trois après 40 ans est concerné.
\item Moins de sel dans la marmite (moins d'une cuillère à café de sel par jour), plus de marche chaque jour, moins d'alcool.
\item Fais mesurer ta tension au centre de santé au moins une fois par an ; si tu es traité, n'arrête jamais tes médicaments sans avis du médecin.
\end{itemize}
\emph{Appel à l'action :} « Rendez-vous samedi au centre de santé du quartier : mesure gratuite de la tension pour tous. »\\
\textbf{Conclusion :} la lutte contre l'HTA repose sur la prévention (réduction du sel, activité physique), le dépistage régulier et l'observance du traitement : la sensibilisation communautaire est un outil de santé publique essentiel.
\end{exoresolu}

\begin{attention}[Les 5 erreurs qui coûtent des points au Bac]
\begin{enumerate}
\item \textbf{Confondre glycémie et diabète dans la rédaction.} On mesure la \emph{glycémie} (un taux, en g/L) ; le diabète est une \emph{maladie} diagnostiquée à partir de glycémies anormales répétées. Écrire « le glucomètre mesure le diabète » est une faute de vocabulaire sanctionnée.
\item \textbf{Oublier les unités et les valeurs de référence.} Une glycémie sans unité (« 1,4 ») ou une comparaison sans seuil (norme à jeun : 0,70 -- \SI{1,10}{g/L} ; seuil du diabète : \SI{1,26}{g/L} ; HTA : $\geq 14/9$ cmHg) rend l'interprétation non recevable. Toute conclusion exige une comparaison chiffrée à la norme.
\item \textbf{Inverser les rôles de l'insuline et du glucagon.} L'insuline (cellules $\beta$) est \emph{hypoglycémiante} : elle fait baisser la glycémie ; le glucagon (cellules $\alpha$) est \emph{hyperglycémiant}. Erreur classique : attribuer la glycogénolyse à l'insuline.
\item \textbf{Décrire une courbe sans l'expliquer, ou l'expliquer sans citer de valeurs.} Le correcteur attend les deux : relevés chiffrés (pic, temps du retour à la normale) \emph{et} mécanisme hormonal. De même, interpréter une expérience sans comparer au lot témoin invalide la conclusion.
\item \textbf{Diagnostiquer une hypertension sur une mesure unique ou à l'effort.} Une pression élevée après un effort ou un stress est physiologique ; l'HTA exige des mesures $\geq 14/9$ cmHg \emph{au repos} et \emph{répétées}. Oublier ces conditions d'application est l'erreur de raisonnement la plus fréquente.
\end{enumerate}
\end{attention}

\newpage

\coursec{Exercices}

\courssub{Je vérifie mes connaissances}

\begin{exercice}[QCM : glycémie et régulation]
Pour chaque question, choisis la (ou les) bonne(s) réponse(s) en justifiant brièvement ton choix.
\begin{enumerate}
\item La glycémie normale à jeun chez un sujet sain est comprise entre :
\begin{enumerate}
\item \num{0,30} et \SI{0,50}{g/L} ;
\item \num{0,70} et \SI{1,10}{g/L} ;
\item \num{1,50} et \SI{2,00}{g/L} ;
\item \num{2,00} et \SI{3,00}{g/L}.
\end{enumerate}
\item L'hormone hypoglycémiante sécrétée par les cellules $\beta$ des îlots de Langerhans est :
\begin{enumerate}
\item le glucagon ; \item l'adrénaline ; \item l'insuline ; \item le cortisol.
\end{enumerate}
\item Le foie stocke le glucose sous forme de :
\begin{enumerate}
\item cellulose ; \item glycogène ; \item amidon ; \item saccharose.
\end{enumerate}
\item Le glucagon agit principalement en provoquant :
\begin{enumerate}
\item la glycogénogenèse ; \item la glycogénolyse hépatique ; \item l'entrée du glucose dans les cellules musculaires ; \item la glycosurie.
\end{enumerate}
\end{enumerate}
\end{exercice}

\begin{exercice}[Vrai ou faux justifié]
Indique si chaque affirmation est vraie ou fausse, puis justifie en une ou deux phrases.
\begin{enumerate}
\item Le pancréas est à la fois une glande exocrine et une glande endocrine.
\item L'insuline et le glucagon sont des hormones synergiques.
\item La glycosurie (présence de glucose dans les urines) apparaît lorsque la glycémie dépasse environ \SI{1,80}{g/L}.
\item Le diabète de type 1 est dû à une résistance des cellules cibles à l'insuline.
\item Le système nerveux végétatif intervient dans la régulation de la glycémie via l'adrénaline.
\end{enumerate}
\end{exercice}

\begin{exercice}[Définitions]
Définis précisément, en une phrase chacune, les termes suivants : \cle{glycémie} ; \cle{hyperglycémie provoquée} (HGPO) ; \cle{diabète} ; \cle{pression artérielle} ; \cle{hypertension artérielle} ; \cle{barorécepteurs}.
\end{exercice}

\begin{exercice}[Calculs directs]
\begin{enumerate}
\item Un prélèvement sanguin réalisé à jeun chez un patient indique une concentration en glucose de \SI{7,2}{mmol/L}. Sachant que la masse molaire du glucose est $M = \SI{180}{g/mol}$, convertis cette valeur en \SI{}{g/L} et conclus quant à l'état du patient (normoglycémie : \num{0,70} à \SI{1,10}{g/L} ; seuil du diabète à jeun : $\geq \SI{1,26}{g/L}$).
\item Au dispensaire de Bafoussam, on relève chez un adulte une tension de $16/9$ (en cm de Hg). Exprime ces valeurs en mm de Hg, identifie la pression systolique et la pression diastolique, puis calcule la pression différentielle (pression pulsée).
\end{enumerate}
\end{exercice}

\courssub{Je m'entraîne}

\begin{exercice}[Analyse d'une courbe de glycémie après un repas]
Chez un sujet sain à jeun (glycémie initiale \SI{0,90}{g/L}), on mesure la glycémie toutes les 30 minutes après l'ingestion d'un repas riche en glucides :
\begin{center}
\begin{tabularx}{\linewidth}{|l|X|X|X|X|X|X|X|}\hline
\rowcolor{popblueL} Temps (min) & 0 & 30 & 60 & 90 & 120 & 150 & 180 \\ \hline
Glycémie (g/L) & \num{0,90} & \num{1,35} & \num{1,50} & \num{1,20} & \num{0,95} & \num{0,85} & \num{0,90} \\ \hline
\end{tabularx}
\end{center}
\begin{enumerate}
\item Trace la courbe de la glycémie en fonction du temps (échelles au choix, axes légendés avec unités).
\item Décris les deux phases de la courbe et propose une hypothèse explicative pour chacune.
\item Nomme l'hormone responsable du retour à la normale et précise ses trois modes d'action sur le métabolisme du glucose.
\end{enumerate}
\end{exercice}

\begin{exercice}[Schéma fonctionnel de la régulation de la glycémie]
\begin{enumerate}
\item Construis le schéma fonctionnel complet de la régulation de la glycémie faisant apparaître : les deux déviations possibles (hyperglycémie, hypoglycémie), les récepteurs, les glandes endocrines impliquées, les hormones sécrétées, les organes effecteurs (foie, muscles, tissu adipeux) et les rétroactions.
\item Justifie le qualificatif de \og régulation par boucles de rétroaction négative \fg{}.
\item Explique pourquoi, après un effort musculaire intense et prolongé, la glycémie d'un sujet sain ne s'effondre pas.
\end{enumerate}
\end{exercice}

\begin{exercice}[Variations de la pression artérielle au cours d'un 100 m]
Chez un élève de Terminale D, on relève la tension artérielle au repos, immédiatement après un 100 m, puis 10 minutes après l'effort :
\begin{center}
\begin{tabularx}{\linewidth}{|l|X|X|X|}\hline
\rowcolor{popblueL} Situation & Repos & Après l'effort & 10 min après \\ \hline
Tension (cm de Hg) & $12/8$ & $17/9$ & $12{,}5/8$ \\ \hline
Fréquence cardiaque (battements/min) & 70 & 150 & 78 \\ \hline
\end{tabularx}
\end{center}
\begin{enumerate}
\item Calcule la pression différentielle dans chaque situation.
\item Explique l'augmentation de la pression systolique pendant l'effort en t'appuyant sur les paramètres physiologiques (débit cardiaque, fréquence cardiaque, volume d'éjection systolique).
\item Décris les mécanismes nerveux qui ramènent la pression artérielle à sa valeur de repos après l'effort (récepteurs, centres nerveux, nerfs effecteurs).
\end{enumerate}
\end{exercice}

\begin{exercice}[Lecture d'un bilan médical]
Le bilan d'un patient de 52 ans consultant à l'hôpital de district de Nkongsamba donne : glycémie à jeun = \SI{1,85}{g/L} ; glycosurie = positive ; insulinémie à jeun = normale ; tension artérielle = $15/9$ ; surcharge pondérale (IMC = \SI{31}{kg/m^2}).
\begin{enumerate}
\item Analyse chaque paramètre en le comparant aux valeurs de référence.
\item Quel type de diabète ce bilan évoque-t-il ? Justifie ta réponse à partir de l'insulinémie.
\item Propose trois mesures hygiéno-diététiques adaptées au contexte camerounais (alimentation locale, activité physique, suivi médical).
\end{enumerate}
\end{exercice}

\courssub{Je me prépare au Bac}

\begin{exercice}[Type Bac — Les expériences historiques sur le pancréas]
En 1889, von Mering et Minkowski réalisent l'ablation du pancréas chez des chiens (expérience 1). En 1921, Banting et Best ligaturent le canal pancréatique de chiens, ce qui provoque la dégénérescence de la partie exocrine de la glande (expérience 2), puis injectent à des chiens dépancréatés un extrait de pancréas ligaturé (expérience 3). Les glycémies moyennes mesurées sont consignées ci-dessous :
\begin{center}
\begin{tabularx}{\linewidth}{|l|X|X|}\hline
\rowcolor{popblueL} Expérience & Protocole & Glycémie (g/L) \\ \hline
Témoin & Chien sain, à jeun & \num{0,95} \\ \hline
1 & Ablation totale du pancréas & \num{2,60} \\ \hline
2 & Ligature du canal pancréatique & \num{0,98} \\ \hline
3 & Chien dépancréaté + extrait de pancréas ligaturé & \num{1,05} \\ \hline
\end{tabularx}
\end{center}
\begin{enumerate}
\item Analyse les résultats de chaque expérience en les comparant au témoin (calculs des variations à l'appui).
\item Déduis de l'expérience 1 le rôle du pancréas dans la glycémie.
\item Montre, à partir des expériences 2 et 3, que ce rôle est lié à la partie endocrine de la glande et non à la partie exocrine. Nomme la substance active mise en évidence.
\item Complète : on réalise ensuite chez le chien dépancréaté une greffe de pancréas sous la peau, sans connexion nerveuse avec l'organisme ; la glycémie redevient normale. Que prouve cette expérience sur le mode d'action du pancréas ?
\item Construis le schéma fonctionnel de la régulation de la glycémie en y intégrant le rôle du foie et les deux hormones pancréatiques antagonistes.
\end{enumerate}
\end{exercice}

\begin{exercice}[Type Bac — Hyperglycémie provoquée et typologie des diabètes]
Au CHU de Yaoundé, on soumet trois sujets (A : sain ; B et C : diabétiques) à une hyperglycémie provoquée par voie orale (HGPO) : ingestion de \SI{75}{g} de glucose, puis dosage de la glycémie et de l'insulinémie pendant 3 heures. Résultats :
\begin{center}
\begin{tabularx}{\linewidth}{|l|X|X|X|X|}\hline
\rowcolor{popblueL} Temps (min) & 0 & 60 & 120 & 180 \\ \hline
Glycémie A (g/L) & \num{0,90} & \num{1,40} & \num{1,10} & \num{0,90} \\ \hline
Glycémie B (g/L) & \num{1,60} & \num{2,80} & \num{2,60} & \num{2,30} \\ \hline
Glycémie C (g/L) & \num{1,50} & \num{2,50} & \num{2,40} & \num{2,10} \\ \hline
Insulinémie A ($\mu$UI/mL) & 8 & 65 & 30 & 10 \\ \hline
Insulinémie B ($\mu$UI/mL) & 2 & 4 & 3 & 2 \\ \hline
Insulinémie C ($\mu$UI/mL) & 15 & 90 & 70 & 40 \\ \hline
\end{tabularx}
\end{center}
\begin{enumerate}
\item Trace sur un même graphique les trois courbes de glycémie (axes légendés, unités, échelle adaptée).
\item Analyse la courbe du sujet A et explique le retour à la normoglycémie en t'appuyant sur sa courbe d'insulinémie.
\item Compare les insulinémies des sujets B et C à celle du sujet A. Attribue à chaque sujet un type de diabète (type 1 ou type 2) en justifiant rigoureusement.
\item Explique le mécanisme cellulaire à l'origine de l'hyperglycémie chronique dans chacun des deux types.
\item Propose pour chaque sujet diabétique le traitement approprié (insulinothérapie ou antidiabétiques oraux associés à un régime) et justifie ton choix.
\end{enumerate}
\end{exercice}

\begin{exercice}[Type Bac — Hypertension artérielle : diagnostic et sensibilisation]
Monsieur E., 48 ans, commerçant à Douala, consulte pour des maux de tête fréquents, des vertiges et des saignements de nez. Le médecin relève sa tension artérielle au repos, le matin, pendant une semaine :
\begin{center}
\begin{tabularx}{\linewidth}{|l|X|X|X|X|X|X|X|}\hline
\rowcolor{popblueL} Jour & L & M & M & J & V & S & D \\ \hline
Systolique (cm de Hg) & \num{16,5} & 17 & \num{16,0} & \num{17,5} & \num{16,8} & 17 & \num{16,2} \\ \hline
Diastolique (cm de Hg) & \num{10,5} & 11 & 10 & \num{11,5} & \num{10,8} & 11 & \num{10,4} \\ \hline
\end{tabularx}
\end{center}
L'interrogatoire révèle : consommation quotidienne de plats très salés (sel de cuisine et cubes aromatiques), consommation régulière d'alcool, sédentarité, antécédents d'hypertension chez le père.
\begin{enumerate}
\item Calcule les tensions systolique et diastolique moyennes de la semaine. Rappelle les valeurs normales ($12/8$ cm de Hg ; HTA si systolique $\geq 14$ cm de Hg ou diastolique $\geq 9$ cm de Hg au repos, de façon répétée) et conclus au diagnostic.
\item Identifie dans le mode de vie de M. E. quatre facteurs de risque d'hypertension artérielle et explique le mécanisme par lequel l'excès de sel élève la pression artérielle (rétention d'eau, volémie).
\item Explique le rôle des barorécepteurs de l'aorte et des sinus carotidiens et du centre vasomoteur bulbaire dans la régulation nerveuse de la pression artérielle ; précise pourquoi cette régulation devient insuffisante en cas d'HTA installée.
\item Cite trois complications graves de l'HTA non traitée (accident vasculaire cérébral, infarctus du myocarde, insuffisance rénale) et explique brièvement le lien avec l'élévation chronique de la pression.
\item Rédige un court message de sensibilisation (10 à 15 lignes) destiné aux élèves de ton lycée, présentant l'HTA comme un \og tueur silencieux \fg{} et proposant des mesures hygiéno-diététiques concrètes adaptées au contexte camerounais.
\end{enumerate}
\end{exercice}

\newpage

\coursec{Corrigés des exercices}

\begin{corrige}[Exercice 1 — QCM : glycémie et régulation]
\begin{enumerate}
\item \textbf{Bonne réponse : b.} La glycémie normale à jeun est comprise entre \num{0,70} et \SI{1,10}{g/L} (soit environ \SI{1}{g/L} en moyenne). Les valeurs a sont hypoglycémiques ; c et d correspondent à des hyperglycémies franches (diabète).
\item \textbf{Bonne réponse : c.} L'\cle{insuline} est sécrétée par les cellules $\beta$ des îlots de Langerhans (pancréas endocrine) ; c'est la seule hormone \emph{hypoglycémiante}. Le glucagon (cellules $\alpha$), l'adrénaline et le cortisol sont hyperglycémiants.
\item \textbf{Bonne réponse : b.} Le foie (et les muscles) stockent le glucose sous forme de \cle{glycogène}, un polymère de glucose animal. L'amidon est la forme de réserve végétale ; la cellulose est structurale ; le saccharose est un diholoside alimentaire.
\item \textbf{Bonne réponse : b.} Le glucagon, hormone hyperglycémiante, déclenche la \cle{glycogénolyse} hépatique (dégradation du glycogène du foie en glucose libéré dans le sang) et stimule la néoglucogenèse. La glycogénogenèse (a) et l'entrée du glucose dans les cellules (c) sont des effets de l'\emph{insuline} ; la glycosurie (d) est un symptôme, pas une action hormonale.
\end{enumerate}
\end{corrige}

\begin{corrige}[Exercice 2 — Vrai ou faux justifié]
\begin{enumerate}
\item \textbf{Vrai.} Le pancréas est une glande \emph{mixte} : sa partie exocrine (acini) sécrète le suc pancréatique (enzymes digestives) dans le duodénum via le canal de Wirsung ; sa partie endocrine (îlots de Langerhans) déverse insuline et glucagon dans le sang.
\item \textbf{Faux.} Insuline et glucagon sont des hormones \emph{antagonistes} : l'insuline fait baisser la glycémie (hypoglycémiante), le glucagon la fait monter (hyperglycémiant). Des hormones synergiques auraient des effets convergents (ex. glucagon et adrénaline).
\item \textbf{Vrai.} Le rein réabsorbe totalement le glucose filtré jusqu'à un seuil appelé \cle{seuil rénal du glucose}, voisin de \SI{1,80}{g/L}. Au-delà, les transporteurs tubulaires sont saturés et le glucose passe dans les urines (glycosurie).
\item \textbf{Faux.} Le diabète de type 1 résulte de la \emph{destruction auto-immune des cellules $\beta$}, donc d'un \emph{déficit absolu en insuline}. La résistance des cellules cibles à l'insuline (avec insulinémie normale ou élevée) caractérise le diabète de type 2.
\item \textbf{Vrai.} En cas d'hypoglycémie, le système nerveux végétatif (orthosympathique) stimule la médullo-surrénale qui libère de l'\cle{adrénaline}, hormone hyperglycémiante qui accélère la glycogénolyse hépatique : c'est la composante nerveuse (via une glande) de la régulation.
\end{enumerate}
\end{corrige}

\begin{corrige}[Exercice 3 — Définitions]
\begin{itemize}
\item \cle{Glycémie} : taux (concentration) de glucose dissous dans le sang, normalement compris entre \num{0,70} et \SI{1,10}{g/L} à jeun.
\item \cle{Hyperglycémie provoquée} (HGPO) : épreuve diagnostique consistant à faire ingérer à jeun une charge de glucose (\SI{75}{g}) puis à suivre la glycémie pendant 2 à 3 heures afin de tester l'aptitude de l'organisme à réguler sa glycémie.
\item \cle{Diabète} : maladie métabolique chronique caractérisée par une hyperglycémie permanente (glycémie à jeun $\geq \SI{1,26}{g/L}$ à deux reprises), due à un défaut de sécrétion ou d'action de l'insuline.
\item \cle{Pression artérielle} : force exercée par le sang sur la paroi des artères, caractérisée par une pression maximale (systolique, $\approx 12$ cm de Hg) et une pression minimale (diastolique, $\approx 8$ cm de Hg).
\item \cle{Hypertension artérielle} (HTA) : élévation permanente et répétée de la pression artérielle au repos (systolique $\geq 14$ cm de Hg et/ou diastolique $\geq 9$ cm de Hg).
\item \cle{Barorécepteurs} : récepteurs sensibles à l'étirement de la paroi artérielle, situés dans la crosse de l'aorte et les sinus carotidiens ; ils informent le centre vasomoteur bulbaire des variations de la pression artérielle.
\end{itemize}
\end{corrige}

\begin{corrige}[Exercice 4 — Calculs directs]
\begin{enumerate}
\item Conversion : $C_{\text{massique}} = C_{\text{molaire}} \times M$.
\[ C = \num{7,2}\ \text{mmol/L} \times 180\ \text{g/mol} = \num{7,2}\times 10^{-3}\ \text{mol/L} \times 180\ \text{g/mol} = \SI{1,296}{g/L} \approx \SI{1,30}{g/L}. \]
Conclusion : $\num{1,30} > \num{1,26}$ g/L : la glycémie à jeun dépasse le seuil du diabète. Le patient est en \textbf{hyperglycémie} ; un diabète est suspecté et doit être confirmé par un second dosage.
\item Conversion : $16/9$ cm de Hg $= 160/90$ mm de Hg.
\begin{itemize}
\item Pression \textbf{systolique} (maxima) : \SI{160}{mm} de Hg ;
\item Pression \textbf{diastolique} (minima) : \SI{90}{mm} de Hg ;
\item Pression différentielle (pulsée) : $160 - 90 = \SI{70}{mm} \text{ de Hg}$ (soit \SI{7}{cm} de Hg).
\end{itemize}
Ces valeurs sont élevées (norme : $120/80$ mm de Hg) : elles évoquent une hypertension artérielle.
\end{enumerate}
\end{corrige}

\begin{corrige}[Exercice 5 — Analyse d'une courbe de glycémie après un repas]
\begin{enumerate}
\item Courbe (échelles : 1 cm pour 30 min ; 1 cm pour \SI{0,20}{g/L}) :
\begin{popfigure}
\begin{center}
\begin{tikzpicture}
\begin{axis}[width=0.85\linewidth,height=6cm,xlabel={Temps (min)},ylabel={Glycémie (g/L)},xmin=0,xmax=190,ymin=0.7,ymax=1.7,xtick={0,30,60,90,120,150,180},ytick={0.7,0.9,1.1,1.3,1.5,1.7},grid=both,grid style={popline!40},legend style={font=\small}]
\addplot[popcurve,mark=*,mark size=1.5pt] coordinates {(0,0.90) (30,1.35) (60,1.50) (90,1.20) (120,0.95) (150,0.85) (180,0.90)};
\addplot[domain=0:190,dashed,popmuted] {1.10};
\legend{Glycémie du sujet, Limite supérieure de la norme}
\end{axis}
\end{tikzpicture}
\end{center}
\legende{Évolution de la glycémie d'un sujet sain après un repas glucidique.}
\end{popfigure}
\item \textbf{Phase 1 (0 à 60 min) : hyperglycémie post-prandiale.} La glycémie passe de \num{0,90} à \SI{1,50}{g/L} ($+\SI{0,60}{g/L}$, soit $+67\,\%$). Hypothèse : l'absorption intestinale massive du glucose issu de la digestion des glucides fait transiter le glucose dans le sang plus vite qu'il n'est capté par les cellules.
\textbf{Phase 2 (60 à 180 min) : retour progressif à la normale.} La glycémie redescend à \SI{0,90}{g/L}, avec un léger dépassement vers le bas à 150 min (\SI{0,85}{g/L}). Hypothèse : une hormone hypoglycémiante accélère la captation et le stockage du glucose, puis la régulation par rétroaction stabilise la glycémie.
\item L'hormone responsable est l'\cle{insuline}, sécrétée par les cellules $\beta$ des îlots de Langerhans en réponse à l'hyperglycémie. Ses trois modes d'action :
\begin{itemize}
\item elle favorise la \textbf{pénétration et l'utilisation du glucose} par les cellules (muscles, tissu adipeux) ;
\item elle stimule la \textbf{glycogénogenèse} (stockage du glucose en glycogène dans le foie et les muscles) ;
\item elle \textbf{inhibe la glycogénolyse} et la néoglucogenèse hépatiques (et favorise la lipogenèse à partir du glucose excédentaire).
\end{itemize}
\end{enumerate}
\end{corrige}

\begin{corrige}[Exercice 6 — Schéma fonctionnel de la régulation de la glycémie]
\begin{enumerate}
\item Schéma fonctionnel :
\begin{popfigure}
\begin{center}
\begin{tikzpicture}[font=\small, node distance=0.55cm and 1.1cm,
boite/.style={draw=popblue, fill=popblueL, rounded corners, align=center, inner sep=3pt},
horm/.style={draw=poporange, fill=poporangeL, rounded corners, align=center, inner sep=3pt},
eff/.style={draw=popgreen, fill=popgreenL, rounded corners, align=center, inner sep=3pt},
fleche/.style={-{Stealth[length=2.5mm]}, thick, popdark}]
\node[boite, align=center] (hyper){Hyperglycémie\\ $> \SI{1,10}{g/L}$};
\node[boite, below=1.9cm of hyper, align=center] (hypo){Hypoglycémie\\ $< \SI{0,70}{g/L}$};
\node[boite, right=of hyper, align=center] (beta){Cellules $\beta$ des îlots\\ de Langerhans (récepteurs)};
\node[boite, right=of hypo, align=center] (alpha){Cellules $\alpha$ des îlots\\ + centre nerveux (SNA)};
\node[horm, right=of beta] (ins) {Insuline};
\node[horm, right=of alpha] (glu) {Glucagon (+ adrénaline)};
\node[eff, right=of ins, align=center] (eff1){Foie, muscles, tissu adipeux :\\ captation, glycogénogenèse,\\ inhibition glycogénolyse};
\node[eff, right=of glu, align=center] (eff2){Foie :\\ glycogénolyse,\\ néoglucogenèse};
\draw[fleche] (hyper) -- (beta);
\draw[fleche] (hypo) -- (alpha);
\draw[fleche] (beta) -- (ins);
\draw[fleche] (alpha) -- (glu);
\draw[fleche] (ins) -- (eff1);
\draw[fleche] (glu) -- (eff2);
\draw[fleche, popgreen] (eff1.north) .. controls +(0,1.1) and +(1,0.6) .. node[above, align=center]{glycémie $\searrow$\\ (rétroaction $-$)} (hyper.east);
\draw[fleche, popgreen] (eff2.south) .. controls +(0,-1.1) and +(1,-0.6) .. node[below, align=center]{glycémie $\nearrow$\\ (rétroaction $-$)} (hypo.east);
\end{tikzpicture}
\end{center}
\legende{Schéma fonctionnel de la régulation de la glycémie : deux boucles antagonistes à rétroaction négative.}
\end{popfigure}
\item On parle de \og boucles de rétroaction négative \fg{} car la réponse hormonale \emph{s'oppose à la variation qui l'a déclenchée} : l'hyperglycémie provoque la sécrétion d'insuline qui fait baisser la glycémie (donc supprime le stimulus), et l'hypoglycémie provoque la sécrétion de glucagon qui la fait remonter. La variable réglée (glycémie) oscille ainsi autour d'une valeur de consigne proche de \SI{1}{g/L}.
\item Pendant un effort intense et prolongé, les muscles consomment beaucoup de glucose, mais la glycémie ne s'effondre pas car :
\begin{itemize}
\item la baisse naissante de la glycémie stimule les cellules $\alpha$ et le système orthosympathique : sécrétion de \textbf{glucagon} et d'\textbf{adrénaline} ;
\item ces hormones déclenchent la \textbf{glycogénolyse hépatique} (le foie libère son glucose de réserve) puis la \textbf{néoglucogenèse} (fabrication de glucose à partir des acides aminés et du lactate) ;
\item la mobilisation des graisses (lipolyse) fournit en outre un carburant de substitution aux muscles.
\end{itemize}
Le foie compense donc en permanence la consommation musculaire : la glycémie reste dans les limites compatibles avec le fonctionnement du cerveau.
\end{enumerate}
\end{corrige}

\begin{corrige}[Exercice 7 — Variations de la pression artérielle au cours d'un 100 m]
\begin{enumerate}
\item Pression différentielle = systolique $-$ diastolique :
\begin{itemize}
\item Repos : $12 - 8 = \SI{4}{cm}$ de Hg ;
\item Après l'effort : $17 - 9 = \SI{8}{cm}$ de Hg ;
\item 10 min après : $12{,}5 - 8 = \SI{4,5}{cm}$ de Hg.
\end{itemize}
\item Pendant l'effort, les muscles en activité réclament un apport accru de dioxygène et de glucose. L'organisme y répond par :
\begin{itemize}
\item une augmentation de la \textbf{fréquence cardiaque} ($70 \to 150$ battements/min, soit $+114\,\%$) sous l'action du nerf orthosympathique (cardio-accélérateur) et de l'adrénaline ;
\item une augmentation du \textbf{volume d'éjection systolique} (VES) grâce à la contraction ventriculaire plus vigoureuse et au retour veineux accru.
\end{itemize}
Or le débit cardiaque est $Q = f_c \times \text{VES}$ : il augmente fortement, ce qui élève la pression systolique ($12 \to 17$ cm de Hg). La diastolique varie peu ($8 \to 9$) car la vasodilatation des artérioles musculaires diminue les résistances périphériques.
\item Après l'effort, l'hyperpression distend la paroi des artères :
\begin{itemize}
\item les \textbf{barorécepteurs} de la crosse aortique et des sinus carotidiens, étirés, augmentent leur fréquence d'influx nerveux ;
\item ces influx parviennent au \textbf{centre vasomoteur bulbaire} (moelle allongée) via les nerfs sensitifs (glossopharyngien IX et vague X) ;
\item le centre répond en activant le nerf \textbf{parasympathique (vague X, cardio-modérateur)} et en inhibant l'orthosympathique : la fréquence cardiaque chute ($150 \to 78$), les artérioles se contractent modérément ;
\item la pression artérielle revient à sa valeur de repos : c'est une \textbf{boucle de rétroaction négative}.
\end{itemize}
\end{enumerate}
\end{corrige}

\begin{corrige}[Exercice 8 — Lecture d'un bilan médical]
\begin{enumerate}
\item Analyse des paramètres :
\begin{itemize}
\item \textbf{Glycémie à jeun : \SI{1,85}{g/L}} $\gg$ \SI{1,26}{g/L} (seuil du diabète) : hyperglycémie chronique confirmée ;
\item \textbf{Glycosurie positive} : cohérente, car $\num{1,85} > \SI{1,80}{g/L}$ (seuil rénal du glucose) ;
\item \textbf{Insulinémie à jeun normale} : le pancréas sécrète correctement l'insuline ; le défaut se situe donc en aval, au niveau des cellules cibles ;
\item \textbf{Tension $15/9$ cm de Hg} : au seuil de l'HTA (systolique $\geq 14$, diastolique $\geq 9$) : hypertension associée, fréquente dans le diabète de type 2 ;
\item \textbf{IMC $= \SI{31}{kg/m^2}$} : obésité (IMC $\geq 30$), facteur majeur d'insulinorésistance.
\end{itemize}
\item Ce bilan évoque un \textbf{diabète de type 2} (diabète gras, de l'adulte, insulinorésistant). Justification : l'insulinémie est \emph{normale} alors que la glycémie est élevée — l'insuline est produite mais inefficace, ce qui traduit une \textbf{résistance des cellules cibles} (défaut des récepteurs à insuline), aggravée par l'obésité. Dans un diabète de type 1, l'insulinémie aurait été effondrée.
\item Mesures hygiéno-diététiques adaptées au contexte camerounais :
\begin{itemize}
\item \textbf{Alimentation} : réduire les féculents raffinés et les boissons sucrées ; privilégier les légumes locaux (folong, ndolé, gombo), les légumineuses (niébé, arachides non sucrées) et des portions modérées de tubercules (manioc, igname, macabo) ; limiter huile rouge et fritures ;
\item \textbf{Activité physique} : marche rapide quotidienne d'au moins 30 min (trajets à pied, travaux champêtres réguliers) pour améliorer la sensibilité à l'insuline et réduire le poids ;
\item \textbf{Suivi médical} : consultation régulière au centre de santé, contrôle trimestriel de la glycémie à jeun et de la tension, observance du traitement par antidiabétiques oraux s'il est prescrit.
\end{itemize}
\end{enumerate}
\end{corrige}

\begin{corrige}[Exercice 9 — Type Bac : les expériences historiques sur le pancréas]
\begin{enumerate}
\item Analyse comparée au témoin (\SI{0,95}{g/L}) :
\begin{itemize}
\item \textbf{Expérience 1} : $\num{2,60} - \num{0,95} = +\SI{1,65}{g/L}$, soit une glycémie multipliée par $\frac{2,60}{0,95} \approx 2{,}7$ ($+174\,\%$) : hyperglycémie majeure après pancréatectomie ;
\item \textbf{Expérience 2} : $\num{0,98} - \num{0,95} = +\SI{0,03}{g/L}$ ($\approx +3\,\%$) : glycémie quasi normale malgré la destruction de la partie exocrine ;
\item \textbf{Expérience 3} : $\num{1,05} - \num{0,95} = +\SI{0,10}{g/L}$ ($\approx +11\,\%$) : l'injection d'extrait ramène la glycémie du chien dépancréaté de \num{2,60} à \SI{1,05}{g/L}, soit une baisse de $\frac{2,60-1,05}{2,60} \approx 60\,\%$.
\end{itemize}
\item L'ablation totale du pancréas provoque une hyperglycémie permanente (diabète expérimental) : le pancréas exerce donc un \textbf{rôle hypoglycémiant} ; il maintient la glycémie à sa valeur normale.
\item L'expérience 2 montre que la destruction de la partie \emph{exocrine} (dégénérescence des acini après ligature du canal) ne modifie pas la glycémie : le rôle glycorégulateur n'appartient pas à la fonction digestive. L'expérience 3 montre qu'un extrait de pancréas \emph{ne contenant plus que la partie endocrine} (îlots de Langerhans survivants) corrige l'hyperglycémie du chien dépancréaté. Conclusion : le facteur hypoglycémiant est sécrété par les \textbf{îlots de Langerhans} (partie endocrine) ; c'est l'\cle{insuline}, isolée par Banting et Best en 1921.
\item La greffe sous-cutanée, dépourvue de toute connexion nerveuse, normalise la glycémie : le pancréas agit donc par \textbf{voie sanguine (hormonale)} et non par voie nerveuse. Il libère dans le sang une substance (hormone) transportée jusqu'aux organes cibles.
\item Schéma fonctionnel :
\begin{popfigure}
\begin{center}
\begin{tikzpicture}[font=\small, node distance=0.5cm and 0.9cm,
boite/.style={draw=popblue, fill=popblueL, rounded corners, align=center, inner sep=3pt},
horm/.style={draw=poporange, fill=poporangeL, rounded corners, align=center, inner sep=3pt},
eff/.style={draw=popgreen, fill=popgreenL, rounded corners, align=center, inner sep=3pt},
fleche/.style={-{Stealth[length=2.5mm]}, thick, popdark}]
\node[boite] (hyper) {Hyperglycémie};
\node[boite, below=1.7cm of hyper] (hypo) {Hypoglycémie};
\node[boite, right=of hyper, align=center] (beta){Cellules $\beta$\\ (îlots de Langerhans)};
\node[boite, right=of hypo, align=center] (alpha){Cellules $\alpha$\\ (îlots de Langerhans)};
\node[horm, right=of beta] (ins) {Insuline};
\node[horm, right=of alpha] (glu) {Glucagon};
\node[eff, right=of ins, align=center] (foie1){Foie : glycogénogenèse,\\ $\downarrow$ glycogénolyse\\ Muscles, tissu adipeux : captation};
\node[eff, right=of glu, align=center] (foie2){Foie :\\ glycogénolyse,\\ néoglucogenèse};
\draw[fleche] (hyper) -- (beta); \draw[fleche] (hypo) -- (alpha);
\draw[fleche] (beta) -- (ins); \draw[fleche] (alpha) -- (glu);
\draw[fleche] (ins) -- (foie1); \draw[fleche] (glu) -- (foie2);
\draw[fleche, popgreen] (foie1.north) .. controls +(0,1.0) and +(1,0.5) .. node[above]{glycémie $\searrow$ ($-$)} (hyper.east);
\draw[fleche, popgreen] (foie2.south) .. controls +(0,-1.0) and +(1,-0.5) .. node[below]{glycémie $\nearrow$ ($-$)} (hypo.east);
\end{tikzpicture}
\end{center}
\legende{Régulation de la glycémie : rôles antagonistes de l'insuline et du glucagon, le foie comme organe effecteur central.}
\end{popfigure}
\end{enumerate}
\textbf{Barème indicatif (sur 20)} : analyse chiffrée des trois expériences (6) ; rôle du pancréas (2) ; démonstration endocrine + insuline (4) ; preuve de la voie hormonale (3) ; schéma fonctionnel complet et légendé (5).

\textbf{Points de vigilance} : toujours comparer chaque résultat \emph{au témoin} avec un calcul d'écart ; distinguer clairement pancréas exocrine (digestion) et endocrine (îlots) ; la question 4 exige la conclusion \og action humorale \fg{} justifiée par l'\emph{absence de connexion nerveuse} ; dans le schéma, les flèches de rétroaction doivent être annotées ($-$).
\end{corrige}

\begin{corrige}[Exercice 10 — Type Bac : HGPO et typologie des diabètes]
\begin{enumerate}
\item Graphique (échelles : 1 cm pour 30 min ; 1 cm pour \SI{0,40}{g/L}) :
\begin{popfigure}
\begin{center}
\begin{tikzpicture}
\begin{axis}[width=0.85\linewidth,height=6.5cm,xlabel={Temps (min)},ylabel={Glycémie (g/L)},xmin=0,xmax=190,ymin=0.6,ymax=3.1,xtick={0,60,120,180},grid=both,grid style={popline!40},legend style={font=\small,at={(0.98,0.98)},anchor=north east}]
\addplot[popcurve,mark=*,mark size=1.5pt] coordinates {(0,0.90) (60,1.40) (120,1.10) (180,0.90)};
\addplot[poporange,very thick,mark=square*,mark size=1.5pt] coordinates {(0,1.60) (60,2.80) (120,2.60) (180,2.30)};
\addplot[poppurple,very thick,mark=triangle*,mark size=2pt] coordinates {(0,1.50) (60,2.50) (120,2.40) (180,2.10)};
\addplot[domain=0:190,dashed,popmuted] {1.80};
\legend{Sujet A (sain), Sujet B, Sujet C, Seuil rénal}
\end{axis}
\end{tikzpicture}
\end{center}
\legende{Courbes d'hyperglycémie provoquée par voie orale (HGPO, \SI{75}{g} de glucose) chez trois sujets.}
\end{popfigure}
\item \textbf{Sujet A} : glycémie à jeun normale (\SI{0,90}{g/L}), pic modéré à 60 min (\SI{1,40}{g/L}, sous le seuil rénal de \SI{1,80}{g/L}), retour à la normale à 120--180 min. Son insulinémie suit la même dynamique : elle passe de 8 à \SI{65}{\mu UI/mL} à 60 min (multipliée par $\approx 8$) puis redescend (30 puis 10). Interprétation : l'hyperglycémie provoquée stimule les cellules $\beta$, qui libèrent massivement de l'insuline ; celle-ci accélère la captation et le stockage du glucose (glycogénogenèse), ramenant la glycémie à la normale en 2 h ; la baisse de la glycémie coupe ensuite la sécrétion d'insuline (rétroaction négative).
\item Comparaison des insulinémies :
\begin{itemize}
\item \textbf{Sujet B} : insulinémie effondrée (2 à \SI{4}{\mu UI/mL}, soit 4 à 16 fois moins que A) et quasi insensible à la charge glucidique $\Rightarrow$ \textbf{diabète de type 1} : déficit absolu en insuline par destruction des cellules $\beta$ ;
\item \textbf{Sujet C} : insulinémie \emph{supérieure} à celle de A (15 à \SI{90}{\mu UI/mL}) malgré une hyperglycémie persistante $\Rightarrow$ \textbf{diabète de type 2} : l'insuline est sécrétée (voire en excès compensateur) mais inefficace, car les cellules cibles y sont résistantes.
\end{itemize}
\item Mécanismes cellulaires :
\begin{itemize}
\item \textbf{Type 1 (B)} : réaction auto-immune détruisant les cellules $\beta$ des îlots de Langerhans $\Rightarrow$ absence d'insuline $\Rightarrow$ le glucose ne pénètre plus dans les cellules et n'est plus stocké : il s'accumule dans le sang ;
\item \textbf{Type 2 (C)} : diminution du nombre ou de la sensibilité des \textbf{récepteurs membranaires à l'insuline} (insulinorésistance, favorisée par l'obésité) $\Rightarrow$ malgré une insulinémie élevée, la captation cellulaire du glucose est insuffisante et le foie continue de libérer du glucose.
\end{itemize}
\item Traitements :
\begin{itemize}
\item \textbf{Sujet B (type 1)} : \textbf{insulinothérapie} (injections quotidiennes d'insuline), seul moyen de suppléer le déficit absolu ; l'insuline ne peut être prise par voie orale car c'est une protéine digérée ;
\item \textbf{Sujet C (type 2)} : \textbf{antidiabétiques oraux} (qui améliorent la sensibilité à l'insuline ou stimulent la sécrétion résiduelle) associés à un \textbf{régime hypoglucidique, à la perte de poids et à l'exercice physique} ; l'insulinothérapie n'est envisagée qu'en cas d'échec.
\end{itemize}
\end{enumerate}
\textbf{Barème indicatif (sur 20)} : graphique correct (axes, unités, échelle, 3 courbes) (4) ; analyse couplée glycémie/insulinémie de A (4) ; comparaison et attribution des types (4) ; mécanismes cellulaires (4) ; traitements justifiés (4).

\textbf{Points de vigilance} : ne jamais écrire \og le sujet C manque d'insuline \fg{} (son insulinémie est élevée !) ; citer le seuil rénal pour expliquer la glycosurie ; justifier le choix thérapeutique par le \emph{mécanisme} (déficit de sécrétion vs défaut d'action) ; rappeler que l'insuline est une protéine non administrable par voie orale.
\end{corrige}

\begin{corrige}[Exercice 11 — Type Bac : hypertension artérielle, diagnostic et sensibilisation]
\begin{enumerate}
\item Moyennes hebdomadaires :
\[ \bar{P}_{\text{sys}} = \frac{16{,}5 + 17 + 16{,}0 + 17{,}5 + 16{,}8 + 17 + 16{,}2}{7} = \frac{117{,}0}{7} \approx \SI{16,7}{cm} \text{ de Hg} \]
\[ \bar{P}_{\text{dia}} = \frac{10{,}5 + 11 + 10 + 11{,}5 + 10{,}8 + 11 + 10{,}4}{7} = \frac{75{,}2}{7} \approx \SI{10,7}{cm} \text{ de Hg} \]
Valeurs normales au repos : $12/8$ cm de Hg ; HTA si systolique $\geq 14$ cm de Hg ou diastolique $\geq 9$ cm de Hg, de façon répétée. Ici, \emph{toutes} les mesures de la semaine dépassent ces seuils ($\num{16,7} \geq 14$ et $\num{10,7} \geq 9$) : le diagnostic d'\textbf{hypertension artérielle permanente} est confirmé. Les symptômes (céphalées, vertiges, épistaxis) sont compatibles.
\item Quatre facteurs de risque présents chez M. E. :
\begin{itemize}
\item consommation excessive de \textbf{sel} (sel de cuisine + cubes aromatiques) ;
\item consommation régulière d'\textbf{alcool} ;
\item \textbf{sédentarité} (absence d'activité physique) ;
\item \textbf{hérédité} (antécédent paternel d'HTA).
\end{itemize}
Mécanisme du sel : l'excès de NaCl augmente l'osmolarité du plasma, ce qui provoque une \textbf{rétention d'eau} (l'eau suit le sodium) et stimule la soif ; la \textbf{volémie} (volume sanguin) augmente, donc le débit cardiaque et la pression artérielle s'élèvent. À terme, le sel favorise aussi la contraction et l'épaississement de la paroi des artérioles (augmentation des résistances périphériques).
\item Régulation nerveuse : lorsque la pression s'élève, les \textbf{barorécepteurs} de la crosse aortique et des sinus carotidiens, étirés par la distension de la paroi, émettent davantage d'influx vers le \textbf{centre vasomoteur bulbaire}. Celui-ci active le \textbf{nerf vague (parasympathique)}, cardio-modérateur, et inhibe l'orthosympathique : fréquence cardiaque et contractilité diminuent, les artérioles se dilatent $\Rightarrow$ la pression redescend (rétroaction négative). En cas d'HTA installée, cette régulation devient insuffisante car les barorécepteurs \emph{s'adaptent} au nouveau niveau de pression (leur seuil de déclenchement est \og recalé \fg{} vers le haut) et les artérioles, épaissies et rigidifiées, répondent mal aux ordres vasodilatateurs : l'hyperpression devient permanente.
\item Trois complications de l'HTA non traitée :
\begin{itemize}
\item \textbf{Accident vasculaire cérébral (AVC)} : la pression chronique fragilise les artères cérébrales qui peuvent se rompre (hémorragie) ou se boucher sur une plaque d'athérome (infarctus cérébral) ;
\item \textbf{Infarctus du myocarde} : l'HTA accélère l'athérosclérose des artères coronaires et épuise le cœur qui travaille contre une pression élevée (hypertrophie puis insuffisance cardiaque) ;
\item \textbf{Insuffisance rénale} : les artérioles rénales, soumises en permanence à une pression excessive, s'épaississent et se sclérosent (néphroangiosclérose) ; la filtration rénale devient défaillante.
\end{itemize}
\item Message de sensibilisation (exemple) :

\og \emph{L'hypertension artérielle est un tueur silencieux : elle ne fait pas mal, elle ne se voit pas, mais elle détruit peu à peu ton cœur, ton cerveau et tes reins. Au Cameroun, près d'un adulte sur trois est hypertendu, souvent sans le savoir. Toi, élève, tu peux agir dès maintenant : réduis le sel dans tes plats, méfie-toi des cubes aromatiques utilisés en excès, limite les fritures et les boissons sucrées, bouge au moins 30 minutes par jour — marche, sport, danse. Dis non à l'alcool et au tabac. Encourage tes parents à faire mesurer leur tension au centre de santé le plus proche : c'est rapide, indolore et cela peut sauver une vie. Une tension normale, c'est environ 12/8 ; au-delà de 14/9 mesuré plusieurs fois au repos, il faut consulter. Prévenir l'hypertension, c'est protéger ta famille et ton avenir.} \fg{}
\end{enumerate}
\textbf{Barème indicatif (sur 20)} : calculs des moyennes et conclusion diagnostique (4) ; facteurs de risque et mécanisme du sel (4) ; régulation baroréceptrice et recalage (4) ; complications expliquées (4) ; message de sensibilisation structuré et adapté (4).

\textbf{Points de vigilance} : le diagnostic d'HTA exige des mesures \emph{répétées au repos} (une seule mesure ne suffit pas) ; bien distinguer systolique/diastolique et citer les seuils ; pour le sel, enchaîner logiquement : NaCl $\to$ rétention d'eau $\to$ volémie $\uparrow$ $\to$ pression $\uparrow$ ; le message de sensibilisation doit comporter des mesures \emph{concrètes et locales} (cubes aromatiques, marche, centres de santé).
\end{corrige}

\newpage

\coursec{Fiche bilan}

\begin{popfigure}
\begin{tikzpicture}[mindmap, grow cyclic, every node/.style={concept}, text width=2.6cm, align=center,
  root concept/.append style={concept color=popblue, font=\bfseries, text width=3.2cm},
  level 1/.append style={level distance=4.6cm, sibling angle=60, font=\small},
  level 2/.append style={level distance=2.8cm, sibling angle=45, font=\scriptsize, text width=2.2cm}]
\node[root concept]{Glycémie et pression artérielle}
  child[concept color=poporange]{ node {Glycémie : mesure et variations}
    child[concept color=poporangeL]{ node {Valeur normale à jeun : \num{0,7} à \SI{1,1}{g/L}} }
    child[concept color=poporangeL]{ node {Hyperglycémie post-prandiale transitoire} }
  }
  child[concept color=popgreen]{ node {Régulation de la glycémie}
    child[concept color=popgreenL]{ node {Insuline (cellules $\beta$) : hypoglycémiante} }
    child[concept color=popgreenL]{ node {Glucagon (cellules $\alpha$) : hyperglycémiant} }
    child[concept color=popgreenL]{ node {Foie : stockage et libération du glucose} }
  }
  child[concept color=poppurple]{ node {Le diabète}
    child[concept color=poppurpleL]{ node {Type 1 : insulinodépendant (destruction des cellules $\beta$)} }
    child[concept color=poppurpleL]{ node {Type 2 : insulinorésistance (gras, sédentarité)} }
    child[concept color=poppurpleL]{ node {Signes : polyurie, polydipsie, polyphagie} }
  }
  child[concept color=popblue]{ node {Pression artérielle : mesure}
    child[concept color=popblueL]{ node {Sphygmomanomètre ; maxima / minima} }
    child[concept color=popblueL]{ node {Normale : $\approx 12/7$ (cm de Hg)} }
  }
  child[concept color=popgold]{ node {Régulation de la PA}
    child[concept color=popgoldL]{ node {Nerfs cardiaques et vaso-moteurs (bulbe)} }
    child[concept color=popgoldL]{ node {Hormones : adrénaline, ADH, aldostérone} }
  }
  child[concept color=poppink]{ node {HTA et hypotension}
    child[concept color=poppinkL]{ node {HTA : PA $\geq 14/9$ ; « tueur silencieux »} }
    child[concept color=poppinkL]{ node {Lutte : régime hyposodé, sport, dépistage} }
  };
\end{tikzpicture}
\legende{Carte mentale de la leçon : les six branches résument la mesure et la régulation de la glycémie (orange, vert), le diabète (violet), la mesure et la régulation de la pression artérielle (bleu, doré) et les troubles de la tension (rose).}
\end{popfigure}

\begin{bilan}
\textbf{1. Définitions et valeurs clés à connaître par c\oe ur}
\begin{itemize}
  \item \cle{Glycémie} : concentration de glucose dans le sang. Valeur normale \textbf{à jeun} : \num{0,7} à \SI{1,1}{g/L} (soit 3,9 à \SI{6,1}{mmol/L}). On parle d'\cle{hyperglycémie} au-delà de \SI{1,26}{g/L} à jeun (à deux reprises) et d'\cle{hypoglycémie} en dessous de \SI{0,6}{g/L}.
  \item Mesure : \cle{glycomètre} (lecteur) + bandelette sur une goutte de sang capillaire (bout du doigt) ; dosage de laboratoire sur sang veineux.
  \item \cle{Pression artérielle} (PA) : pression exercée par le sang sur la paroi des artères. Elle comporte la \cle{maxima} (pression systolique, contraction du c\oe ur) et la \cle{minima} (pression diastolique, relâchement). Valeur normale au repos : environ $\mathbf{12/7}$ (en cm de Hg, soit 120/70 mm de Hg). Mesure au \cle{sphygmomanomètre} (tensiomètre) avec brassard et stéthoscope.
  \item \cle{Hypertension artérielle} (HTA) : PA $\geq 14/9$ de façon permanente ; \cle{hypotension} : PA $\leq 9/6$.
\end{itemize}

\textbf{2. Régulation de la glycémie (schéma fonctionnel exigible)}
\begin{center}
\begin{tabularx}{\linewidth}{|l|X|X|}
\hline
\rowcolor{popblueL} \textbf{Élément} & \textbf{Insuline} & \textbf{Glucagon} \\
\hline
Origine & cellules $\beta$ des îlots de Langerhans (pancréas) & cellules $\alpha$ des îlots de Langerhans \\
\hline
Effet sur la glycémie & \textbf{hypoglycémiante} & \textbf{hyperglycémiante} \\
\hline
Mécanismes & entrée du glucose dans les cellules ; \textbf{glycogénogenèse} hépatique et musculaire & \textbf{glycogénolyse} et néoglucogenèse hépatiques \\
\hline
Moment de sécrétion & après un repas (glycémie élevée) & à jeun, à l'effort (glycémie basse) \\
\hline
\end{tabularx}
\end{center}
Le \cle{foie} est l'organe effecteur central : il stocke le glucose sous forme de glycogène (sous l'action de l'insuline) et le libère (sous l'action du glucagon). Le \cle{pancréas} est à la fois capteur et glande endocrine. La régulation nerveuse passe par le \cle{système nerveux végétatif} (centres bulbaires) et l'\cle{adrénaline} médullo-surrénalienne (hyperglycémiante).

\textbf{3. Le diabète}
\begin{itemize}
  \item \cle{Diabète de type 1} (insulinodépendant, DID) : destruction auto-immune des cellules $\beta$ $\Rightarrow$ absence d'insuline ; sujet jeune ; traitement : injections d'insuline à vie.
  \item \cle{Diabète de type 2} (non insulinodépendant, DNID) : \cle{insulinorésistance} des cellules cibles ; sujet adulte, souvent obèse et sédentaire ; traitement : hygiène de vie, antidiabétiques oraux, parfois insuline.
  \item Signes cardinaux : \cle{polyurie}, \cle{polydipsie}, \cle{polyphagie}, amaigrissement ; présence de glucose dans les urines (\cle{glycosurie}) quand la glycémie dépasse le seuil rénal ($\approx$ \SI{1,8}{g/L}).
  \item Complications : cécité (rétinopathie), insuffisance rénale, neuropathies, infarctus, amputations. Lutte : dépistage, alimentation équilibrée pauvre en sucres rapides, activité physique régulière.
\end{itemize}

\textbf{4. Régulation de la pression artérielle}
\begin{itemize}
  \item La PA dépend du \cle{débit cardiaque} (fréquence $\times$ volume d'éjection systolique) et des \cle{résistances périphériques} (diamètre des artérioles) : PA $=$ débit cardiaque $\times$ résistances.
  \item \cle{Régulation nerveuse rapide} : barorécepteurs de la crosse aortique et des sinus carotidiens $\rightarrow$ centre vaso-moteur du bulbe rachidien $\rightarrow$ nerfs cardiaques (accélérateur sympathique / ralentisseur parasympathique) et nerfs vaso-moteurs (vasoconstriction / vasodilatation).
  \item \cle{Régulation hormonale lente} : adrénaline et noradrénaline (vasoconstriction, tachycardie), ADH (rétention d'eau), système rénine--angiotensine--aldostérone (rétention de Na$^+$ et d'eau $\Rightarrow$ augmentation du volume sanguin).
\end{itemize}

\textbf{5. Méthodes express}
\begin{itemize}
  \item \emph{Interpréter une courbe de glycémie} : repérer la valeur à jeun, le pic post-prandial ($\approx$ \SI{1,4}{g/L} vers 1 h chez le sain), le retour à la normale en $\approx$ 2 h grâce à l'insuline ; chez le diabétique, pic plus haut et retour très lent.
  \item \emph{Schéma fonctionnel} : toujours boucler la boucle (glycémie $\uparrow$ $\rightarrow$ pancréas $\rightarrow$ insuline $\rightarrow$ foie/cellules $\rightarrow$ glycémie $\downarrow$) avec flèches $+$ et $-$.
  \item \emph{Lire une tension} : citer les deux valeurs (maxima/minima), comparer à 12/7, conclure (normotension, HTA ou hypotension) et proposer une conduite à tenir.
\end{itemize}
\end{bilan}

\begin{aretenir}[Checklist avant le Bac]
\begin{itemize}
  \item $\square$ Je connais la glycémie normale à jeun (\num{0,7}--\SI{1,1}{g/L}) et le seuil du diabète (\SI{1,26}{g/L}).
  \item $\square$ Je sais décrire la mesure de la glycémie au glycomètre et celle de la PA au sphygmomanomètre.
  \item $\square$ Je situe les cellules $\alpha$ (glucagon) et $\beta$ (insuline) des îlots de Langerhans.
  \item $\square$ Je sais que l'insuline est la seule hormone hypoglycémiante et je cite deux hormones hyperglycémiantes (glucagon, adrénaline).
  \item $\square$ Je maîtrise les rôles du foie : glycogénogenèse, glycogénolyse, néoglucogenèse.
  \item $\square$ Je sais tracer et légender le schéma fonctionnel complet de la régulation de la glycémie.
  \item $\square$ Je distingue sans hésiter diabète de type 1 et de type 2 (cause, âge, traitement).
  \item $\square$ Je connais les trois signes cardinaux du diabète et le seuil rénal de la glycosurie.
  \item $\square$ Je sais que la PA normale est $\approx 12/7$ et que l'HTA commence à $14/9$.
  \item $\square$ J'explique la régulation nerveuse de la PA (barorécepteurs, bulbe, nerfs cardiaques et vaso-moteurs).
  \item $\square$ Je cite les hormones qui élèvent la PA (adrénaline, ADH, aldostérone) et leur mode d'action.
  \item $\square$ Je sais proposer des mesures de lutte contre l'HTA et le diabète adaptées au contexte camerounais (régime hyposodé, moins de sucres rapides, marche et sport, dépistage en centre de santé).
\end{itemize}
\end{aretenir}

\findecours

\end{document}
